\pdfoutput=1
\documentclass[11pt]{amsart}
\usepackage{amsmath,amssymb,amsthm}
\usepackage[margin=1.2in]{geometry}
\usepackage[colorlinks=true,linkcolor=blue,citecolor=blue,urlcolor=blue]{hyperref}
\hypersetup{pdfauthor={Bernd Johannes Wuebben},
  pdftitle={Small Dolbeault eigenvalues along Hermitian Yang--Mills flow}}

\newtheorem{theorem}{Theorem}[section]
\newtheorem{proposition}[theorem]{Proposition}
\newtheorem{corollary}[theorem]{Corollary}
\newtheorem{lemma}[theorem]{Lemma}
\newtheorem{question}[theorem]{Question}
\theoremstyle{definition}
\newtheorem{definition}[theorem]{Definition}
\theoremstyle{remark}
\newtheorem{remark}[theorem]{Remark}
\newtheorem{example}[theorem]{Example}

\newcommand{\CC}{\mathbb C}
\newcommand{\End}{\operatorname{End}}
\newcommand{\tr}{\operatorname{tr}}
\newcommand{\rk}{\operatorname{rk}}
\newcommand{\id}{\operatorname{Id}}
\newcommand{\op}{\mathrm{op}}
\newcommand{\dvol}{\,dV_\omega}
\newcommand{\db}{\bar\partial}
\newcommand{\Hom}{\operatorname{Hom}}
\newcommand{\diag}{\operatorname{diag}}
\newcommand{\Ker}{\operatorname{ker}}
\hyphenation{semi-sta-ble poly-sta-ble Her-mit-ian}

\title[Small eigenvalues along Hermitian Yang--Mills flow]{Small Dolbeault eigenvalues\\
along Hermitian Yang--Mills flow}
\author{Bernd Johannes Wuebben}
\subjclass[2020]{53C07 (primary); 32L05, 32Q15, 58J50 (secondary)}
\keywords{Hermitian Yang--Mills flow, Dolbeault Laplacian, holomorphic vector
bundles, spectral degeneration, weighted graph Laplacian, slope stability}
\date{September 20, 2026 (initial version); September 30, 2026 (this version)}

\begin{document}

\begin{abstract}
We study the small eigenvalues of the Dolbeault Laplacian on trace-free
endomorphisms of a semistable holomorphic bundle along Hermitian
Yang--Mills flow, normalized as $h^{-1}\partial_\tau h=-2(K_h-\kappa)$ with
$K_h$ the mean curvature and $\kappa$ its average. The
bundles considered are simple bundles on a compact curve whose
Jordan--H\"older factors are pairwise nonisomorphic, and extensions of one
polystable bundle by another on compact K\"ahler manifolds, again with
pairwise nonisomorphic stable factors. From every
smooth initial metric, each small eigenvalue is asymptotic to an eigenvalue
of a weighted graph Laplacian, whose conductances evolve by an explicit
finite-dimensional flow; the relative error tends to zero. The graph records
which stable factors are joined by nonzero extension classes. The
asymptotics are governed by the weight grading of Haiden, Katzarkov,
Kontsevich, and Pandit, the solution of a quadratic program on this graph.
Suppose some Lagrange multiplier vector of this program is strictly positive
on every active constraint. We then determine the decay exponent and the
leading coefficient of every small eigenvalue, and the limits of the
spectral subspaces. The eigenvalues of order $\tau^{-1}$ have coefficients
independent of the initial metric, and $1/2$ is always one of them. At each
faster rate the coefficients depend on the initial metric, and their sum is
unbounded over metrics on one fixed bundle. A key step shows that the
bounded remainder in the asymptotics of the finite-dimensional flow
converges. Examples include an eigenvalue decaying like $\tau^{-4/3}$ and,
without the positivity condition, an eigenvalue of order
$(\tau\log\tau)^{-1}$. We also treat families of harmonic extensions with
fixed metrics, including repeated factors
and recovery of the graded object, and we show that the constants in a slope
inequality relating the lowest eigenvalue to the averaged trace-free mean
curvature are sharp in every rank.
\end{abstract}

\maketitle

\section{Introduction}\label{sec:intro}

The summand projections of a polystable direct sum are holomorphic
endomorphisms and hence zero modes of its Dolbeault Laplacian. Under a
deformation of the holomorphic structure, some of these modes can acquire
positive eigenvalues. We ask how the extension data determine the small
eigenvalues and their limiting spectral subspaces, and what these limits
reveal about the stable summands. Our main concern is the spectrum along
Hermitian Yang--Mills flow, where the time parameter is intrinsic rather than
chosen.
For every simple bundle on a curve whose Jordan--H\"older factors are
pairwise nonisomorphic, and for two-step extensions on compact K\"ahler
manifolds, we reduce the small spectrum to a finite-dimensional flow. Under
a positivity condition on the associated Lagrange multipliers, we determine
the exponent and the leading coefficient of every small eigenvalue, the
limits of the spectral subspaces, and which coefficients retain information
about the initial metric.

On a compact curve, every stable bundle carries a Hermitian--Einstein
metric, whose Chern connection is projectively flat
\cite{NarasimhanSeshadri1965,Donaldson1983}, and the Yang--Mills
functional was analysed from the point of view of Morse theory by
Atiyah and Bott \cite{AtiyahBott1983}. Hermitian Yang--Mills flow
\cite{Donaldson1985,Simpson1988} carries a semistable bundle on a curve
to a Hermitian--Einstein connection on its graded bundle, modulo unitary
gauge \cite{Daskalopoulos1992,Rade1992}. In higher dimensions the
limiting reflexive sheaf is the double dual of the graded object of the
Harder--Narasimhan--Seshadri filtration
\cite{DaskalopoulosWentworth2004,Jacob2015,Sibley2015}. When that graded
object is locally free, the convergence is smooth on the whole manifold
\cite{SibleyWentworth2015}.

The exponents are governed by the weight grading of Haiden, Katzarkov,
Kontsevich, and Pandit \cite{HKKP2023}, a solution of a quadratic program
on the directed graph of nonzero extension entries. When the Lagrange multipliers of that program are strictly positive,
which is the hypothesis of \cite[Theorem~3.6]{HKKP2023}, we show the
following. The eigenvalues decaying like $\tau^{-1}$ have coefficients that
do not depend on the initial metric, and $1/2$ is always one of them. At
each faster rate the coefficients depend on the initial metric, and their
sum is unbounded over metrics on one fixed bundle. The proof has three parts: a
comparison of every small eigenvalue with a finite-dimensional diagonal
flow, with relative error tending to zero on the entire future interval; a
proof that the bounded remainder in the asymptotics of that flow converges;
and a contraction rule for weighted graph Laplacians whose weights decay at
different rates. The finite-dimensional models, the exponents, and the
logarithmic example come from \cite{HKKP2023,HKKP2018}. Our contribution is
the transfer of precise coefficients to the Dolbeault spectrum of the
bundle, the convergence of the remainder, and the resulting dichotomy.
Section~\ref{sec:earlier-work} gives the detailed attribution.

The spectral input is a relative comparison for harmonic extension
families with fixed background metrics. Scaled harmonic extensions of
this kind were used by Sektnan and Tipler in an adiabatic Hermitian
Yang--Mills problem \cite[Section~5]{SektnanTipler2024}. A finite-dimensional commutator
form describes the small spectrum; for pairwise nonisomorphic stable
factors it is a weighted graph Laplacian. The error is relative to each
eigenvalue, so the estimate remains effective when different modes
vanish at different rates.

These asymptotics also give quantitative information about slope
stability. The classical subsheaf degree formula
\cite{Kobayashi1987,Lubke1983} yields an inequality
relating slopes, the lowest endomorphism eigenvalue, and the trace-free
mean curvature. Corrected two-factor families show that its constants are
sharp in every rank. Thus the explicit spectral families
both describe the approach to a polystable splitting and test the
optimality of a general stability estimate.

We use the following conventions throughout. Let $(X,\omega)$ be a
compact connected K\"ahler manifold of complex dimension $n\ge1$, and set
\[
 dV_\omega=\frac{\omega^n}{n!},\qquad V=\int_X dV_\omega.
\]
For a torsion-free coherent sheaf $\mathcal F$ of positive rank, define
\[
 \deg_\omega\mathcal F
 =\int_X c_1(\mathcal F)\wedge\frac{\omega^{n-1}}{(n-1)!},
 \qquad
 \mu_\omega(\mathcal F)=\frac{\deg_\omega\mathcal F}{\rk\mathcal F}.
\]
Here $c_1(\mathcal F)=c_1(\det\mathcal F)$. Following Mumford and
Takemoto \cite{Mumford1963,Takemoto1972}, the bundle $E$ is
slope-stable if $\mu_\omega(\mathcal F)<\mu_\omega(E)$ for every
coherent subsheaf with $0<\rk\mathcal F<\rk E$; semistability uses
the non-strict inequality; see \cite[Chapter~V]{Kobayashi1987}. It suffices to consider saturated subsheaves,
for which $E/\mathcal F$ is torsion-free.

For a holomorphic bundle $E$ of rank $r\ge2$ with a smooth Hermitian
metric $h$, let $\db_{\End E}$ denote the induced Dolbeault operator on
$\End E$, and define
\begin{equation}\label{eq:lambda}
 \lambda_1(h)=
 \inf_{0\ne u\in C^\infty(X,\End_0 E)}
 \frac{\lVert\db_{\End E}u\rVert_{L^2}^2}{\lVert u\rVert_{L^2}^2}.
\end{equation}
All $L^2$ norms use $h$, $\omega$, and the Hilbert--Schmidt norm on
endomorphisms. Thus $\lambda_1(h)$ is the lowest eigenvalue, including
zero, of $\db_{\End E}^{*}\db_{\End E}$ on trace-free endomorphisms;
it is not the first positive eigenvalue when the kernel is nonzero.
The bundle is simple if $H^0(X,\End E)=\CC\id_E$. By elliptic theory,
$\lambda_1(h)>0$ if and only if $E$ is simple.

\subsection{The small spectrum along the flow}\label{sec:intro-flow}

Let $X$ be a compact curve, and let $E$ be a simple semistable bundle whose
Jordan--H\"older factors $F_1,\ldots,F_m$, $m\ge2$, are pairwise
nonisomorphic, of ranks $r_i$. The \emph{socle} of a semistable bundle is
the sum of its stable subbundles of the same slope
\cite[Section~1.5]{HuybrechtsLehn2010}. The socle filtration
$S_1\subset S_2\subset\cdots\subset E$ takes $S_1$ to be the socle of $E$
and $S_{k+1}/S_k$ to be the socle of $E/S_k$. Each $F_i$ occurs in exactly
one quotient $S_k/S_{k-1}$, its \emph{socle layer}. A \emph{two-step
extension} is one with two socle layers,
\[
 0\longrightarrow S\longrightarrow E\longrightarrow U\longrightarrow0,
\]
with $S$ and $U$ polystable of slope $\mu(E)$.

By Lemma~\ref{lem:socle-normal-form}, $E$ is smoothly identified with
$(\bigoplus_iF_i,\db_0+N)$, where $N=\sum\beta_{ij}$, each
$\beta_{ij}\in\Omega^{0,1}(\Hom(F_j,F_i))$ is harmonic, and $\beta_{ij}\ne0$
only when $F_i$ lies in a lower socle layer than $F_j$. Let $\Gamma$ be the
directed graph with an edge $i\to j$ for each nonzero $\beta_{ij}$, and put
$M_i=Vr_i$. We call $\sum_ix_i\id_{F_i}$, $x_i\in\CC$, a
\emph{block-scalar} endomorphism. Write $K_h=\sqrt{-1}\Lambda_\omega F_h$
for the mean curvature, and run the normalized metric flow
\[
 h^{-1}\partial_\tau h=-2(K_h-\kappa\id),\qquad \kappa=2\pi\mu(E)/V,
\]
from any smooth metric. The flow converges to the Hermitian--Einstein
connection of the graded bundle \cite{Daskalopoulos1992,Rade1992}. Exactly
$m-1$ eigenvalues $\lambda_1\le\cdots\le\lambda_{m-1}$ of the Dolbeault
Laplacian on trace-free endomorphisms therefore tend to zero.

The finite-dimensional comparison object is the \emph{diagonal flow}. It is
the flow on block-scalar metrics $g^*h_0$, $g=\diag(e^{y_i}\id_{F_i})$, in
which the mean curvature is replaced by its block-scalar average:
\begin{equation}\label{eq:intro-diag-flow}
 \dot y_i=-\frac1{M_i}\Bigl(\sum_{i\to j}k_{ij}-\sum_{l\to i}k_{li}\Bigr),
 \qquad k_{ij}=\lVert\beta_{ij}\rVert_{L^2}^2e^{2(y_i-y_j)}.
\end{equation}
The edge weight, or \emph{conductance}, $k_{ij}$ is the squared norm of the
rescaled extension entry. This is the weighted-graph flow of
\cite[(3.10)]{HKKP2023}, and we call $y$ its logarithmic coordinates.

The \emph{weight grading} $w$ of \cite{HKKP2023} is the minimizer of
$\frac12\sum_iM_iw_i^2$ subject to $w_j-w_i\ge\frac12$ for every edge
$i\to j$ (Definition~\ref{def:weight-grading}). An edge is \emph{tight}
when equality holds, that is, when its constraint is active; the tight edges form the \emph{tight subgraph}
$\mathcal T$, and its connected components, isolated vertices included, are
the \emph{tight components}. The Lagrange multipliers of this program vanish
off $\mathcal T$. We call $(\Gamma,(r_i))$ \emph{nondegenerate} when some
multiplier vector is strictly positive on every tight edge
(Definition~\ref{def:nondegenerate}). In optimization this is strict
complementarity; it is the hypothesis of \cite[Theorem~3.6]{HKKP2023}.

Each edge receives the exponent $a_{ij}=2(w_j-w_i)\ge1$. For each value $a$
of these exponents, contract the connected components of the edges with
smaller exponent, adding the masses, and join the resulting vertices by the
edges of exponent $a$. Call this the \emph{$a$-quotient graph}, and let
$d_a$ be the number of its positive Laplacian eigenvalues. The small
eigenvalues are assigned exponents accordingly: $d_a$ of them receive the
exponent $a$, larger exponents taking smaller indices. Write $a_j$ for the
exponent of $\lambda_j$.

\begin{theorem}\label{thm:intro-main}
Assume $(\Gamma,(r_i))$ is nondegenerate. For every smooth initial metric:
\begin{enumerate}
\item $\tau^{a_j}\lambda_j(h(\tau))\to c_j(h)>0$ for every $j<m$.
\item The coefficients with $a_j=1$ are the positive eigenvalues of the
 Laplacian of $\mathcal T$ with masses $M_i$ and conductances $u^*_{ij}$.
 Here $u^*$ is the unique multiplier vector of the form
 $u^*_{ij}=\lVert\beta_{ij}\rVert^2e^{2(z_i-z_j)}$ on the edges of
 $\mathcal T$, with $z\in\mathbb R^m$, and zero elsewhere. These coefficients do
 not depend on the initial metric. The value $\frac12$ is always among
 them; for each tight component $\mathcal C$ with at least two vertices it
 has the eigenvector $\sum_{i\in\mathcal C}w_i\id_{F_i}$.
\item For every exponent $a>1$ with $d_a\ge1$, the coefficients with
 $a_j=a$ are the positive eigenvalues of the $a$-quotient graph with
 conductances $k^\infty_{ij}(h)=\lim\tau^{a_{ij}}k_{ij}(g_\tau)$. Here
 $g_\tau$ is the diagonal part of $h(\tau)$ in the normal form
 (Lemma~\ref{lem:socle-matching}), and these limits depend on the initial
 metric. The sum of the coefficients is unbounded as the initial metric
 varies on $E$.
\item In unitary gauges in which the flow converges, the spectral
 projection onto the eigenvalues with a given exponent and a given
 coefficient converges smoothly. The limit is the corresponding eigenspace
 of the Laplacian of the $a$-quotient graph ($\mathcal T$ when $a=1$),
 embedded as block-scalar endomorphisms constant on the contracted
 vertices.
\end{enumerate}
The same holds for two-step extensions on compact K\"ahler manifolds of any
dimension, when the factors are vector bundles.
\end{theorem}

Theorem~\ref{thm:intro-main} combines Theorems~\ref{thm:complete-spectrum}
and~\ref{thm:higher-dim}. Its basis is a comparison that needs no
nondegeneracy (Theorem~\ref{thm:transfer}): for every sufficiently large
$T$ there are a solution $g^{[T]}$ of \eqref{eq:intro-diag-flow} on
$[T,\infty)$ and $\delta_T\to0$ with
\[
 e^{-\delta_T}\nu_j(g^{[T]}(\tau))\le\lambda_j(h(\tau))\le
 e^{\delta_T}\nu_j(g^{[T]}(\tau)),\qquad\tau\ge T,
\]
where the $\nu_j$ are the positive eigenvalues of the weighted graph
Laplacian with masses $M_i$ and conductances $k_{ij}$. Without
nondegeneracy, iterated logarithms can appear in the diagonal flow.

Three examples show the range of behaviour.
\begin{itemize}
\item For a nonsplit extension of two stable bundles,
 $\tau\lambda_1(h(\tau))\to\frac12$. The limit is independent of the ranks,
 the extension class, the initial metric, and the dimension
 (Theorems~\ref{thm:flow-two-factor} and~\ref{thm:higher-dim}).
\item Take four stable factors of ranks $(2,1,1,2)$ with nonzero entries
 exactly on the edges $1\to2$, $3\to2$, $3\to4$ (a zigzag). Two eigenvalues
 are asymptotic to $1/(2\tau)$, and $\lambda_1\sim c(h(0))\tau^{-4/3}$. With
 $V=1$, unit-norm entries, and the initial metrics $h_\rho=g_\rho^*h_0$ with
 $g_\rho=\diag(\rho^{1/2},\rho^{-1/2},\rho^{1/2},\rho^{-1/2})$, one has
 $c(h_\rho)=\frac2{27}\rho^{-2/3}(1+O(\rho))$
 (Proposition~\ref{prop:unequal-rank-family}).
\item For four line bundles with the same zigzag, the multiplier on the edge
 $3\to2$ is forced to vanish, so $(\Gamma,(r_i))$ is degenerate. Here
 $\lambda_1\sim(2\tau\log\tau)^{-1}$, an eigenvalue separated from the
 others by a logarithmic factor, and $\lambda_2\sim\lambda_3\sim(2\tau)^{-1}$
 (Theorem~\ref{thm:zigzag-flow}).
\end{itemize}
The exponent $4/3$ and the logarithmic model are features of the weighted
graph flows of \cite[Section~3.2 and Example~5.13]{HKKP2023}. The theorems
here determine the corresponding Dolbeault coefficients along the actual
flow.

The finite-dimensional input is Theorem~\ref{thm:diag-convergence}. For
nondegenerate $(\Gamma,(r_i))$, every solution of \eqref{eq:intro-diag-flow}
satisfies
\[
 y(\tau)=w\log\tau+\zeta^*+\sum_{\mathcal C}Z_{\mathcal C}\mathbf 1_{\mathcal C}
 +O(\tau^{-\delta}).
\]
Here $\zeta^*$ is the minimizer of an explicit strictly convex function on
the tight subgraph, determined by $\Gamma$, the ranks, $V$ and the norms
$\lVert\beta_{ij}\rVert$; $\mathbf 1_{\mathcal C}$ is the indicator vector
of the tight component $\mathcal C$. The constants $Z_{\mathcal C}$ depend on
the solution. This refines the bounded remainder of
\cite[Theorem~3.6]{HKKP2023}. It also yields convergence of the bounded term
in the metric asymptotics along the actual flow
(Corollary~\ref{cor:metric-asymptotics}). The constants $Z_{\mathcal C}$
are exactly what makes the faster coefficients depend on the initial metric.

\subsection{Harmonic extensions and their small spectrum}

First suppose that $X$ is a compact curve. Fix pairwise nonisomorphic
stable bundles $F_1,\ldots,F_m$, $m\ge2$, of equal slope and ranks $r_i$,
and equip $E_0=\bigoplus_iF_i$ with a product Hermitian--Einstein metric.
Choose harmonic forms
$\beta_{ij}\in\Omega^{0,1}(\Hom(F_j,F_i))$, $i<j$, at least one of
which is nonzero. For strictly decreasing integers
$\ell_1>\cdots>\ell_m=0$, set
\[
 p_{ij}=\ell_i-\ell_j,\qquad
 \db_t=\db_0+\sum_{i<j}t^{p_{ij}}\beta_{ij},\qquad t>0.
\]
Each entry occupies its indicated upper-triangular block. Integrability
is automatic on a curve; denote the resulting holomorphic bundle by
$E_t$. The diagonal transformation
$g_t=\diag(t^{\ell_i}\id_{F_i})$ identifies all positive parameters with
the same holomorphic bundle. We retain the fixed product metric when
computing the spectrum. Equivalently, $g_t$ transports it to a family
of metrics on that fixed holomorphic bundle.

The split trace-free holomorphic endomorphisms form the space
\[
 \mathcal K=\left\{\diag(x_i\id_{F_i}):\sum_i r_ix_i=0\right\}.
\]
Associate an undirected graph to the nonzero entries $\beta_{ij}$, with
vertex masses $Vr_i$ and edge conductances
$t^{2p_{ij}}\lVert\beta_{ij}\rVert_{L^2}^2$. Its weighted Laplacian
$L_t$ on $\mathcal K$ is defined by
\[
 \langle x,L_tx\rangle
 =\sum_{i<j}t^{2p_{ij}}\lVert\beta_{ij}\rVert_{L^2}^2|x_i-x_j|^2,
 \qquad \lVert x\rVert^2=V\sum_i r_i|x_i|^2.
\]
Let $\nu_1(t)\le\cdots\le\nu_{m-1}(t)$ be its eigenvalues and
$\lambda_1(t)\le\lambda_2(t)\le\cdots$ those of the Dolbeault Laplacian
on trace-free endomorphisms of $E_t$, counting complex multiplicity and including
zero. Put $p_* =\min_{\beta_{ij}\ne0}p_{ij}$. The following statement
is proved in Section~\ref{sec:several}.

\begin{theorem}[Spectral comparison]\label{thm:intro-spectral}
Exactly $m-1$ eigenvalues tend to zero as $t\to0$, and the remaining
spectrum stays uniformly bounded away from zero. For every positive
$\nu_j(t)$,
\[
 \frac{\lambda_j(t)}{\nu_j(t)}=1+O(t^{2p_*}),
\]
with an error uniform in $j$ for the fixed family. If the graph has $c$
connected components, exactly $c-1$ of the small eigenvalues are
identically zero. Let $\Pi_t$ be the orthogonal projection onto the full
small spectral subspace and $P$ the orthogonal projection onto
$\mathcal K$. For every $k\ge0$,
\[
 \lVert\Pi_t-P\rVert_{L^2\to C^k}=O(t^{2p_*}).
\]
\end{theorem}

Theorem~\ref{thm:orders} determines every leading decay order by a
finite construction. At an edge exponent $p$, contract the
components joined by edges with $p_{ij}<p$ and add their vertex
masses. Join the resulting vertices using entries with $p_{ij}=p$,
assigning conductances $\lVert\beta_{ij}\rVert_{L^2}^2$ and adding
contributions between the same pair of vertices. The positive
eigenvalues $\rho$ of this quotient graph give
precisely the coefficients of the modes
$\lambda(t)=t^{2p}(\rho+o(1))$. Their spectral subspaces converge
smoothly to the corresponding quotient eigenspaces, embedded as
endomorphisms constant on each contracted component. For a repeated
coefficient, the conclusion concerns the whole eigenspace. For the
matrix family $L_t$ alone, the contraction rule follows from the
reduction process of analytic perturbation theory
\cite[Chapter~II, \S2.3]{Kato1995}; Theorem~\ref{thm:orders} transfers
it, together with the limiting subspaces, to the Dolbeault spectrum.

The limiting full subspace has a geometric interpretation. After
adjoining the scalar identity to $\mathcal K$, one obtains the algebra
$\bigoplus_i\CC\id_{F_i}$. Its minimal nonzero idempotents are the
projections onto the stable summands. The limiting spectral subspace,
viewed as a space of endomorphisms with pointwise multiplication,
therefore identifies the unordered summands of this split
limit.

For a concrete example, take three pairwise nonisomorphic degree-zero
line bundles on a curve of genus at least two and volume one. Choose
harmonic forms of unit $L^2$ norm on the consecutive blocks $12,23$,
scale them by $t,t^2$, and set the $13$ entry to zero. The two small
eigenvalues satisfy
\[
 \lambda_1(t)=\tfrac32t^4(1+O(t^2)),\qquad
 \lambda_2(t)=2t^2(1+O(t^2)).
\]
The limiting eigensection for the smaller eigenvalue distinguishes
$F_1\oplus F_2$ from $F_3$; that for the larger eigenvalue distinguishes
$F_1$ from $F_2$. An absolute error of order $t^4$
would obscure the first leading coefficient, whereas the relative
comparison determines both. Section~\ref{sec:examples} derives these
formulas and gives a five-factor example with a repeated coefficient.

\subsection{Realization for semistable bundles on curves}

The harmonic families can be constructed from every simple strictly
semistable bundle $E$ on a compact curve. Choose a Jordan--H\"older
filtration with stable quotients $F_1,\ldots,F_m$, allowing repetitions.
Its graded bundle
$G=\operatorname{gr}(E)=\bigoplus_iF_i$ is determined up to isomorphism
by $E$ \cite{Seshadri1967}; see also
\cite[Section~1.5]{HuybrechtsLehn2010}. After choosing
Hermitian--Einstein metrics on the quotients, compatible with fixed
identifications between isomorphic copies,
Lemma~\ref{lem:normal-form} gives a filtration-preserving smooth
identification in which all upper-triangular extension entries are
harmonic. The construction corrects successive upper diagonals and
retains the changes they induce in higher entries. It requires no
smallness assumption on the original extension.

Write $G=\bigoplus_\alpha H_\alpha\otimes\CC^{m_\alpha}$ with
pairwise nonisomorphic stable $H_\alpha$. Since stable bundles are
simple and a nonzero homomorphism between stable bundles of equal slope
is an isomorphism
\cite[Chapter~V, Proposition~7.11 and Corollaries~7.12, 7.14]{Kobayashi1987},
its holomorphic endomorphism algebra is
$\bigoplus_\alpha M_{m_\alpha}(\CC)$.
Theorem~\ref{thm:matrix-relative} extends the relative comparison
to the quadratic form below, where $N_t$ is the harmonic extension
form after diagonal scaling:
\[
 \langle x,L_tx\rangle=\lVert[N_t,x]\rVert_{L^2}^2,
 \qquad x\in H^0(X,\End_0G).
\]
Theorem~\ref{thm:curve-recovery} constructs approximate
Hermitian--Einstein metrics $h_t$
(cf.\ \cite[Chapter~VI, Theorem~10.13]{Kobayashi1987} and
\cite{Jacob2014}) and explicit gauges \eqref{eq:recovery-gauges} in which this comparison
applies. Exactly $d=\sum_\alpha m_\alpha^2-1$ positive
eigenvalues tend to zero, the remaining spectrum has a uniform gap,
and the full small spectral projection converges smoothly to the
projection onto $H^0(X,\End_0G)$. Adjoining the scalar identity
recovers the endomorphism algebra of the graded bundle. Its central
idempotents identify the isotypical components, its matrix sizes give
the multiplicities, and the images of its minimal idempotents give
the stable factors up to isomorphism. Individual repeated copies
have no canonical projections. For example, Section~\ref{sec:repeated-example}
constructs a simple bundle with graded object $A\oplus B\oplus A\oplus C$
and five small modes, with leading terms $t^4,t^2,t^2,2t^2,2t^2$.
These recovery statements concern the constructed metric families.

\subsection{Slope inequalities and their sharpness}

We now return to an arbitrary smooth Hermitian metric $h$ on a bundle
$E$ over a compact connected K\"ahler manifold. The implication from
irreducible Hermitian--Einstein metrics to slope stability rests on
the subbundle curvature formula and its extension to coherent
subsheaves \cite{Kobayashi1987,Lubke1983}. The converse is the
Donaldson--Uhlenbeck--Yau theorem
\cite{Donaldson1985,UhlenbeckYau1986}. For a general metric, the same
curvature formula gives a quantitative estimate with an error term.

Write $F_h$ for its Chern curvature, with the convention
$c_1(E,h)=\frac{\sqrt{-1}}{2\pi}\tr F_h$, and put
\begin{equation}\label{eq:curvature}
 K_h=\sqrt{-1}\Lambda_\omega F_h,
 \qquad K_h^0=K_h-\frac{\tr K_h}{r}\id_E.
\end{equation}
Both endomorphisms are self-adjoint with respect to $h$. The quantity used
to measure the curvature error is
\begin{equation}\label{eq:epsilon}
 \varepsilon_1(h)=\frac1V\int_X\lVert K_h^0\rVert_{\op}\dvol.
\end{equation}
\begin{theorem}\label{thm:main}
For every saturated coherent subsheaf $\mathcal S\subset E$ of rank
$0<s<r$, one has
\begin{equation}\label{eq:main}
 \mu_\omega(E)-\mu_\omega(\mathcal S)
 \ge \frac{V(r-s)}{2\pi r}
 \left(\lambda_1(h)-\frac{r}{\max\{s,r-s\}}\varepsilon_1(h)\right).
\end{equation}
In particular, if
\begin{equation}\label{eq:criterion}
 \lambda_1(h)>c_r\varepsilon_1(h),
 \qquad c_r=\frac{r}{\lceil r/2\rceil},
\end{equation}
then $E$ is slope-stable.
\end{theorem}

The proof combines the classical subsheaf degree formula with the
Rayleigh principle; we have not found the inequality stated in this form. Its test section is the orthogonal projection onto
the subsheaf with its scalar part removed; its pointwise squared norm
is $s(r-s)/r$. A trace-free matrix estimate gives the displayed rank
dependence. Since $c_r\le2$ and
$\varepsilon_1(h)\le\lVert K_h^0\rVert_{L^\infty,\op}$, the simpler
condition $2\lVert K_h^0\rVert_{L^\infty,\op}<\lambda_1(h)$ also
implies stability. No simplicity assumption is needed: the strict
inequality already implies $\lambda_1(h)>0$.

A strictly semistable bundle has a proper saturated subsheaf of equal
slope. Theorem~\ref{thm:main} therefore forces $\lambda_1(h_j)\to0$
whenever $\varepsilon_1(h_j)\to0$, in particular along approximate
Hermitian--Einstein metrics. This conclusion concerns arbitrary such
sequences; it does not require the harmonic form of the families above.

The relation to the explicit spectral asymptotics is clearest for two
factors. Let $F,G$ be nonisomorphic stable bundles of equal slope on a
compact K\"ahler manifold, of ranks $a,b$, with fixed
Hermitian--Einstein metrics and a nonzero harmonic extension form
$\beta\in\Omega^{0,1}(\Hom(G,F))$. Scaling the extension by $t$ gives
exactly one small eigenvalue. Theorem~\ref{thm:two-factor} gives
\[
 \lambda_1(t)=\frac{(a+b)\lVert\beta\rVert_{L^2}^2}{Vab}\,t^2
 -dt^4+O(t^6),\qquad d\ge0,
\]
with an explicit formula for $d$ in terms of the inverse split
Laplacian on the complement of its kernel. In the gauge $g_t$ of
\eqref{eq:two-gauge}, the
normalized eigensection converges smoothly to the normalized trace-free
part of a summand projection. This two-factor result holds in every
complex dimension.

A second-order correction of the metric, by the Green operator applied to
the trace-free part of the curvature, makes both estimates in the proof of
Theorem~\ref{thm:main} asymptotically sharp. On a curve of genus at least
two, for every $r\ge2$ and $1\le s<r$, this gives nonsplit extensions of a
stable bundle of rank $r-s$ by one of rank $s$, and metrics with
\[
 \frac{\lambda_1}{\varepsilon_1}\longrightarrow\frac r{\max\{s,r-s\}}
\]
(Corollary~\ref{cor:sharpness}). These bundles are strictly semistable, so
no coefficient in \eqref{eq:main} can be decreased, and $c_r$ is optimal in
every rank. The same families show that $c_r$ is optimal for the criterion
with $\lVert K_h^0\rVert_{L^\infty,\op}$ in place of $\varepsilon_1$.

\subsection{Organization}

Sections~\ref{sec:two-factor}--\ref{sec:examples} treat harmonic families
with fixed metrics. They contain the two-factor expansion, the relative
comparison with the graph Laplacian, and the contraction rule for decay
orders. Section~\ref{sec:normal-form} gives the harmonic normal form on
curves. Sections~\ref{sec:flow}--\ref{sec:higher-dim} contain the flow
results.
\begin{itemize}
\item Section~\ref{sec:flow} treats two factors directly.
\item Section~\ref{sec:diagonal-flow} proves the sharp asymptotics of the
 diagonal flow.
\item Section~\ref{sec:transfer} compares the bundle flow with diagonal
 flows begun at arbitrarily late times.
\item Section~\ref{sec:complete-spectrum} proves
 Theorem~\ref{thm:intro-main}, Section~\ref{sec:flow-examples} gives the
 examples, and Section~\ref{sec:higher-dim} treats higher dimensions.
\end{itemize}
The comparison at late times is essential. A comparison begun at a single
fixed time gives only two-sided bounds of order $\tau^{-a}$, not the
coefficients. Sections~\ref{sec:repeated} and~\ref{sec:curve-recovery}
return to fixed metrics, with repeated factors and a recovery theorem. The
slope inequality is proved at the end.

\subsection{Relation to earlier work and scope}\label{sec:earlier-work}

The scaled-extension construction is classical in this setting.
Sektnan--Tipler \cite[Section~5]{SektnanTipler2024} use it in an
adiabatic Hermitian Yang--Mills problem; their Proposition~41 computes
the action on summand identities, and Lemma~52 uses a weighted quadratic
form in their differences. Their linearized operator is a connection
Laplacian with a varying base metric. Here we work with the nonnegative
Dolbeault Laplacian and fixed metrics, derive a relative comparison for
each small eigenvalue, and prove convergence of the full small spectral
subspace. For pairwise nonisomorphic factors, graph contraction also
determines every leading coefficient and its limiting subspace.
The two-factor expansion uses standard analytic perturbation theory
\cite[Chapter~VII]{Kato1995}.

For the flow theorems, existence and convergence to the graded bundle are
classical inputs \cite{Daskalopoulos1992,Rade1992}; in higher dimensions we
use Sibley and Wentworth's identification of the singular set
\cite[Theorem~1.1]{SibleyWentworth2015}. Haiden, Katzarkov, Kontsevich, and
Pandit \cite{HKKP2018} describe metric growth along the flow on curves by
weight filtrations and iterated logarithms, with a bounded remainder. Their
weighted-graph analysis \cite[Section~3]{HKKP2023} supplies the weight
grading, the multiplier condition, the convexity argument, the regime giving
the exponent $4/3$, and the logarithmic model of Example~5.13. Our proofs
use their Green operator correction and order comparison
\cite[Sections~3.5 and~4.2]{HKKP2018}.

Lam \cite[Theorems~1.1 and~8.2, Proposition~8.6]{Lam2026} extends the
metric asymptotics to compact K\"ahler manifolds with locally free graded
object, including repeated factors, again with a bounded remainder. A fixed
multiplicative bound between metrics leaves a fixed uncertainty in each
leading coefficient. Theorem~\ref{thm:transfer} removes that uncertainty by
comparing with diagonal metrics begun at arbitrarily late times. The
convergence of the remainder in the diagonal flow,
Theorem~\ref{thm:diag-convergence}, then determines the coefficients. That
convergence also sharpens the metric asymptotics in the nondegenerate case.

After a linear change of variables, the diagonal flow is gradient flow for
an exponential loss with linearly separable data, and the weight grading is
the hard-margin direction (Remark~\ref{rem:implicit-bias}). For gradient
descent with unit weights, Soudry, Hoffer, Nacson, Gunasekar, and Srebro
\cite{Soudry2018} prove convergence of the remainder when the support vectors
span the data. Here that means there is a single tight component. Nacson et
al.\ \cite[Theorem~10]{Nacson2019} bound the component of the remainder
outside the span of the support vectors without that assumption.
In our setting there can be several tight components; the argument for
convergence in that case is short. The dependence on the initial metric
enters exactly through those components.

The harmonic gauge condition itself also appears in
\cite[Section~5.1]{SektnanTipler2024}; our triangular construction on a
curve supplies the normal form needed to apply the spectral comparison to
the bundle $E$.

We state the scope of the fixed-metric results once. They concern the
constructed families: harmonic extension data, a chosen filtration,
quotient metrics, and scaling exponents, with rates referring to the chosen
parameter $t$. The flow results use the intrinsic time $\tau$. The
identification of stable factors uses the limiting space of endomorphisms
with its multiplication, not the eigenvalue list alone, and it does not
select a canonical Jordan--H\"older filtration. For arbitrary approximate
Hermitian--Einstein sequences, the slope inequality gives eigenvalue
collapse, but convergence of spectral subspaces is not asserted.

\section{An extension of two stable bundles}\label{sec:two-factor}

Let $(X,\omega)$ again be a compact connected K\"ahler manifold. Let
$F,G$ be nonisomorphic stable bundles of equal slope, with ranks $a,b$,
and put $r=a+b$. Fix Hermitian--Einstein metrics $h_F,h_G$, whose common
Einstein constant is $\kappa=2\pi\mu_\omega(F)/V$, and a nonzero harmonic
form
\[
 \beta\in\Omega^{0,1}(\Hom(G,F)),\qquad
 \db\beta=0,\qquad \db^*\beta=0.
\]
All inner products in this section use $h_0=h_F\oplus h_G$ and the fixed
K\"ahler form. We use the convention
$\langle u,v\rangle=\int_X\tr(u^*v)\dvol$.

On the fixed smooth bundle $F\oplus G$ define
\begin{equation}\label{eq:two-family}
 N=\begin{pmatrix}0&\beta\\0&0\end{pmatrix},\qquad
 \db_t=\db_0+tN,\qquad t\in\mathbb R.
\end{equation}
This is integrable in every complex dimension because $\db_0N=0$ and
$N\wedge N=0$. Denote the resulting holomorphic bundle by $E_t$.
On trace-free endomorphisms put
\[
 D_t=D_0+tB,\qquad D_0=\db_{0,\End},\qquad
 Bu=[N,u],\qquad \Delta_t=D_t^*D_t.
\]
These nonnegative self-adjoint elliptic operators have the fixed domain
$H^2$ in $L^2$. Eigenvalues count complex multiplicity and include zero.
Let $P_F$ be the orthogonal projection onto $F$ and set
\begin{equation}\label{eq:two-q}
 q=P_F-\frac ar\id
 =\diag\left(\frac br\id_F,-\frac ar\id_G\right),\qquad
 V_q=\lVert q\rVert_{L^2}^2=\frac{Vab}{r}.
\end{equation}

Write $\beta=\sum_\nu\beta_\nu\bar\theta^\nu$ in a local unitary
coframe, and define the global self-adjoint trace-free endomorphism
\begin{equation}\label{eq:two-C}
 C=\diag\left(\sum_\nu\beta_\nu\beta_\nu^*,
                  -\sum_\nu\beta_\nu^*\beta_\nu\right).
\end{equation}
Stability gives $\Ker\Delta_0=\CC q$, as verified below. Let $P$ be the
$L^2$ projection onto this kernel, $Q=\id-P$, and
$R=(\Delta_0|_{QH^2})^{-1}:QL^2\to QH^2$. Define
\begin{equation}\label{eq:two-coefficients}
 c=\frac{r\lVert\beta\rVert_{L^2}^2}{Vab},\qquad
 z=C-cq=QC,\qquad d=\frac{\langle z,Rz\rangle}{V_q}.
\end{equation}

\begin{theorem}\label{thm:two-factor}
For the family \eqref{eq:two-family}, $E_t$ is simple and strictly
semistable for every $t\ne0$. Exactly one eigenvalue of $\Delta_t$ tends
to zero, and
\begin{equation}\label{eq:two-expansion}
 \lambda_1(t)=ct^2-dt^4+O(t^6),\qquad c>0,\quad d\ge0.
\end{equation}
There are $\delta,t_0>0$ such that $\lambda_2(t)\ge\delta$ for
$|t|<t_0$. The normalized first eigensection, with phase chosen so that
$\langle q,u_t\rangle>0$, satisfies for every $k\ge0$
\begin{equation}\label{eq:two-section}
 u_t=\frac{q}{\sqrt{V_q}}-\frac{t^2}{\sqrt{V_q}}Rz+O_{C^k}(t^3).
\end{equation}
The mean curvature and its averaged trace-free norm are exactly
\begin{equation}\label{eq:two-curvature}
 K_t=\kappa\id+t^2C,\qquad
 \varepsilon_1(t)=\frac{t^2}{V}\int_X\lVert C\rVert_{\op}\dvol.
\end{equation}
Moreover, $\lambda_1(t)\le ct^2$ for every real $t$.
\end{theorem}

\begin{proof}
Nonisomorphism and stability at equal slope imply
\[
 H^0(\Hom(F,G))=H^0(\Hom(G,F))=0,\qquad
 H^0(\End F)=\CC\id_F,\quad H^0(\End G)=\CC\id_G.
\]
Thus the trace-free holomorphic endomorphisms of the split bundle are
precisely $\CC q$. This proves the assertion about $\Ker\Delta_0$ and
the existence of $R$.

For $t\ne0$, an endomorphism of $E_t$ preserves $F$, since its composite
$F\to E_t\to G$ vanishes. It induces scalars $\alpha,\eta$ on $F,G$,
and its upper-right block $v$ satisfies
\[
 \db v+t(\eta-\alpha)\beta=0.
\]
Pairing with $\beta$ gives $\eta=\alpha$. The block $v$ is then
holomorphic and hence zero, proving simplicity. For any subsheaf of
$E_t$, its intersection with $F$ and image in $G$ have slopes at most
the common slope whenever their ranks are nonzero. Degree and rank
additivity prove semistability, and the equal-slope subbundle $F$
rules out stability.

Expand on the common domain:
\begin{equation}\label{eq:two-operator}
 \Delta_t=\Delta_0+tA_1+t^2A_2,\qquad
 A_1=D_0^*B+B^*D_0,\quad A_2=B^*B.
\end{equation}
Since $D_0q=0$, $Bq=-N$, and $D_0^*N=0$ by harmonicity, one has
$A_1q=0$. The adjoint formula
$B^*v=\sum_\nu[N_\nu^*,v_\nu]$ gives
\begin{equation}\label{eq:two-action}
 A_2q=C,\qquad \Delta_tq=t^2C,\qquad
 \langle q,C\rangle=\lVert\beta\rVert_{L^2}^2.
\end{equation}
In particular $PC=cq$, and the Rayleigh quotient of $q$ gives
$\lambda_1(t)\le ct^2$.

The zero eigenvalue of $\Delta_0$ is isolated and simple. The analytic
family \eqref{eq:two-operator} has compact resolvent and constant domain,
so its spectral projection near zero is analytic and of rank one for
small $|t|$ \cite[Chapter~VII]{Kato1995}. The remaining spectrum stays
bounded away from zero. For clarity, on a contour separating zero from
the positive spectrum, $(\Delta_0-\zeta)^{-1}$ maps $L^2$ to $H^2$,
and
$(tA_1+t^2A_2)(\Delta_0-\zeta)^{-1}$ has norm $O(|t|)$ on $L^2$.
A Neumann series constructs the perturbed resolvent and its contour
integral. The same argument between $H^s$ and $H^{s+2}$ gives Taylor
expansions of the eigensection in every $C^k$ norm.

Choose the analytic eigensection $v(t)$ with
$\langle q,v(t)\rangle=V_q$ and write
\[
 v(t)=q+tv_1+t^2v_2+t^3v_3+\cdots,\qquad v_j\perp q.
\]
The order-$t$ equation gives $\lambda'(0)=0$ and $v_1=0$. At order
$t^2$, projection onto $q$ and $q^\perp$ gives respectively the
coefficient $c$ and
\[
 \Delta_0v_2+C=cq,\qquad v_2=-Rz.
\]
The order-$t^3$ coefficient vanishes because $A_1q=0$. If
$\lambda_{[4]}$ denotes the coefficient of $t^4$, then
\[
 V_q\lambda_{[4]}
 =\langle q,A_1v_3+A_2v_2\rangle
 =-\langle z,Rz\rangle.
\]
This coefficient is $-d$, with $d\ge0$ because $R$ is positive.
Conjugation by $J=\diag(-\id_F,\id_G)$ is unitary and intertwines
$D_t$ with $D_{-t}$. Uniqueness of the isolated branch makes its
eigenvalue even in $t$, proving the $O(t^6)$ remainder in
\eqref{eq:two-expansion}. Finally,
$\lVert v(t)\rVert^2=V_q+O(t^4)$, so normalization gives
\eqref{eq:two-section}.

The Chern connection is $\nabla_t=\nabla_0+t(N-N^*)$, where the adjoint
also changes the form type. The K\"ahler identity
$\db^*\beta=-\sqrt{-1}\Lambda_\omega\partial\beta=0$ and its
adjoint kill the linear term in the contracted curvature. In a unitary
coframe with $\omega=\sqrt{-1}\sum\theta^\nu\wedge\bar\theta^\nu$,
the quadratic term is
\[
 \sqrt{-1}\Lambda_\omega(N-N^*)\wedge(N-N^*)=C.
\]
Its trace is zero, proving \eqref{eq:two-curvature}.
\end{proof}

To view this degeneration on the fixed holomorphic bundle $E_1$, set
$g_t=\diag(t\id_F,\id_G)$ for $t>0$. Then
\begin{equation}\label{eq:two-gauge}
 \db_t=g_t\db_1g_t^{-1},\qquad
 h_t=g_t^*h_0=t^2h_F\oplus h_G.
\end{equation}
The map $g_t:(E_1,h_t)\to(E_t,h_0)$ is a holomorphic isometry and
identifies the endomorphism Laplacians unitarily. Thus $h_t$ is an
approximate Hermitian--Einstein family on $E_1$. The convergence in
\eqref{eq:two-section} is in the gauge $g_t$ of \eqref{eq:two-gauge}. Its limit $u_0$
has two distinct fiberwise eigenvalues, and
\[
 P_F=\sqrt{V_q}\,u_0+\frac ar\id,\qquad P_G=\id-P_F.
\]


\section{Several stable summands on a curve}\label{sec:several}

Let $(X,\omega)$ now be a compact connected Riemann surface. Fix pairwise
nonisomorphic stable bundles $F_1,\ldots,F_m$, $m\ge2$, of equal slope,
with ranks $r_i$ and Hermitian--Einstein metrics $h_i$. On
$E_0=\bigoplus_iF_i$ use the metric $h_0=\bigoplus_i h_i$.
Choose strictly decreasing integers $\ell_1>\cdots>\ell_m=0$, the
\emph{scaling exponents}, and harmonic forms
\[
 \beta_{ij}\in\Omega^{0,1}(\Hom(F_j,F_i)),\qquad i<j.
\]
Some forms may vanish, but at least one is nonzero. Set
\begin{equation}\label{eq:several-family}
 p_{ij}=\ell_i-\ell_j,\qquad
 N_t=\sum_{i<j}t^{p_{ij}}\beta_{ij},\qquad
 \db_t=\db_0+N_t,\qquad t>0,
\end{equation}
where each summand occupies its indicated upper-triangular block.
Integrability is automatic on a curve. The diagonal transformation
$g_t=\diag(t^{\ell_i}\id_{F_i})$ satisfies
$\db_t=g_t\db_1g_t^{-1}$. Thus all $E_t$, $t>0$, are holomorphically
isomorphic, and $g_t$ is an isometry from $(E_1,g_t^*h_0)$ to $(E_t,h_0)$.

On the fixed Hilbert space $H=L^2(X,\End_0 E_0)$ define
\[
 D_t=D_0+B_t,\qquad D_0=\db_{0,\End},\qquad
 B_tu=[N_t,u],\qquad \Delta_t=D_t^*D_t.
\]
The form domain is $H^1$ and the operator domain is $H^2$. Stability
and nonisomorphism give
\begin{equation}\label{eq:split-kernel}
 \mathcal K=\Ker D_0
 =\left\{\diag(x_i\id_{F_i}):\sum_i r_ix_i=0\right\},
 \qquad \dim_{\CC}\mathcal K=m-1.
\end{equation}
Let $P:H\to\mathcal K$ be the orthogonal projection, $Q=\id-P$, and
let $\gamma>0$ be the lowest eigenvalue of $\Delta_0$ on $QH$. Put
\[
 \eta_t=\lVert B_t\rVert_{L^2\to L^2(\Omega^{0,1})},\qquad
 p_* =\min_{\beta_{ij}\ne0}p_{ij}.
\]
The coefficients of $B_t$ are $O(t^{p_*})$ in every differentiability
norm; in particular $\eta_t=O(t^{p_*})$.

\subsection{The weighted graph and the relative estimate}

Associate to \eqref{eq:several-family} the undirected graph with vertices
$1,\ldots,m$ and edges $ij$ for which $\beta_{ij}\ne0$. Define its vertex
masses and conductances by
\begin{equation}\label{eq:graph-data}
 M_i=Vr_i,\qquad c_{ij}=\lVert\beta_{ij}\rVert_{L^2}^2,
 \qquad k_{ij}(t)=c_{ij}t^{2p_{ij}}.
\end{equation}
Absent conductances are zero. The weighted graph Laplacian $L_t$ on
$\mathcal K$ is defined by
\begin{equation}\label{eq:graph-form}
 A_t(x):=\langle x,L_tx\rangle
 =\sum_{i<j}k_{ij}(t)|x_i-x_j|^2,
 \qquad \lVert x\rVert^2=\sum_iM_i|x_i|^2.
\end{equation}
Equivalently, $(L_tx)_i=M_i^{-1}\sum_{j\ne i}k_{ij}(t)(x_i-x_j)$,
restricted to $\sum_iM_ix_i=0$. Write
$\nu_1(t)\le\cdots\le\nu_{m-1}(t)$ for its eigenvalues and
$\lambda_1(t)\le\lambda_2(t)\le\cdots$ for those of $\Delta_t$ on $H$.
Both lists count multiplicity and include zero.

\begin{theorem}\label{thm:relative}
If $\eta_t\le\sqrt\gamma/4$ and $\gamma_t=(\sqrt\gamma-\eta_t)^2$,
then
\begin{equation}\label{eq:relative}
 \left(1-\frac{2\eta_t^2}{\gamma_t}\right)\nu_j(t)
 \le\lambda_j(t)\le\nu_j(t),\qquad 1\le j\le m-1,
 \qquad \lambda_m(t)\ge\frac\gamma4.
\end{equation}
There are exactly $m-1$ eigenvalues below $\gamma/4$, including zeros.
For every positive graph eigenvalue,
\begin{equation}\label{eq:relative-ratio}
 \frac{\lambda_j(t)}{\nu_j(t)}=1+O(t^{2p_*}).
\end{equation}
The error is uniform in $j$ for this fixed family, regardless of the
ratios between its positive graph eigenvalues. If the graph has $c$
connected components, then $\Delta_t$ has exactly $c-1$ zero eigenvalues
on trace-free endomorphisms for every $t>0$. In particular, $E_t$ is
simple if and only if the graph is connected.
\end{theorem}

\begin{proof}
For $x\in\mathcal K$, the $ij$ block of $B_tx$ is
$t^{p_{ij}}(x_j-x_i)\beta_{ij}$. The different blocks are orthogonal, so
\[
 \lVert B_tx\rVert^2=A_t(x).
\]
Harmonicity gives the cancellation
\begin{equation}\label{eq:harmonic-cancellation}
 D_0^*B_tx=0\qquad(x\in\mathcal K).
\end{equation}
For $y\in QH^1$, the split spectral gap gives
\begin{equation}\label{eq:complement-gap}
 \lVert D_ty\rVert\ge\lVert D_0y\rVert-\lVert B_ty\rVert
 \ge(\sqrt\gamma-\eta_t)\lVert y\rVert=\sqrt{\gamma_t}\,\lVert y\rVert.
\end{equation}
Decompose $u=x+y$ with $x\in\mathcal K$, $y\in QH^1$. By
\eqref{eq:harmonic-cancellation},
\[
 \lVert D_tu\rVert^2=A_t(x)
  +2\operatorname{Re}\langle B_tx,B_ty\rangle+\lVert D_ty\rVert^2.
\]
The cross term has absolute value at most
$2\eta_t A_t(x)^{1/2}\lVert y\rVert$. Young's inequality therefore gives
\begin{equation}\label{eq:form-lower}
 \lVert D_tu\rVert^2\ge
 \left(1-\frac{2\eta_t^2}{\gamma_t}\right)A_t(x)+\frac{\gamma_t}2\lVert y\rVert^2.
\end{equation}

The min--max principle on subspaces of $\mathcal K$ gives the upper
bound in \eqref{eq:relative}. For the lower bound, compare with the
orthogonal sum of $(1-2\eta_t^2/\gamma_t)L_t$ on $\mathcal K$ and
$(\gamma_t/2)\id$ on $QH$. Its variational values below $\gamma_t/2$ are the graph
eigenvalues with the indicated prefactor. Since
\[
 \nu_{m-1}(t)\le\eta_t^2\le\gamma/16,
 \qquad \gamma_t/2\ge9\gamma/32,
\]
min--max proves \eqref{eq:relative}, including
$\lambda_m(t)\ge \gamma_t/2>\gamma/4$. This also proves
\eqref{eq:relative-ratio}; its constants may depend on the fixed bundles
and forms.

The kernel of the graph form consists of vectors constant on each
connected component. Its intersection with $\mathcal K$ has dimension
$c-1$. Formula \eqref{eq:form-lower} shows that these are exactly the
zero modes of $\Delta_t$ for small $t$. The positive-parameter complex
gauge equivalence gives the same kernel dimension for all $t>0$.
Adding the scalar identity proves the assertion about simplicity.
\end{proof}

\begin{remark}\label{rem:harmonicity-needed}
The cancellation \eqref{eq:harmonic-cancellation} is needed. Without it,
complementary directions can change the leading term of the form restricted
to $\mathcal K$. Already
in finite dimensions, $D_0=(0\ \ 1)$ and $B_t=(t\ \ 0)$ have restricted
energy $t^2$ on $\Ker D_0$, while $(D_0+B_t)^*(D_0+B_t)$ has $(-1,t)$ in its
kernel. So the relative comparison does not follow from smallness of $B_t$
alone.
\end{remark}

\subsection{The complementary equation and smooth convergence}

On $QH$ let $T_t=Q\Delta_tQ$, with domain $QH^2$, be the operator of
the restricted closed quadratic form. It satisfies $T_t\ge \gamma_t$, where
$\gamma_t$ is as in Theorem~\ref{thm:relative}. By \eqref{eq:harmonic-cancellation}, the
off-diagonal block is
\begin{equation}\label{eq:off-diagonal}
 R_t=Q\Delta_tP=QB_t^*B_tP,
 \qquad \lVert R_tx\rVert\le\eta_t A_t(x)^{1/2}.
\end{equation}
For $0\le\lambda<\gamma_t/2$, solving the $Q$ component of the eigenvalue
equation gives
\begin{equation}\label{eq:schur}
 y=-(T_t-\lambda)^{-1}R_tx,\qquad
 S_t(\lambda)x=\lambda x,\qquad
 S_t(\lambda)=L_t-R_t^*(T_t-\lambda)^{-1}R_t.
\end{equation}
The correction in this Schur complement satisfies the relative form bound
\begin{equation}\label{eq:schur-bound}
 0\le L_t-S_t(\lambda)\le\frac{2\eta_t^2}{\gamma_t}L_t.
\end{equation}
Cauchy--Schwarz for the positive form on the left yields, for $x,z\in
\mathcal K$,
\begin{equation}\label{eq:polarized-schur}
 |\langle z,(L_t-S_t(\lambda))x\rangle|
 \le\frac{2\eta_t^2}{\gamma_t}A_t(z)^{1/2}A_t(x)^{1/2}.
\end{equation}

\begin{proposition}\label{prop:subspace}
For every $k\ge0$, a unit eigensection $u$ belonging to one of the
$m-1$ small eigenvalues $\lambda$ satisfies
\begin{equation}\label{eq:smooth-complement}
 \lVert Qu\rVert_{C^k}\le C_k t^{p_*}\sqrt\lambda.
\end{equation}
If $\Pi_t$ denotes the orthogonal projection onto the full small
spectral subspace, then
\begin{equation}\label{eq:full-projection}
 \lVert\Pi_t-P\rVert_{L^2\to C^k}=O(t^{2p_*}).
\end{equation}
\end{proposition}

\begin{proof}
Write $u=x+y$ as above and set $\epsilon_t=2\eta_t^2/\gamma_t$.
Formula \eqref{eq:form-lower} implies
$A_t(x)\le\lambda/(1-\epsilon_t)$. Equations
\eqref{eq:off-diagonal}--\eqref{eq:schur} first give
$\lVert y\rVert_{L^2}\le C\eta_t\sqrt\lambda$.

To obtain the same estimate in stronger norms, observe that the nonzero
$\beta_{ij}$ form a fixed finite collection in distinct orthogonal
blocks. Therefore
$\lVert B_tx\rVert_{H^s}\le C_s\lVert B_tx\rVert_{L^2}$ for
$x\in\mathcal K$, uniformly in its scalar coefficients and in $t$.
The coefficients of $B_t$ and all their derivatives are $O(t^{p_*})$,
so
\[
 \lVert R_tx\rVert_{H^s}\le C_s t^{p_*}A_t(x)^{1/2}.
\]
Uniform elliptic estimates, together with the $L^2$ inverse bound,
give uniform $H^s\to H^{s+2}$ bounds for $(T_t-\lambda)^{-1}$ when
$0\le\lambda\le\gamma/16$. The projections $P,Q$ introduce only
smooth finite-rank terms. Apply these estimates to \eqref{eq:schur}
and then use Sobolev embedding to obtain \eqref{eq:smooth-complement}.

The largest small eigenvalue is $O(t^{2p_*})$. Applying
\eqref{eq:smooth-complement} to an orthonormal eigenbasis bounds $Q$
on $\operatorname{ran}\Pi_t$, as a map to $C^k$, by $O(t^{2p_*})$.
The spaces $\operatorname{ran}\Pi_t$ and $\mathcal K$ have the same
finite dimension. Thus $P$ restricts to an isomorphism between them for
small $t$, and $\operatorname{ran}\Pi_t$ is the graph of a map
$\mathcal K\to QH$ of $C^k$ norm $O(t^{2p_*})$. The formula for
the orthogonal projection onto this graph proves
\eqref{eq:full-projection}.
\end{proof}

After adjoining the scalar identity, the limiting subspace is the
algebra $\bigoplus_i\CC\id_{F_i}$. Its minimal nonzero idempotents are
the projections onto the $F_i$, and hence determine the unordered
summands of the split limit $E_0$. The small spectral subspace at
positive $t$ need not be an algebra.

\subsection{Mean curvature of the family}

The common Einstein constant of the $h_i$ is $\kappa$. The Chern
connection is $\nabla_t=\nabla_0+N_t-N_t^*$. Harmonicity and the
K\"ahler identities eliminate the linear term in its mean curvature,
so
\begin{equation}\label{eq:several-curvature}
 K_t=\kappa\id+\sqrt{-1}\Lambda_\omega
       (N_t-N_t^*)\wedge(N_t-N_t^*),
 \qquad K_t-\kappa\id=O_{C^k}(t^{2p_*}).
\end{equation}
In a unitary coframe on the curve, write $N_t=n_t\bar\theta$.
The quadratic term is $[n_t,n_t^*]$, so it is trace-free; it may have
off-diagonal blocks. Thus $g_t^*h_0$ is an approximate
Hermitian--Einstein family on the fixed bundle $E_1$. The
upper-triangular filtration, whose quotients are stable of the same
slope, proves semistability by degree additivity. Its first summand
excludes stability. Theorem~\ref{thm:relative} supplies the precise
spectral rates for this family.

\section{Decay orders and limiting spectral subspaces}\label{sec:orders}

The graph comparison reduces every leading decay order to a finite
construction. Let $p^{(1)}<\cdots<p^{(s)}$ be the distinct exponents
$p_{ij}$ on nonzero edges. For such an exponent $p$, let
$\mathcal C_{<p}$ be the connected components of the graph that
retains only edges with $p_{ij}<p$, and define
$\mathcal C_{\le p}$ similarly. Isolated vertices count as
components. Define a quotient graph on the vertices
$A\in\mathcal C_{<p}$ by the masses and conductances
\begin{equation}\label{eq:quotient-data}
 M_A=\sum_{i\in A}M_i,\qquad
 c^{(p)}_{AB}
 =\sum_{\substack{ij\text{ joins }A\text{ to }B\\p_{ij}=p}}c_{ij},
 \qquad A\ne B.
\end{equation}
Each original unordered edge is counted once. Internal edges make no
contribution. Let $L^{(p)}$ be the weighted Laplacian of this
quotient graph, and list its positive eigenvalues as
$\rho_{p,1},\ldots,\rho_{p,d_p}$, counting multiplicity.
Their number is
\[
 d_p=|\mathcal C_{<p}|-|\mathcal C_{\le p}|.
\]
The quotient vertex space embeds isometrically in the original vertex
space as the vectors constant on each $A\in\mathcal C_{<p}$.

\begin{theorem}\label{thm:orders}
For each $p$, exactly $d_p$ eigenvalues of $\Delta_t$ have
the asymptotics
\begin{equation}\label{eq:all-orders}
 \lambda(t)=t^{2p}(\rho_{p,k}+o(1)),
 \qquad 1\le k\le d_p.
\end{equation}
Together with the $c-1$ identically zero eigenvalues, these account for
all $m-1$ small eigenvalues. For each fixed positive eigenvalue $\rho$
of $L^{(p)}$, the corresponding spectral subspace converges to
the $\rho$-eigenspace, embedded as block-scalar endomorphisms constant
on the components in $\mathcal C_{<p}$. More precisely, the
orthogonal projections converge in $L^2\to C^k$ norm for every $k$.
For a repeated coefficient the assertion concerns the entire subspace.
\end{theorem}

\begin{proof}
First work on the full graph vertex space, including its scalar zero
mode, with the mass inner product. Let $W_p$ consist of the vectors
constant on $\mathcal C_{<p}$. All terms with exponent less than
$p$ vanish on $W_p$ and, by self-adjointness, have zero cross
blocks relative to $W_p\oplus W_p^\perp$. On
$W_p^\perp$, their sum divided by $t^{2p}$ tends uniformly
to positive infinity on the unit sphere. Indeed, if stronger edges
exist and $p_-<p$ is their largest exponent, the graph gaps
inside their connected components give a lower bound
$c't^{2(p_- -p)}$ on that complement.

After division by $t^{2p}$, the terms with exponent at least
$p$ have bounded cross blocks. Their restriction to $W_p$
converges to $L^{(p)}$, with masses exactly as in
\eqref{eq:quotient-data}. For a bounded spectral parameter, elimination
of $W_p^\perp$ introduces an error tending to zero, because the
inverse of its diagonal block tends to zero. The finite-dimensional
min--max principle then shows that the $\dim W_p$ bounded
eigenvalues of $t^{-2p}L_t$ converge to those of $L^{(p)}$;
all remaining eigenvalues tend to infinity. If no stronger edge exists,
this is ordinary matrix convergence, with no complementary block.

The positive quotient eigenvalues account for exactly $d_p$
modes at this scale. The component counts telescope to $m-c$ over all
exponents, leaving $c$ exact graph zero modes. Remove the scalar mode
and apply the relative estimate \eqref{eq:relative} to obtain
\eqref{eq:all-orders} for $\Delta_t$.

Now let $u_t=x_t+y_t$ be a unit eigensection with
$\lambda(t)/t^{2p}\to\rho>0$. By \eqref{eq:form-lower},
$A_t(x_t)=O(t^{2p})$. Every stronger edge therefore satisfies
\[
 |(x_t)_i-(x_t)_j|=O(t^{p-p_{ij}}).
\]
Every limit of $x_t$ lies in $W_p\cap\mathcal K$, and
$y_t\to0$ in every $C^k$ norm by \eqref{eq:smooth-complement}.
Test the second equation in \eqref{eq:schur} against a fixed
$z\in W_p\cap\mathcal K$. Its energy is $O(t^{2p})$, so
\eqref{eq:polarized-schur} bounds the correction by
$O(t^{2p+2p_*})$. After division by $t^{2p}$ this correction
vanishes. The stronger graph terms vanish on $z$, the terms with
exponent greater than $p$ vanish in the limit, and the remaining
terms give
\[
 L^{(p)}x_0=\rho x_0.
\]
This equality on the full quotient space follows from the tested
identity on its mean-zero subspace, since both sides have mean zero.

Choose a small interval about $\rho$ containing no other quotient
eigenvalue and not containing zero. The eigenvalue count above gives
the dimension of the spectral subspace for which
$\lambda/t^{2p}$ lies in this interval. Limits of orthonormal
eigensections remain orthonormal, because their $Q$ components tend to
zero. They therefore span the entire $\rho$-eigenspace. Compactness
in the fixed finite-dimensional space $\mathcal K$, together with
\eqref{eq:smooth-complement}, gives convergence of the corresponding
orthogonal projections in every $L^2\to C^k$ norm.
\end{proof}

\section{Examples with unequal decay orders}\label{sec:examples}

\subsection{Three factors}

Take $m=3$, $\ell=(p+q,q,0)$, and integers $0<p<q$. Assume
$\beta_{12},\beta_{23}\ne0$; the entry $\beta_{13}$ may vanish. Set
\[
 A=\lVert\beta_{12}\rVert_{L^2}^2,\qquad
 B=\lVert\beta_{23}\rVert_{L^2}^2,\qquad
 C=\lVert\beta_{13}\rVert_{L^2}^2.
\]
Thus $k_{12}=At^{2p}$, $k_{23}=Bt^{2q}$, and
$k_{13}=Ct^{2(p+q)}$. The two positive graph eigenvalues are the roots
of $\nu^2-T\nu+J=0$, where
\begin{align}
 T={}&k_{12}(M_1^{-1}+M_2^{-1})
      +k_{23}(M_2^{-1}+M_3^{-1})
      +k_{13}(M_1^{-1}+M_3^{-1}),\label{eq:three-trace}\\
 J={}&\frac{M_1+M_2+M_3}{M_1M_2M_3}
       (k_{12}k_{23}+k_{12}k_{13}+k_{23}k_{13}).\label{eq:three-product}
\end{align}
These identities are the trace and sum of principal two-by-two minors
of the mass-normalized graph Laplacian. With
$\sigma=2\min\{p,q-p\}>0$, Theorem~\ref{thm:relative} gives
\begin{align}
 \lambda_1(t)&=B\left(\frac1{M_1+M_2}+\frac1{M_3}\right)
      t^{2q}\bigl(1+O(t^\sigma)\bigr),\label{eq:three-slow}\\
 \lambda_2(t)&=A\left(\frac1{M_1}+\frac1{M_2}\right)
      t^{2p}\bigl(1+O(t^\sigma)\bigr),
 \qquad\lambda_3(t)\ge\gamma/4.\label{eq:three-fast}
\end{align}
To check the error terms directly, put $\eta=t^{2(q-p)}$ and
$\vartheta=t^{2p}$. The leading terms of $T,J$ have relative errors
$O(\eta)$ and $O(\vartheta)$ respectively, and $J/T^2=O(\eta)$.
The larger root is $T(1+O(\eta))$ and the smaller is
$(J/T)(1+O(\eta))$. The operator comparison adds $O(\vartheta)$.
At the $t^{2q}$ scale, the stronger $12$ edge has combined vertices
$1,2$, explaining the mass $M_1+M_2$ in \eqref{eq:three-slow}.

Write $r_{12}=r_1+r_2$ and $r=r_{12}+r_3$. The limiting eigensections
are the normalized versions of
\begin{align*}
 q_{\rm slow}&=\diag\left(\frac{r_3}r\id_{F_1},
                  \frac{r_3}r\id_{F_2},-\frac{r_{12}}r\id_{F_3}\right),
 &\lVert q_{\rm slow}\rVert^2&=\frac{Vr_{12}r_3}r,\\
 q_{\rm fast}&=\diag\left(\frac{r_2}{r_{12}}\id_{F_1},
                  -\frac{r_1}{r_{12}}\id_{F_2},0\right),
 &\lVert q_{\rm fast}\rVert^2&=\frac{Vr_1r_2}{r_{12}}.
\end{align*}
Theorem~\ref{thm:orders} gives smooth convergence after fixing the
phases by positive inner products with these respective limits.
The stronger extension separates the internal $F_1,F_2$ mode at the
larger eigenvalue; the weaker extension separates their combined block
from $F_3$ at the smaller eigenvalue.

For three line bundles, $V=1$, $A=B=1$, and $(p,q)=(1,2)$, the
formulas become
\begin{equation}\label{eq:three-lines}
 \lambda_1(t)=\tfrac32t^4(1+O(t^2)),\qquad
 \lambda_2(t)=2t^2(1+O(t^2)).
\end{equation}
These choices are realized geometrically on a curve of genus at least
two. Choose three pairwise nonisomorphic degree-zero line bundles.
Every nontrivial degree-zero Hom line bundle has $h^1=g-1>0$, by the
Riemann--Roch argument in Corollary~\ref{cor:sharpness}. Choose harmonic
forms of unit $L^2$ norm for the consecutive edges and set
$\beta_{13}=0$. The resulting bundle is simple and strictly semistable.

\subsection{Five factors and a repeated coefficient}

Take five pairwise nonisomorphic degree-zero line bundles on such a
curve, with $V=1$ and $\ell=(7,6,3,1,0)$. Choose unit-norm forms on the
path edges $12,23,34,45$, whose exponents are $1,3,2,1$. One may also
include a form of norm $2$ on edge $15$, of exponent $7$. All remaining
entries vanish. The quotient graphs have the following positive spectra:
\[
\begin{array}{c|c|c}
 p&\text{components before adding edges of exponent }p
        &\text{positive coefficients}\\ \hline
 1&\{1\},\{2\},\{3\},\{4\},\{5\}&2,\ 2\\
 2&\{1,2\},\{3\},\{4,5\}&3/2\\
 3&\{1,2\},\{3,4,5\}&5/6\\
 7\text{ (if present)}&\{1,2,3,4,5\}&\text{none}
\end{array}
\]
At exponent $2$, edge $34$ joins masses $1,2$, giving $1+1/2$.
At exponent $3$, edge $23$ joins masses $2,3$, giving $1/2+1/3$.
The optional edge of exponent $7$ is internal to a component already
connected by stronger edges. Consequently
\begin{equation}\label{eq:five-lines}
 \lambda_1(t)\sim\tfrac56t^6,\qquad
 \lambda_2(t)\sim\tfrac32t^4,\qquad
 \lambda_3(t)\sim2t^2,\qquad\lambda_4(t)\sim2t^2.
\end{equation}
For the last two eigenvalues, the limiting two-dimensional subspace is
spanned by the internal differences on $F_1\oplus F_2$ and
$F_4\oplus F_5$. Theorem~\ref{thm:orders} specifies this subspace,
without choosing individual eigenvectors within it.

\section{A harmonic normal form on curves}\label{sec:normal-form}

Every filtered bundle on a curve can be put in a form in which all
extension entries are harmonic. The normal form follows from Dolbeault
Hodge theory and a finite sequence of upper-triangular gauge
transformations. The condition $\db_0^*\beta=0$ is used in the
scaled-extension construction of Sektnan--Tipler
\cite[Section~5.1, equation~(28)]{SektnanTipler2024}; on a curve it makes
every extension entry harmonic. We give the triangular construction
explicitly to account for the interactions between successive extensions.


Let $X$ be a compact connected Riemann surface with a fixed K\"ahler
form. Suppose a holomorphic bundle $E$ has a filtration by holomorphic
subbundles
\begin{equation}\label{eq:curve-filtration}
 0=E^{(0)}\subset E^{(1)}\subset\cdots\subset E^{(m)}=E,
 \qquad F_i=E^{(i)}/E^{(i-1)}.
\end{equation}
Put $G=\bigoplus_i F_i$, let $\db_0$ be its split Dolbeault operator,
and fix Hermitian metrics on the $F_i$. Harmonicity of a
$\Hom(F_j,F_i)$-valued form always refers to these metrics and the
fixed K\"ahler form.

\begin{lemma}[Harmonic upper-triangular form]\label{lem:normal-form}
There is a smooth isomorphism $\Phi:G\to E$ preserving the
filtration and inducing the identity on each quotient $F_i$, such that
\begin{equation}\label{eq:harmonic-normal-form}
 \Phi^{-1}\db_E\Phi=\db_0+\beta,\qquad
 \beta=\sum_{i<j}\beta_{ij},\qquad
 \db_0\beta_{ij}=\db_0^*\beta_{ij}=0.
\end{equation}
No stability hypothesis or smallness assumption on the original
extension entries is needed for this normal form.
\end{lemma}

\begin{proof}
Choose a smooth splitting $\sigma:G\to E$ of the filtration,
inducing the identity on its quotients. The pulled-back operator has
the form $\sigma^{-1}\db_E\sigma=\db_0+A$, where $A$ is strictly
upper triangular. For $\ell\ge1$, define the subbundle
\[
 \mathfrak n^\ell
 =\bigoplus_{j-i\ge\ell}\Hom(F_j,F_i)\subset\End G.
\]
Then $\mathfrak n^a\mathfrak n^b\subset\mathfrak n^{a+b}$,
$\mathfrak n^m=0$, and $\db_0$ preserves these subbundles.

Suppose that the entries of a current upper-triangular form $A$
with $j-i<\ell$ have already been made harmonic. Every smooth
$\Hom(F_j,F_i)$-valued $(0,1)$-form on a curve has a Hodge
decomposition
\begin{equation}\label{eq:curve-hodge}
 A_{ij}=H_{ij}A_{ij}+\db_0v_{ij},
\end{equation}
where $H_{ij}$ is harmonic projection and $v_{ij}$ is a smooth
section. Indeed, the $(0,2)$-term in the Dolbeault Hodge decomposition
vanishes in complex dimension one. Smoothness follows from elliptic
regularity. Choose $u\in\Gamma(\mathfrak n^\ell)$ with entries
$u_{ij}=-v_{ij}$ when $j-i=\ell$ and zero elsewhere, and set
$U_\ell=\id+u$. This is a globally invertible smooth transformation:
its inverse is the finite sum $\id-u+u^2-\cdots$.

The transformed extension form is exactly
\[
 A'=U_\ell^{-1}AU_\ell+U_\ell^{-1}\db_0U_\ell.
\]
Since $A\in\Omega^{0,1}(\mathfrak n^1)$ and
$u\in\Gamma(\mathfrak n^\ell)$, the multiplication rule gives
\begin{equation}\label{eq:triangular-gauge}
 A'\equiv A+\db_0u
 \pmod{\Omega^{0,1}(\mathfrak n^{\ell+1})}.
\end{equation}
For the derivative term, we use $2\ell\ge\ell+1$.
Thus entries with $j-i<\ell$ are unchanged, and each entry with
$j-i=\ell$ becomes $H_{ij}A_{ij}$. Changes to higher entries are
retained for the subsequent steps.

Starting with $\ell=1$ and ending with $\ell=m-1$ makes every
entry harmonic. If $U=U_1\cdots U_{m-1}$ is the product of the
successive transformations, then $\Phi=\sigma U$ has all the
asserted properties. Each transformation is the identity on the
quotients, and the process terminates after finitely many steps.
\end{proof}

The higher entries in \eqref{eq:harmonic-normal-form} need not be
the harmonic projections of the original entries of $A$. For three
blocks, let the first correction have entries $u_{12},u_{23}$ and
zero $13$ entry, so that
$\beta_{12}=A_{12}+\db_0u_{12}$ and
$\beta_{23}=A_{23}+\db_0u_{23}$ are harmonic. Before the second
correction, the $13$ entry is
\begin{equation}\label{eq:three-normal-form}
 A'_{13}=A_{13}+A_{12}u_{23}-u_{12}\beta_{23}.
\end{equation}
It is this modified entry whose harmonic part defines $\beta_{13}$.
The induction therefore includes the nonlinear interaction of the
successive extensions.


\section{The flow for two factors: the model case}\label{sec:flow}

This section treats a nonsplit extension of two stable bundles directly.
It introduces the three tools used in general in
Sections~\ref{sec:diagonal-flow}--\ref{sec:complete-spectrum}: comparison
of metrics along the flow, a corrected approximate solution with integrable
error, and matching of the flow with the extension orbit at late times.

Let $X$ be a compact connected Riemann surface, and retain $F,G$, the
nonzero harmonic extension form $\beta$, and the notation $q,V_q,c,C,z,R$
of Section~\ref{sec:two-factor}. Thus $E=E_1$ is the nonsplit extension
of the nonisomorphic stable equal-slope bundles $G$ by $F$. We use flow
time $\tau$ to distinguish it from the scaling parameter of the earlier
sections. Starting from any smooth Hermitian metric on $E$, consider
the normalized Hermitian Yang--Mills flow
\begin{equation}\label{eq:metric-flow}
 h^{-1}\partial_\tau h=-2(K_h-\kappa\id),\qquad
 \kappa=\frac{2\pi\mu(E)}{V}.
\end{equation}
The factor $2$ fixes the time normalization in the result below.

We use the classical convergence theorem for Yang--Mills flow on
a curve: the solution exists for all time, and in unitary gauges its
Chern connections converge to the Hermitian--Einstein connection on
the graded bundle $F\oplus G$
\cite{Daskalopoulos1992,Rade1992}. For smooth initial data the
convergence, modulo gauge, is smooth; the higher regularity follows
from parabolic regularity. In metric language there are holomorphic
isometries from $(E,h(\tau))$ to $(F\oplus G,h_0,\db^{(\tau)})$
with $\db^{(\tau)}\to\db_0$ smoothly. This analytic convergence
and the identification of the graded limit are inputs here.

The metric asymptotics of this flow have been studied much more
generally by Haiden--Katzarkov--Kontsevich--Pandit
\cite{HKKP2018}. Their construction uses harmonic extension forms,
a Green operator correction, and the order-preserving property of
the metric flow. We use these ingredients in the two-factor case,
together with a comparison begun at arbitrarily late times, to
determine the leading coefficient of the small Dolbeault eigenvalue.

\begin{theorem}\label{thm:flow-two-factor}
For every smooth initial metric in \eqref{eq:metric-flow}, exactly one
eigenvalue on trace-free endomorphisms tends to zero, and
\begin{equation}\label{eq:flow-rate}
 \lambda_1(h(\tau))=\frac{1+o(1)}{2\tau},\qquad
 \lambda_2(h(\tau))\ge\delta>0
\end{equation}
for sufficiently large $\tau$.

In the unitary identifications supplied by smooth convergence of the
connections, the first spectral projection converges in every
$L^2(h_0)\to C^k(h_0)$ norm to the projection onto $\CC q$.
Equivalently, the unit first eigensection, with its phase fixed by
positive inner product with $q$, tends smoothly to $q/\sqrt{V_q}$.
Its limit recovers the two summand projections by
\[
 P_F=\sqrt{V_q}\,u_\infty+\frac ar\id,\qquad P_G=\id-P_F.
\]
\end{theorem}

The proof separates metric comparison, a corrected approximate
solution, and a local description of the fixed extension orbit.

\subsection{Order comparison and the endomorphism spectrum}

If two metrics $h_1,h_2$ on a fixed holomorphic bundle satisfy
$e^{-\epsilon}h_1\le h_2\le e^\epsilon h_1$, then their induced
norms on endomorphisms satisfy
\[
 e^{-2\epsilon}|u|_{h_1}^2
 \le |u|_{h_2}^2\le e^{2\epsilon}|u|_{h_1}^2.
\]
Indeed, in an $h_1$-orthonormal frame diagonalizing $h_2$, the
coefficient of $|u_{ij}|^2$ is the ratio of its $i$th and $j$th
eigenvalues. The same bounds hold for endomorphism-valued $(0,1)$-forms,
since the base metric is fixed. The Dolbeault operator and its
trace-free domain are independent of the bundle metric, so min--max
gives
\begin{equation}\label{eq:metric-spectral-comparison}
 e^{-4\epsilon}\lambda_j(h_1)
 \le\lambda_j(h_2)\le e^{4\epsilon}\lambda_j(h_1).
\end{equation}
A positive scalar multiple of either metric leaves these eigenvalues
unchanged.

The order comparison for the metric flow has a useful version with
an error term. Suppose a smooth family $\widehat h(\tau)$ has residual
\begin{equation}\label{eq:flow-residual}
 \mathcal E_{\widehat h}
 :=\widehat h^{-1}\partial_\tau\widehat h
       +2(K_{\widehat h}-\kappa\id),\qquad
 \lVert\mathcal E_{\widehat h}\rVert_{L^\infty,\op,\widehat h}
 \le f(\tau).
\end{equation}
If $h$ is an exact solution and
$e^{-\epsilon}\widehat h(T)\le h(T)\le e^\epsilon\widehat h(T)$,
then
\begin{equation}\label{eq:flow-barriers}
 e^{-\epsilon-\int_T^\tau f(s)\,ds}\widehat h(\tau)
 \le h(\tau)\le
 e^{\epsilon+\int_T^\tau f(s)\,ds}\widehat h(\tau).
\end{equation}
To see this, multiply $\widehat h$ by
$e^{\pm\epsilon\pm\int_T^\tau f}$.
A spatially constant scalar does not change Chern curvature, and its
logarithmic derivative adds $\pm f\id$ to the residual. The two
metrics are therefore respectively a subsolution and a supersolution.
The maximum principle for Hermitian metrics gives the assertion;
this is the order comparison of
\cite[Section~4.2, Propositions~4.3 and~4.5]{HKKP2018}, with the
error measured in the normalization \eqref{eq:flow-residual}.

\subsection{An approximate solution with a small total error}

Use determinant-one diagonal scaling
\begin{equation}\label{eq:flow-diagonal}
 g_\rho=\rho^q
 =\diag(\rho^{b/r}\id_F,\rho^{-a/r}\id_G),\qquad
 h_\rho=g_\rho^*h_0.
\end{equation}
It still transports the extension entry to $\rho\beta$, and its
endomorphism spectrum is that of Theorem~\ref{thm:two-factor}.
The uncorrected curvature is $\kappa\id+\rho^2C$.
Its component $\rho^2z$ perpendicular to $q$ generally prevents
$h_\rho$ from following \eqref{eq:metric-flow} under a scalar
change of parameter.

Put
\begin{equation}\label{eq:flow-correction}
 v=Rz,\qquad k_\rho=\exp(-\tfrac12\rho^2v),\qquad
 \widehat g_\rho=k_\rho g_\rho,
 \qquad \widehat h_\rho=\widehat g_\rho^*h_0.
\end{equation}
The endomorphism $v$ is smooth, self-adjoint, trace-free, and block
diagonal, so it commutes with $q$. Block diagonality follows because
the split Laplacian preserves the diagonal blocks. Self-adjointness
follows from the K\"ahler identities: at a Hermitian--Einstein
connection the endomorphism Dolbeault Laplacian is half the connection
Laplacian and therefore preserves adjoints.

\begin{lemma}\label{lem:flow-approximation}
Let $0<\rho_0\le\rho_*$ for a fixed sufficiently small $\rho_*$,
and define, for $\tau\ge T$,
\begin{equation}\label{eq:flow-rho}
 \rho(\tau)=\bigl(\rho_0^{-2}+2c(\tau-T)\bigr)^{-1/2}.
\end{equation}
There is a constant $A$, independent of $T$ and $\rho_0$, such that
$\widehat h(\tau)=\widehat h_{\rho(\tau)}$ satisfies
\begin{equation}\label{eq:flow-error-bound}
 \lVert\mathcal E_{\widehat h}\rVert_{L^\infty,\op,\widehat h}
 \le A\rho(\tau)^3,
 \qquad
 \int_T^\infty A\rho(\tau)^3\,d\tau=\frac{A}{c}\rho_0.
\end{equation}
For this approximate solution,
\begin{equation}\label{eq:flow-approx-spectrum}
 \lambda_1(\widehat h_\rho)=c\rho^2(1+O(\rho^2)),\qquad
 \lim_{\tau\to\infty}\tau\lambda_1(\widehat h(\tau))=\tfrac12.
\end{equation}
\end{lemma}

\begin{proof}
Transport everything by $\widehat g_\rho$ to the fixed metric $h_0$.
At the split structure the derivative of mean curvature under the
metric change $h_0e^{\epsilon s}$ is $\Delta_0s$. For example, this
follows by differentiating the Chern curvature to obtain
$\sqrt{-1}\Lambda\db_0\partial_0s=\Delta_0s$.
Smooth Taylor expansion in $\rho$, with
$k_\rho=\id-\tfrac12\rho^2v+O_{C^k}(\rho^4)$, therefore gives
\begin{equation}\label{eq:flow-corrected-curvature}
 \widehat K_\rho
 :=\widehat g_\rho K_{\widehat h_\rho}\widehat g_\rho^{-1}
 =\kappa\id+\rho^2(C-\Delta_0v)+O_{C^k}(\rho^3)
 =\kappa\id+c\rho^2q+O_{C^k}(\rho^3).
\end{equation}
Terms involving both the change of metric, of order $\rho^2$,
and the extension, of order $\rho$, are included in the cubic
remainder. All estimates are uniform for $0<\rho\le\rho_*$.

Since $q$ and $v$ commute and $\dot\rho=-c\rho^3$, the logarithmic
derivative of the gauge is exactly
\[
 \dot{\widehat g}_\rho\widehat g_\rho^{-1}
 =-c\rho^2q+c\rho^4v.
\]
For a metric $g^*h_0$, its logarithmic time derivative, transported
by $g$, is $\dot g g^{-1}+(\dot g g^{-1})^*$.
Combining this with \eqref{eq:flow-corrected-curvature} cancels the
quadratic terms in \eqref{eq:flow-residual} and proves the cubic
bound. Integration using $\dot\rho=-c\rho^3$ gives the stated
total error.

Also $e^{-B\rho^2}h_\rho\le\widehat h_\rho\le e^{B\rho^2}h_\rho$
for some fixed $B$. Apply \eqref{eq:metric-spectral-comparison}
and Theorem~\ref{thm:two-factor} to get
$\lambda_1(\widehat h_\rho)=c\rho^2(1+O(\rho^2))$.
Equation~\eqref{eq:flow-rho} proves the limit.
\end{proof}

\subsection{Matching the actual flow at late times}

The following observation strengthens bounded comparison to comparison
with a factor tending to one. It uses the fact that the extension has
two nonisomorphic stable factors.

\begin{lemma}\label{lem:flow-matching}
For the solution $h(T)$ of \eqref{eq:metric-flow}, there are
$\rho_T>0$, $A_T>0$, and $\epsilon_T\ge0$, with
$\rho_T\to0$ and $\epsilon_T\to0$, such that
\begin{equation}\label{eq:flow-matching}
 e^{-\epsilon_T}A_T\widehat h_{\rho_T}
 \le h(T)\le e^{\epsilon_T}A_T\widehat h_{\rho_T}.
\end{equation}
\end{lemma}

\begin{proof}
Use the smooth convergence of the flow in unitary gauges, and write
$D_T=\db^{(T)}\to\db_0$ on $(F\oplus G,h_0)$, with each
$(F\oplus G,D_T)$ holomorphically isomorphic to $E$.
The spaces of holomorphic maps $F\to E$ and $E\to G$ both have
dimension one: apply $\Hom(F,\cdot)$ and $\Hom(\cdot,G)$ to the
extension and use stability and $\Hom(F,G)=0$. Their counterparts
at $F\oplus G$ are also one-dimensional, generated by the inclusion
and projection.

The Dolbeault Laplacians on these two Hom bundles converge smoothly.
The spectral projections onto their one-dimensional kernels therefore
converge in every smooth norm, by elliptic resolvent estimates.
Choose holomorphic maps $i_T:F\to(F\oplus G,D_T)$ and
$p_T:(F\oplus G,D_T)\to G$ converging respectively to inclusion
and projection. For large $T$, $i_T$ is an embedding and $p_T$ is
surjective. Moreover $p_Ti_T=0$, since $\Hom(F,G)=0$.

Set $j_T=p_T^*(p_Tp_T^*)^{-1}$ and
$\sigma_T(f,g)=i_Tf+j_Tg$. This smooth splitting tends to the
identity and has $p_T\sigma_T(f,g)=g$. Holomorphicity of $i_T,p_T$
implies
\[
 \sigma_T^{-1}D_T\sigma_T
 =\begin{pmatrix}\db_F&\alpha_T\\0&\db_G\end{pmatrix},
 \qquad \alpha_T\to0\quad\hbox{in }C^\infty.
\]
Hodge decomposition writes
$\alpha_T=H\alpha_T+\db v_T$, with $v_T\to0$ smoothly.
The unipotent change of splitting with upper-right entry $-v_T$
makes the extension entry $H\alpha_T$, while the resulting
isomorphism $S_T$ still tends to the identity.

The harmonic form $H\alpha_T$ is a nonzero scalar multiple
$z_T\beta$. Indeed, an isomorphism with $E$ identifies its inclusion
of $F$ and quotient map to $G$ up to nonzero scalars, because the
two Hom spaces above have dimension one. It therefore identifies
the extension classes up to their scalar ratio. Harmonic
representatives preserve this equality. Nonvanishing follows from
nonsplitting, and $z_T\to0$ follows from $\alpha_T\to0$.

Set $\rho_T=|z_T|$. A block-scalar unitary transformation removes
the phase of $z_T$. The pulled-back metric on
$(F\oplus G,\db_0+\rho_TN)$ tends to $h_0$ smoothly, since
$S_T\to\id$ and the phase transformation is unitary and parallel.
Its holomorphic identification with $(E,h(T))$ differs from
$g_{\rho_T}:E\to E_{\rho_T}$ by a scalar automorphism of $E$,
using simplicity from Theorem~\ref{thm:two-factor}. Consequently,
for some $A_T>0$,
\[
 h(T)=a_Tg_{\rho_T}^*\widetilde h_T,\qquad
 \widetilde h_T\longrightarrow h_0\quad\hbox{in }C^\infty.
\]
Thus $h(T)$ is within a multiplicative factor tending to one of
$a_Th_{\rho_T}$. The same is true of $\widehat h_{\rho_T}$
relative to $h_{\rho_T}$ by \eqref{eq:flow-correction}, which
proves \eqref{eq:flow-matching}.
\end{proof}

\begin{proof}[Proof of Theorem~\ref{thm:flow-two-factor}]
For a sufficiently large $T$, start the approximate solution of
Lemma~\ref{lem:flow-approximation} with $\rho_0=\rho_T$.
Multiply it by $A_T$, which changes neither its curvature nor its
endomorphism spectrum. Lemma~\ref{lem:flow-matching} supplies its
initial comparison with $h(T)$. Equations
\eqref{eq:flow-barriers} and \eqref{eq:flow-error-bound} give,
for every $\tau\ge T$,
\[
 e^{-b_T}A_T\widehat h(\tau)
 \le h(\tau)\le e^{b_T}A_T\widehat h(\tau),\qquad
 b_T=\epsilon_T+(A/c)\rho_T\longrightarrow0.
\]
It follows from \eqref{eq:metric-spectral-comparison} and
\eqref{eq:flow-approx-spectrum} that
\[
 \tfrac12e^{-4b_T}
 \le\liminf_{\tau\to\infty}\tau\lambda_1(h(\tau))
 \le\limsup_{\tau\to\infty}\tau\lambda_1(h(\tau))
 \le\tfrac12e^{4b_T}.
\]
Now let $T\to\infty$ to prove \eqref{eq:flow-rate} for the first
eigenvalue. A comparison begun at a single fixed time would give
only two-sided bounds of order $\tau^{-1}$, not the coefficient.

In the unitary gauges of smooth connection convergence, the
endomorphism Laplacians converge smoothly as operators with domain
$H^2$. Their limit is $\Delta_0$, whose kernel on trace-free
endomorphisms is $\CC q$. A contour enclosing only zero therefore
gives spectral projections of rank one for large $\tau$ and a
uniform gap above them. Resolvent convergence on the Sobolev scale
gives smooth convergence of the normalized eigensections and hence
of the finite-rank projections in $L^2\to C^k$ norm for every $k$.
This proves the remaining claims and the recovery formulas.
\end{proof}

The leading coefficient in \eqref{eq:flow-rate} is independent of
the ranks, extension norm, and initial metric under the normalization
\eqref{eq:metric-flow}. These data enter the intermediate constant
$c$, which cancels between the model spectrum $c\rho^2$ and the
flow law $\rho^2\sim(2c\tau)^{-1}$. The theorem does not assert
a next-order eigenvalue expansion or a quantitative rate of
eigensection convergence along the actual flow. In particular, the
fourth-order coefficient in \eqref{eq:two-expansion} cannot simply
be transferred by substituting $\rho=(2c\tau)^{-1/2}$ into the
uncorrected harmonic family. Sections~\ref{sec:diagonal-flow}
to~\ref{sec:complete-spectrum} extend the argument to every simple bundle on
a curve with pairwise nonisomorphic Jordan--H\"older factors, including
eigenvalues with distinct scales.

\section{The diagonal flow}\label{sec:diagonal-flow}

This section is finite-dimensional. It determines the exact asymptotics
of the flow on block-scalar metrics that governs the small spectrum in
Sections~\ref{sec:transfer} and~\ref{sec:complete-spectrum}.

\subsection{Setting and the weight grading}

Let $\Gamma$ be a finite directed acyclic graph with vertex set
$\{1,\ldots,m\}$ and edge set $\Gamma_1$. An edge $\alpha\in\Gamma_1$ from $i$ to
$j$ is written $\alpha=(i\to j)$. Fix masses $M_i>0$ and constants
$c_\alpha>0$, and assume that the underlying undirected graph is connected.
In the geometric application the vertices are the stable factors, an edge
$i\to j$ records a nonzero harmonic entry $\beta_{ij}$ with $F_i$ on the
subbundle side, $M_i=V\rk F_i$, and $c_\alpha=\lVert\beta_{ij}\rVert_{L^2}^2$.

On the hyperplane $\sum_iM_iy_i=0$ of $\mathbb R^m$ put
\begin{equation}\label{eq:diag-flow}
 k_\alpha(y)=c_\alpha e^{2(y_i-y_j)},\qquad
 M_ip_i(y)=\sum_{\alpha\text{ out of }i}k_\alpha
          -\sum_{\alpha\text{ into }i}k_\alpha,\qquad
 \dot y_i=-p_i(y).
\end{equation}
With $d_i=e^{y_i}$ this is the flow $\dot gg^{-1}=-\diag(p_i\id_{F_i})$
for $g=\diag(d_i\id_{F_i})$ studied in Section~\ref{sec:transfer}. Put
$\mathfrak e=\sum_\alpha k_\alpha$, the \emph{total conductance}. The positive eigenvalues
$\nu_1\le\cdots\le\nu_{m-1}$ of the weighted graph form
$\sum_\alpha k_\alpha|x_i-x_j|^2$ on $\{\sum_iM_ix_i=0\}$, with norm
$\sum_iM_i|x_i|^2$, are the \emph{graph eigenvalues} of $y$.

\begin{definition}\label{def:weight-grading}
The \emph{weight grading} is the unique minimizer $w\in\mathbb R^m$ of
$\frac12\sum_iM_iw_i^2$ subject to $w_j-w_i\ge\frac12$ for every edge
$i\to j$. An edge is \emph{tight} if equality holds, and its \emph{slack}
is $\varepsilon_\alpha=2(w_j-w_i)-1\ge0$. The \emph{tight subgraph}
$\mathcal T\subset\Gamma_1$ consists of the tight edges, and its connected components
(on all $m$ vertices, isolated vertices included) are the \emph{tight
components}.
\end{definition}

The constraint set is nonempty, since half the length of a longest path
into each vertex is admissible, and the objective is strictly convex.
The constraints are invariant under adding constants, so
$\sum_iM_iw_i=0$. Since they are linear, there are Lagrange multipliers
$u_\alpha\ge0$, vanishing off $\mathcal T$, with
\begin{equation}\label{eq:kkt}
 M_iw_i=\sum_{\alpha\text{ into }i}u_\alpha
        -\sum_{\alpha\text{ out of }i}u_\alpha .
\end{equation}
Summing over a tight component $\mathcal C$ gives $\sum_{i\in\mathcal C}M_iw_i=0$, because
$u$ is supported on edges inside tight components. A vertex with no tight
edge therefore has $w_i=0$. At least one edge is tight: otherwise every
multiplier would vanish, so $w=0$ by \eqref{eq:kkt}, which is not
admissible.

This is the weight grading of Haiden, Katzarkov, Kontsevich, and Pandit
\cite[Section~3]{HKKP2023}. Their flow \cite[(3.10)]{HKKP2023} is
\eqref{eq:diag-flow} after the substitutions
\[
 x=-2y,\qquad t=2\tau,\qquad v=-2w,\qquad m_i=M_i,
\]
and their tight edges $v_i-v_j=1$ are ours.

\begin{definition}\label{def:nondegenerate}
The pair $(\Gamma,M)$ is \emph{nondegenerate} if some multiplier vector
satisfying \eqref{eq:kkt} is strictly positive on every tight edge.
\end{definition}

In the language of optimization this is strict complementarity. Replacing $M$
by a positive multiple leaves $w$ unchanged and rescales the multipliers,
so for $M_i=V\rk F_i$ nondegeneracy depends only on $\Gamma$ and the ranks.
We then say that $(\Gamma,(r_i))$ is nondegenerate.

This is exactly the hypothesis of \cite[Theorem~3.6]{HKKP2023}, under which
$y=w\log\tau+O(1)$. It fails, for example, for the zigzag
$1\to2\leftarrow3\to4$ with masses satisfying $M_1M_4=M_2M_3$ (compare
\cite[(3.23)--(3.28)]{HKKP2023}),
where iterated logarithms can appear. The equal-rank zigzag of
Theorem~\ref{thm:zigzag-flow} lies on that quadric. In a nondegenerate
graph the tight edges are exactly the edges carrying positive multipliers.

\subsection{Convergence of the remainder}

\begin{theorem}\label{thm:diag-convergence}
Assume $(\Gamma,M)$ is nondegenerate. The tight potential
\[
 \Phi(z)=\sum_iM_iw_iz_i+\frac12\sum_{\alpha\in\mathcal T}c_\alpha e^{2(z_i-z_j)}
\]
has a minimizer $\zeta^*$, unique subject to $\sum_{i\in\mathcal C}M_i\zeta^*_i=0$
for every tight component $\mathcal C$. Every solution of \eqref{eq:diag-flow}
satisfies
\begin{equation}\label{eq:diag-asymptotics}
 y(\tau)=w\log\tau+\zeta^*+\sum_{\mathcal C}Z_{\mathcal C}\mathbf 1_{\mathcal C}+O(\tau^{-\delta})
\end{equation}
for some $\delta>0$ and constants $Z_{\mathcal C}$ with $\sum_{\mathcal C}M_{\mathcal C}Z_{\mathcal C}=0$, where
$M_{\mathcal C}=\sum_{i\in\mathcal C}M_i$ and $\mathbf 1_{\mathcal C}$ is the indicator vector of $\mathcal C$.
The vector $\zeta^*$ depends only on $(\Gamma,M,c)$. The constants $Z_{\mathcal C}$ depend
on the solution.
\end{theorem}

\begin{proof}
Set $s=\log\tau$ for $\tau\ge1$ and $z(s)=y(e^s)-ws$. Then
$\sum_iM_iz_i=0$, and $\tau k_\alpha=c_\alpha e^{2t_\alpha}e^{-\varepsilon_\alpha s}$
with $t_\alpha=z_i-z_j$. Define
\[
 \Psi_s(z)=\sum_iM_iw_iz_i
 +\frac12\sum_\alpha c_\alpha e^{2t_\alpha}e^{-\varepsilon_\alpha s}.
\]
Since $\partial\Psi_s/\partial z_i=M_iw_i+M_i\tau p_i$ and
$dz_i/ds=-\tau p_i-w_i$, the curve $z$ is the gradient flow
$dz/ds=-M^{-1}\nabla\Psi_s(z)$. Because $\partial_s\Psi_s\le0$,
\begin{equation}\label{eq:psi-monotone}
 \frac{d}{ds}\Psi_s(z(s))
 =-\sum_i\frac{(\partial_i\Psi_s)^2}{M_i}
  -\frac12\sum_\alpha\varepsilon_\alpha\tau k_\alpha\le0.
\end{equation}

Fix multipliers $u$ satisfying \eqref{eq:kkt} with $u_\alpha>0$ on $\mathcal T$.
By \eqref{eq:kkt}, $\sum_iM_iw_iz_i=-\sum_\alpha u_\alpha t_\alpha$, so
\[
 \Psi_s(z)=\sum_{\alpha\in\mathcal T}\Bigl(\frac{c_\alpha}2e^{2t_\alpha}
 -u_\alpha t_\alpha\Bigr)
 +\frac12\sum_{\alpha\notin\mathcal T}c_\alpha e^{2t_\alpha}e^{-\varepsilon_\alpha s}.
\]
Each tight term is bounded below by $\frac{u_\alpha}2(1-\log(u_\alpha/c_\alpha))$
and is proper in $t_\alpha$; the other terms are nonnegative. Hence
$\Psi_s(z(s))$ decreases to a finite limit, every tight difference
$t_\alpha$ is bounded along the solution, and \eqref{eq:psi-monotone} gives
\begin{equation}\label{eq:slack-integrable}
 \int_0^\infty\sum_{\alpha\notin\mathcal T}\varepsilon_\alpha\tau k_\alpha\,ds<\infty.
\end{equation}
Within a tight component, $z_i-z_j$ is therefore bounded.

Let $Z_{\mathcal C}(s)=M_{\mathcal C}^{-1}\sum_{i\in\mathcal C}M_iz_i$. In
$\sum_{i\in\mathcal C}\partial_i\Psi_s$ every edge inside $\mathcal C$ cancels, and
$\sum_{i\in\mathcal C}M_iw_i=0$. Edges between different tight components are not
tight, so their slack is at least
$\varepsilon_*=\min_{\alpha\notin\mathcal T}\varepsilon_\alpha>0$. Hence
\[
 \Bigl|\frac{dZ_{\mathcal C}}{ds}\Bigr|
 \le\frac1{M_{\mathcal C}\varepsilon_*}\sum_{\alpha\notin\mathcal T}\varepsilon_\alpha\tau k_\alpha,
\]
which is integrable by \eqref{eq:slack-integrable}. So $Z_{\mathcal C}(s)$ converges,
and $z$ is bounded. Consequently $\tau k_\alpha\le Ce^{-\varepsilon_*s}$ for
$\alpha\notin\mathcal T$. If every edge is tight there is one component, and these
terms are absent.

Let $H$ be the $M$-orthogonal complement of the span of the $\mathbf 1_{\mathcal C}$,
and write $z=\zeta+\sum_{\mathcal C}Z_{\mathcal C}\mathbf 1_{\mathcal C}$ with $\zeta\in H$. The tight
potential satisfies $\Phi(z+\sum_{\mathcal C}\theta_{\mathcal C}\mathbf 1_{\mathcal C})=\Phi(z)$, so
$M^{-1}\nabla\Phi$ takes values in $H$, and
\[
 \frac{d\zeta}{ds}=-M^{-1}\nabla\Phi(\zeta)+R(s),\qquad
 |R(s)|\le Ce^{-\varepsilon_*s}.
\]
On $H$ the map $\zeta\mapsto(t_\alpha)_{\alpha\in\mathcal T}$ is injective, so
$\Phi|_H$ is strictly convex and proper, with a unique minimizer $\zeta^*$,
and its Hessian $2\sum_{\alpha\in\mathcal T}c_\alpha e^{2t_\alpha}(dt_\alpha)^2$ is
uniformly positive on bounded subsets of $H$.

Let $B\subset H$ be a closed ball containing $\zeta^*$ and the bounded curve
$\zeta([0,\infty))$. On $B$, in the norm $|x|_M^2=\sum_iM_ix_i^2$, the
Hessian of $\Phi|_H$ is bounded below by some $c_0>0$. Put
$W=\Phi(\zeta)-\Phi(\zeta^*)\ge0$. For $\zeta\in B$, strong convexity gives
\[
 \tfrac{c_0}2|\zeta-\zeta^*|_M^2\le W\le\tfrac1{2c_0}|\nabla\Phi(\zeta)|_{M^{-1}}^2,
\]
where $|v|_{M^{-1}}^2=\sum_iv_i^2/M_i$. Along the solution,
\[
 \frac{dW}{ds}=-|\nabla\Phi|_{M^{-1}}^2+\langle\nabla\Phi,R\rangle
 \le-\tfrac12|\nabla\Phi|_{M^{-1}}^2+\tfrac12|R|_M^2
 \le-c_0W+C^2e^{-2\varepsilon_*s}.
\]
By Gr\"onwall's inequality, $W(s)\le C'(1+s)e^{-\min(c_0,2\varepsilon_*)s}$,
and hence $|\zeta(s)-\zeta^*|_M=O(e^{-\delta s})$ for every
$\delta<\frac12\min(c_0,2\varepsilon_*)$. Since $z$ is bounded,
$|dZ_{\mathcal C}/ds|\le Ce^{-\varepsilon_*s}$, and so
$|Z_{\mathcal C}(s)-\lim Z_{\mathcal C}|\le Ce^{-\varepsilon_*s}/\varepsilon_*$.
In the variable $\tau=e^s$ this is \eqref{eq:diag-asymptotics}.
\end{proof}

The critical equation of $\Phi$ on $H$ states that
\begin{equation}\label{eq:ustar}
 u^*_\alpha=c_\alpha e^{2(\zeta^*_i-\zeta^*_j)}\qquad(\alpha\in\mathcal T)
\end{equation}
satisfies \eqref{eq:kkt}. It is the unique multiplier vector of this form.
When $\mathcal T$ is a forest, \eqref{eq:kkt} determines $u^*$, which is then
independent of $c$.

\begin{remark}\label{rem:implicit-bias}
After the substitution $\theta=M^{1/2}y$, the flow \eqref{eq:diag-flow} is
gradient flow for the exponential loss $\frac12\sum_\alpha c_\alpha
e^{-\langle x_\alpha,\theta\rangle}$ of the linearly separable vectors
$x_\alpha=2M^{-1/2}(e_j-e_i)$. The weight grading is the corresponding
hard-margin direction, and tight edges are support vectors. For gradient
descent with unit weights, Soudry, Hoffer, Nacson, Gunasekar, and Srebro
prove convergence of the remainder when the support vectors span the data
\cite[Theorems~3 and~4]{Soudry2018}. Here that means there is a single tight
component. Nacson et al.\ \cite[Theorem~10]{Nacson2019} bound the component
of the remainder outside the span of the support vectors without that
assumption. The iterated-logarithm expansion of Soudry et al.\ in the degenerate case
\cite[Theorem~13]{Soudry2018} parallels \cite[Section~5]{HKKP2023}. The
case of several tight components, where the constants $Z_{\mathcal C}$ carry the
dependence on the solution, is the case needed below. The monotone quantity
$\Psi_s$ handles it directly.
\end{remark}

\subsection{Conductances and graph eigenvalues}

Put $a_\alpha=1+\varepsilon_\alpha$ and
$u^0_\alpha=c_\alpha e^{2(\zeta^*_i-\zeta^*_j)}$ for every edge; on tight
edges $u^0_\alpha=u^*_\alpha$. By
Theorem~\ref{thm:diag-convergence},
\begin{equation}\label{eq:edge-asymptotics}
 \tau^{a_\alpha}k_\alpha(y(\tau))\longrightarrow
 k^\infty_\alpha=u^0_\alpha e^{2(Z_{\mathcal C(i)}-Z_{\mathcal C(j)})},
\end{equation}
where $\mathcal C(i)$ is the tight component of $i$. In particular
$\tau k_\alpha\to u^*_\alpha$ on tight edges, whatever the solution.

For an exponent $a$, contract the connected components of the subgraph of
edges with $a_\alpha<a$, adding masses. Join the resulting vertices by the
edges with $a_\alpha=a$ that connect distinct contracted vertices, with
conductances $k^\infty_\alpha$ summed over parallel edges. Let $d_a$ be the number
of positive eigenvalues of this \emph{$a$-quotient graph}. Edges with
$a_\alpha=a$ inside one contracted vertex are omitted; such edges can
occur.

\begin{theorem}\label{thm:diag-spectrum}
Assume $(\Gamma,M)$ is nondegenerate, and let $y$ solve \eqref{eq:diag-flow}.
\begin{enumerate}
\item For each exponent $a$, exactly $d_a$ graph eigenvalues satisfy
 $\nu=\rho\,\tau^{-a}(1+o(1))$, where $\rho$ runs through the positive
 eigenvalues of the $a$-quotient graph, with multiplicity. The
 corresponding spectral subspaces converge to the quotient eigenspaces,
 embedded as vectors constant on the contracted vertices.
\item At $a=1$ the quotient graph is the tight subgraph, which is
 nonempty, with conductances $u^*$ and masses $M_i$. These coefficients are the same for every
 solution. The value $\frac12$ occurs with multiplicity at least the number
 of tight components with at least two vertices. For each such component $\mathcal C$,
 $w\mathbf 1_{\mathcal C}$ is a $\frac12$-eigenvector.
\item For each $a>1$ with $d_a\ge1$, the sum of the coefficients at exponent $a$
 is
 \[
  \sum_{\alpha}k^\infty_\alpha\bigl(M_A^{-1}+M_B^{-1}\bigr),
 \]
 summed over edges with $a_\alpha=a$ that join distinct contracted vertices
 $A\ne B$. As the initial condition varies, this sum is unbounded, and so is
 the largest coefficient at exponent $a$.
\end{enumerate}
\end{theorem}

\begin{proof}
By \eqref{eq:edge-asymptotics}, each conductance lies within a factor
$1+O(\tau^{-\delta})$ of $k^\infty_\alpha\tau^{-a_\alpha}$. A common relative
perturbation of all conductances moves every graph eigenvalue by at most the same
factor, by min--max. The proof of Theorem~\ref{thm:orders} uses only the
ordering of the exponents. With $t=\tau^{-1}$ and real exponents
$a_\alpha$ in place of $2p_{ij}$, it gives the eigenvalue and eigenspace
statements in (1); the $1+o(1)$ factors change each quotient by $o(1)$.

In (2) nothing is contracted at $a=1$, and \eqref{eq:ustar} identifies the
weights. For every tight edge, $w_j-w_i=\frac12$. By \eqref{eq:kkt},
\[
 \sum_{\alpha\in\mathcal T,\;\alpha\ni i}u^*_\alpha(w_i-w_{\alpha(i)})
 =\frac12\Bigl(\sum_{\alpha\text{ into }i}u^*_\alpha
 -\sum_{\alpha\text{ out of }i}u^*_\alpha\Bigr)=\frac12M_iw_i,
\]
where $\alpha(i)$ is the other endpoint. This identity holds row by row, and
the tight Laplacian is block diagonal over tight components. So
$w\mathbf 1_{\mathcal C}$ is a $\frac12$-eigenvector whenever it is nonzero. It is
nonzero when $\mathcal C$ has an edge, and it is $M$-orthogonal to $\mathbf 1_{\mathcal C}$.

For (3), the trace of the Laplacian of the $a$-quotient graph with respect to the
masses is the displayed sum. Its terms are
$u^0_\alpha e^{2(Z_{\mathcal C(i)}-Z_{\mathcal C(j)})}$ with $\mathcal C(i)\ne \mathcal C(j)$, and the linear
forms $Z_{\mathcal C(i)}-Z_{\mathcal C(j)}$ are nonzero on $\{\sum M_{\mathcal C}Z_{\mathcal C}=0\}$. So the sum is
unbounded as $Z$ ranges over this hyperplane.

Every such $Z$ is approximately attained. Start at a late time $s_0$ from
$\zeta^*+\sum_{\mathcal C}\theta_{\mathcal C}\mathbf 1_{\mathcal C}$, with
$\sum_{\mathcal C}M_{\mathcal C}\theta_{\mathcal C}=0$, and put $\Theta=\max|\theta_{\mathcal C}|$.
Let $s_2\in(s_0,\infty]$ be the supremum of the $s$ for which
$|z(\sigma)-z(s_0)|_\infty\le1$ on $[s_0,s]$. On $[s_0,s_2)$, every
difference $t_\alpha$ is at most $2\Theta+2+\max|\zeta^*_i-\zeta^*_j|$. So every
non-tight conductance satisfies $\tau k_\alpha\le C_1e^{4\Theta}e^{-\varepsilon_*s}$,
where $C_1$ does not depend on $(\theta_{\mathcal C})$ or $s_0$. Let $B$ now
be the closed ball in $H$ about $\zeta^*$ containing all $\zeta$ with
$|z-z(s_0)|_\infty\le1$, for $z(s_0)$ of the above form. Its convexity
constant $c_0$ depends only on $\Gamma$, $M$, $c$ and $\Theta$. There:
\begin{itemize}
\item $|dZ_{\mathcal C}/ds|\le C_2e^{4\Theta}e^{-\varepsilon_*s}$, so
 $|Z_{\mathcal C}(s)-\theta_{\mathcal C}|\le C_2e^{4\Theta-\varepsilon_*s_0}/\varepsilon_*$;
\item the transverse part starts at the critical point, $\zeta(s_0)=\zeta^*$.
 The estimate above, with $R=O(e^{4\Theta-\varepsilon_*s})$, gives
 $|\zeta(s)-\zeta^*|_M\le C_3e^{4\Theta-\varepsilon_*s_0}$.
\end{itemize}
The deviation $|z(s)-z(s_0)|_\infty$ is at most
$C_2e^{4\Theta-\varepsilon_*s_0}/\varepsilon_*
+(\min_iM_i)^{-1/2}C_3e^{4\Theta-\varepsilon_*s_0}$. If $s_0$ is so large
that this is below $\frac12$, the deviation stays below $\frac12$ on
$[s_0,s_2)$. By continuity this
forces $s_2=\infty$. Hence
\[
 \lim_{s\to\infty}Z_{\mathcal C}(s)=\theta_{\mathcal C}+O(e^{4\Theta-\varepsilon_*s_0}).
\]
The flow is autonomous, and a shift of time does not change the leading
coefficient of $\tau^{-a}$. So every vector $(\theta_{\mathcal C})$ is, up to an
arbitrarily small error, the vector of limits $Z_{\mathcal C}$ of some
solution.
\end{proof}

\section{Uniform spectral transfer}\label{sec:transfer}

Throughout this section $X$ is a compact curve. Let $E$ be a simple
semistable holomorphic bundle whose Jordan--H\"older factors
$F_1,\ldots,F_m$, $m\ge2$, are pairwise nonisomorphic, with ranks $r_i$
and Hermitian--Einstein metrics $h_i$. Put $E_0=\bigoplus_iF_i$,
$h_0=\bigoplus_ih_i$, $M_i=Vr_i$, and $|M|=\sum_iM_i$.

\subsection{The socle filtration and the harmonic normal form}

The \emph{socle} $\operatorname{soc}(E)$ is the sum of the stable
subbundles of $E$ of slope $\mu(E)$; it is polystable. Put
$S_1=\operatorname{soc}(E)$ and let $S_{k+1}\supset S_k$ be the preimage
of $\operatorname{soc}(E/S_k)$. This \emph{socle filtration}
$0=S_0\subset S_1\subset\cdots\subset S_L=E$ is canonical, and each
$S_k/S_{k-1}$ is a direct sum of some of the $F_i$. Since the factors are
pairwise nonisomorphic, each $F_i$ occurs in exactly one
$S_{k}/S_{k-1}$; we call $k$ the \emph{socle layer} of $F_i$. Two-step
extensions are the case $L=2$.

\begin{lemma}\label{lem:socle-normal-form}
There is a smooth identification $E\cong(E_0,\db_0+N)$, preserving the
socle filtration, in which $N=\sum\beta_{ij}$ with
$\beta_{ij}\in\Omega^{0,1}(\Hom(F_j,F_i))$ harmonic, and $\beta_{ij}\ne0$
only if the layer of $F_i$ is smaller than that of $F_j$. The directed
graph $\Gamma$ with an edge $i\to j$ for each nonzero $\beta_{ij}$ is
connected.
\end{lemma}

\begin{proof}
Refine the socle filtration to a Jordan--H\"older filtration and apply
Lemma~\ref{lem:normal-form}. Each layer occupies a consecutive range of
Jordan--H\"older indices, and the corresponding diagonal block of $N$ is
the normal form of the polystable, hence split, subquotient
$S_k/S_{k-1}$. Its entries between adjacent indices are harmonic
representatives of zero classes, so they vanish. Suppose the entries of
index gap less than $\sigma$ inside a layer vanish. Then the subquotient
on indices $a,\ldots,a+\sigma$ is the direct sum of the intermediate
factors and the extension of $F_{a+\sigma}$ by $F_a$ given by the harmonic
entry. Being a subquotient of a polystable bundle, it is split, so that
entry vanishes as well. Connectedness is equivalent to simplicity by the
argument of Theorem~\ref{thm:relative}.
\end{proof}

A factor in layer $k\ge2$ has a nonzero entry from layer $k-1$.
Otherwise it would be a holomorphic subbundle of $E/S_{k-2}$ and would lie
in layer $k-1$.

\subsection{Diagonal metrics and their correction}

For $g=\diag(d_i\id_{F_i})$ with $d_i>0$ and $\prod_id_i^{r_i}=1$, put
$N_g=gNg^{-1}$, $h_g=g^*h_0$,
\[
 k_{ij}(g)=\lVert(N_g)_{ij}\rVert_{L^2}^2,\qquad
 M_ip_i(g)=\sum_{j}k_{ij}(g)-\sum_{l}k_{li}(g),\qquad
 \mathfrak e_g=\sum_{ij}k_{ij}(g),
\]
and let $\nu_1(g)\le\cdots\le\nu_{m-1}(g)$ be the positive graph
eigenvalues of Section~\ref{sec:diagonal-flow} for these conductances. We
call $\mathfrak e_g=\lVert N_g\rVert_{L^2}^2$ the \emph{total conductance},
and write $D_g=\db_{0,\End}+[N_g,\cdot\,]$ for the Dolbeault operator of
$(E_0,\db_0+N_g)$ on endomorphisms. The
\emph{diagonal flow} is \eqref{eq:diag-flow}, that is
$\dot gg^{-1}=-\diag(p_i(g)\id_{F_i})$.

The mean curvature of $(\db_0+N_g,h_0)$ is $\kappa\id+C_g$, where
$C_g=\sqrt{-1}\Lambda_\omega(N_g-N_g^*)\wedge(N_g-N_g^*)$, as in
\eqref{eq:several-curvature}. Taking block traces shows
$PC_g=\diag(p_i(g)\id_{F_i})$. Theorem~\ref{thm:relative} holds with
independent edge amplitudes, since its proof uses only harmonicity and
$\eta=\lVert[N_g,\cdot]\rVert\le C\sqrt{\mathfrak e_g}$:
\begin{equation}\label{eq:general-independent-comparison}
 \lambda_j(h_g)=\nu_j(g)(1+O(\mathfrak e_g)),\qquad 1\le j\le m-1,
\end{equation}
uniformly as $\mathfrak e_g\to0$, with no condition on the ratios of the conductances.

\begin{lemma}\label{lem:layer-energy}
Along the diagonal flow, $\dot{\mathfrak e}=-2\sum_iM_ip_i^2$ and
\begin{equation}\label{eq:layer-energy}
 \dot{\mathfrak e}\le-\frac{8\mathfrak e^2}{(L-1)^2|M|},\qquad
 \int_T^\infty\mathfrak e^{3/2}\,d\tau\le\frac{(L-1)^2|M|}{4}\sqrt{\mathfrak e(T)}.
\end{equation}
The flow exists for all time and preserves $\prod_id_i^{r_i}=1$.
\end{lemma}

\begin{proof}
Each conductance satisfies $\dot k_{ij}=-2(p_i-p_j)k_{ij}$, and summing gives
the formula for $\dot{\mathfrak e}$. Let $\varphi_i$ be minus the layer of $F_i$,
shifted so that $\sum_iM_i\varphi_i=0$. Every edge $i\to j$ has
$\varphi_i-\varphi_j\ge1$, so
\[
 \mathfrak e\le\sum_{i\to j}k_{ij}(\varphi_i-\varphi_j)=\sum_iM_ip_i\varphi_i
 \le\Bigl(\sum_iM_ip_i^2\Bigr)^{1/2}\Bigl(\sum_iM_i\varphi_i^2\Bigr)^{1/2}.
\]
The shifted values lie in an interval of length $L-1$, so
$\sum_iM_i\varphi_i^2\le(L-1)^2|M|/4$. This proves the differential
inequality, and integration gives the integral bound. The total conductance
decreases,
so the $p_i$ are bounded on finite intervals, and the positive $d_i$ cannot
reach zero or infinity in finite time. Since $\sum_ir_ip_i=0$, the weighted
determinant is preserved.
\end{proof}

\begin{lemma}\label{lem:general-correction}
For sufficiently small $\mathfrak e_g$, set $v_g=RQC_g$, $k_g=\exp(-v_g/2)$, and
let $\widehat h_g=(k_gg)^*h_0$ be the \emph{corrected diagonal metric}. Along the diagonal flow the residual
\eqref{eq:flow-residual} satisfies
$\lVert\mathcal E_{\widehat h_g}\rVert_{L^\infty,\op}\le A\mathfrak e_g^{3/2}$.
Moreover $e^{-C\mathfrak e_g}h_g\le\widehat h_g\le e^{C\mathfrak e_g}h_g$ and
$\lambda_j(\widehat h_g)=\nu_j(g)(1+O(\mathfrak e_g))$. The constants depend only on
the fixed forms and $h_0$.
\end{lemma}

\begin{proof}
The split Laplacian preserves adjoints, so $v_g$ is self-adjoint, and
$v_g=O_{C^k}(\mathfrak e_g)$ for every $k$. The first variation of mean curvature at
the split metric is $s\mapsto\Delta_0s$. Taylor expansion therefore gives,
in the gauge $k_gg$,
\[
 \widehat K_g=C_g-\Delta_0v_g+O_{C^k}(\mathfrak e_g^{3/2})
 =PC_g+O_{C^k}(\mathfrak e_g^{3/2}).
\]
Uniformity follows from the fixed finite collection of harmonic forms:
$N_g=O_{C^k}(\sqrt{\mathfrak e_g})$, and the mixed terms between $N_g$ and $v_g$ are
$O_{C^k}(\mathfrak e_g^{3/2})$.

The diagonal flow gives $\dot N_g=-[PC_g,N_g]=O_{C^k}(\mathfrak e_g^{3/2})$, hence
$\dot v_g=O_{C^k}(\mathfrak e_g^2)$. Although $v_g$ need not commute with $PC_g$,
\[
 \partial_\tau(k_gg)(k_gg)^{-1}=\dot k_gk_g^{-1}-k_g\,PC_g\,k_g^{-1}
 =-PC_g+O_{C^k}(\mathfrak e_g^2).
\]
Its sum with its adjoint is the transported metric derivative, so the terms
$\mp2PC_g$ cancel in the residual \eqref{eq:flow-residual}. This proves the
residual bound. The metric comparison follows from $v_g=O(\mathfrak e_g)$, and the
spectral statement from \eqref{eq:metric-spectral-comparison} and
\eqref{eq:general-independent-comparison}.
\end{proof}

\subsection{Matching the flow at late times}

\begin{lemma}\label{lem:socle-matching}
Let $h$ solve \eqref{eq:metric-flow} on $E$. For every large $T$,
\[
 h(T)=A_Tg_T^*\widetilde h_T,\qquad A_T>0,\qquad
 \widetilde h_T\longrightarrow h_0\ \text{in }C^\infty,\qquad
 \mathfrak e_{g_T}\longrightarrow0,
\]
where $g_T$ is diagonal with $\prod_id_i^{r_i}=1$, and $\widetilde h_T$ is
a metric on $(E_0,\db_0+N_{g_T})$. We call the maps $g_T$ the
\emph{diagonal gauges} of $h$.
\end{lemma}

\begin{proof}
In the smooth unitary flow gauges write $D_T\to\db_0$ on $E_0$. Each
$(E_0,D_T)$ is isomorphic to $E$.

\emph{Hom dimensions.} Let $F_i$ lie in layer $k$. A nonzero map from $F_i$
to a semistable bundle of slope $\mu$ is injective with saturated image,
and that image lies in the socle. Hence $\dim\Hom(F_i,E/S_{k-1})=1$. At
the split limit every such Hom space also has dimension one.

\emph{Splitting.} For $F_i$ in layer one, the Hom-bundle Laplacians
converge smoothly, so their one-dimensional kernel projections do too. This
gives holomorphic maps $\iota_{i,T}:F_i\to(E_0,D_T)$ converging to the
standard inclusions. For large $T$ their sum is fiberwise injective, and
its image is the socle of $(E_0,D_T)$. The $h_0$-orthogonal complement
carries a quotient holomorphic structure on the remaining factors that
converges smoothly to $\db_0$ and is isomorphic to $E/S_1$. Repeating this
$L$ times gives a smooth identification $\sigma_T\to\id$ transporting $D_T$
to $\db_0+\alpha_T$. Here $\alpha_T$ is strictly block upper-triangular in
layers, $\alpha_T\to0$ smoothly, and the diagonal blocks are exactly
$\db_{F_i}$.

\emph{Harmonic form.} Apply the triangular Hodge induction in the proof of
Lemma~\ref{lem:normal-form}, with the filtration by layer gap in place of
the Jordan--H\"older index gap. This is legitimate because $\alpha_T$ has no
entries inside a layer, and products of entries of layer gaps $a$ and $b$
have layer gap $a+b$. At each step the entries of the next layer gap are
replaced by their harmonic parts, using the gauge $\db_0^*\mathbb G$ of the
current entries, where $\mathbb G$ is the Green operator of the Dolbeault
Laplacian. The new product terms are polynomials in entries that tend to
zero in $C^\infty$. Hence the gauge tends to the identity, and the resulting
harmonic normal form $\db_0+N'_T$ has $N'_T\to0$.

\emph{Uniqueness up to diagonal maps.} Let
$\Phi:(E_0,\db_0+N)\to(E_0,\db_0+N'_T)$ be a holomorphic isomorphism. It
preserves the socle filtration, which is the standard layer flag on both
sides. On $N'_T$ this is because the splitting sends the flag to the socle
and the Hodge gauge is unipotent in layers. Inside a layer the gauge
equation reduces to $\db_0\Phi_{ij}=0$, so $\Phi$ is $\diag(z_i)$ there.
Write $\Phi=Z+U$ with $U$ strictly upper-triangular in layers. The
$(i,j)$ entry of $\db_0\Phi+N'_T\Phi-\Phi N=0$ is
\[
 \db_0U_{ij}+z_jN'_{T,ij}-z_iN_{ij}
 +\sum_{k}\bigl(N'_{T,ik}U_{kj}-U_{ik}N_{kj}\bigr)=0,
\]
summed over $k$ in layers strictly between those of $i$ and $j$. By
induction on the layer gap, the sum vanishes. The exact form
$\db_0U_{ij}$ then equals a harmonic form, so both vanish, and $U_{ij}$ is
a holomorphic section of $\Hom(F_j,F_i)$, hence zero. Therefore $\Phi=Z$
and $N'_T=ZNZ^{-1}$.

A parallel diagonal unitary map removes the phases of the $z_i$. The
moduli $d_i=|z_i|$, normalized so that $\prod_id_i^{r_i}=1$, define $g_T$;
this does not change $N_{g_T}$. The metric $\widetilde h_T$ tends to $h_0$,
because the splitting and the Hodge gauge tend to the identity and the
unitary maps are parallel. Also $\mathfrak e_{g_T}\to0$, because $N'_T\to0$. The
resulting identification with $E$ differs from $g_T$ by an automorphism of
$E$, which is scalar since $E$ is simple. Its squared modulus is $A_T$.
\end{proof}

\subsection{The transfer theorem}

\begin{theorem}\label{thm:transfer}
Start the flow \eqref{eq:metric-flow} on $E$ from any smooth metric. For
every sufficiently large $T$ there are a solution $g^{[T]}$ of the diagonal
flow on $[T,\infty)$ and a number $\delta_T\to0$ such that
\begin{equation}\label{eq:transfer}
 e^{-\delta_T}\nu_j(g^{[T]}(\tau))\le\lambda_j(h(\tau))
 \le e^{\delta_T}\nu_j(g^{[T]}(\tau)),\qquad\tau\ge T,\quad1\le j<m.
\end{equation}
Moreover $\liminf\lambda_m(h(\tau))>0$, and in smooth unitary flow gauges
the full small spectral projection converges in every $L^2\to C^k$ norm to
the projection onto $\mathcal K$.
\end{theorem}

\begin{proof}
At time $T$ apply Lemma~\ref{lem:socle-matching}. Start the diagonal flow
at $g_T$, and compare $h$ with $A_T\widehat h_{g^{[T]}(\tau)}$. The initial
discrepancy is $O(\epsilon_T+\mathfrak e_{g_T})$, where
$e^{-\epsilon_T}h_0\le\widetilde h_T\le e^{\epsilon_T}h_0$.
Lemmas~\ref{lem:layer-energy} and~\ref{lem:general-correction} bound the
total residual by $O(\sqrt{\mathfrak e_{g_T}})$. The barriers \eqref{eq:flow-barriers}
then give
\begin{equation}\label{eq:general-restarted-barrier}
 e^{-b_T}A_T\widehat h_{g^{[T]}(\tau)}\le h(\tau)\le
 e^{b_T}A_T\widehat h_{g^{[T]}(\tau)},\qquad
 b_T=\epsilon_T+C\mathfrak e_{g_T}+\frac{A(L-1)^2|M|}{4}\sqrt{\mathfrak e_{g_T}}.
\end{equation}
Combine this with \eqref{eq:metric-spectral-comparison} and
Lemma~\ref{lem:general-correction} to obtain \eqref{eq:transfer} with
$\delta_T\le4b_T+C\mathfrak e_{g_T}$. The remaining assertions follow by
resolvent convergence, as in Theorem~\ref{thm:flow-two-factor}.
\end{proof}

\begin{corollary}\label{cor:coefficient-transfer}
Fix $j<m$ and a positive function $f$ with $f(\tau+\sigma)/f(\tau)\to1$ for
every fixed $\sigma$. Suppose every solution $g$ of the diagonal flow satisfies
$\nu_j(g(\tau))/f(\tau)\to c_j(g)>0$. Then every solution of
\eqref{eq:metric-flow} has $\lambda_j(h(\tau))/f(\tau)\to c_j(h)>0$, with
$|\log(c_j(h)/c_j(g^{[T]}))|\le\delta_T$.
\end{corollary}

\begin{proof}
At one fixed large $T$, the comparison \eqref{eq:transfer} gives positive
finite bounds on the lower and upper limits $\ell,u$ of $\lambda_j/f$. For
each larger $T$, the comparison on $[T,\infty)$ gives
$e^{-\delta_T}c_j(g^{[T]})\le\ell\le u\le e^{\delta_T}c_j(g^{[T]})$. Hence
$u/\ell\le e^{2\delta_T}$ for all large $T$, and $u=\ell$. The shift
condition makes the choice of time origin of $g^{[T]}$ immaterial.
\end{proof}

\begin{lemma}\label{lem:instantaneous-edges}
Let $g_\tau$ be the diagonal gauges of Lemma~\ref{lem:socle-matching} at
time $\tau$. For every edge and every large $T$,
\[
 e^{-2b_T-o_\tau(1)}\le\frac{k_{ij}(g_\tau)}{k_{ij}(g^{[T]}(\tau))}
 \le e^{2b_T+o_\tau(1)}.
\]
\end{lemma}

\begin{proof}
Evaluate \eqref{eq:general-restarted-barrier} and the representation
$h(\tau)=A_\tau g_\tau^*\widetilde h_\tau$ on a nonzero vector of $F_i$.
Since $g$ acts on $F_i$ by the scalar $d_i$, and
$\widehat h_g\in e^{\pm C\mathfrak e_g}h_g$, this gives $d_i(g_\tau)^2/d_i(g^{[T]}(\tau))^2$
up to the factor $A_T/A_\tau$ and $e^{\pm(b_T+o_\tau(1))}$. Dividing the
bounds for the two endpoints of an edge cancels $A_T/A_\tau$.
\end{proof}

Two choices of the diagonal gauges $g_\tau$, with
$\widetilde h,\widetilde h'\to h_0$, have conductance ratios tending to one,
by the same argument. The limits of conductances along $g_\tau$ in
Section~\ref{sec:complete-spectrum} therefore do not depend on these
choices.

\section{The complete leading small spectrum}\label{sec:complete-spectrum}

We keep the setting of Section~\ref{sec:transfer}: a compact curve, and a
simple bundle $E$ with pairwise nonisomorphic Jordan--H\"older factors,
written in the normal form of Lemma~\ref{lem:socle-normal-form} with graph
$\Gamma$ and masses $M_i=Vr_i$. We use the weight grading $w$, the tight subgraph
$\mathcal T$, the slacks $\varepsilon_\alpha$, the exponents
$a_\alpha=1+\varepsilon_\alpha$, and the $a$-quotient graphs of
Section~\ref{sec:diagonal-flow}. These depend only on $\Gamma$ and the ranks.
List the small eigenvalues as $\lambda_1\le\cdots\le\lambda_{m-1}$. Assign
exponents to indices so that larger exponents take smaller indices: the
first $d_{a_{\max}}$ indices receive $a_{\max}$, and so on. Write $a_j$ for
the exponent of index $j$.

\begin{theorem}\label{thm:complete-spectrum}
Assume $(\Gamma,M)$ is nondegenerate. For every smooth initial metric on $E$:
\begin{enumerate}
\item $\tau^{a_j}\lambda_j(h(\tau))\to c_j(h)>0$ for every $j<m$.
\item The coefficients with $a_j=1$ are the positive eigenvalues of the
 tight subgraph with conductances $u^*$ and masses $M_i$, with multiplicity.
 They do not depend on the initial metric. The value $\frac12$ occurs with
 multiplicity at least the number of tight components with at least two
 vertices.
\item For every edge the limit $k^\infty_\alpha(h)=\lim\tau^{a_\alpha}
 k_\alpha(g_\tau)$ exists along the diagonal gauges of
 Lemma~\ref{lem:socle-matching}. For $a_j=a>1$, the coefficients $c_j(h)$
 are the positive eigenvalues of the $a$-quotient graph with conductances
 $k^\infty(h)$, with multiplicity. For every $a>1$ with $d_a\ge1$, their sum,
 and hence the largest of them, is unbounded as $h(0)$ ranges over metrics
 on the single bundle $E$.
\item In smooth unitary flow gauges, for each exponent $a$ and each
 distinct value $\rho$ among the coefficients with $a_j=a$, the spectral
 projection onto $\{\lambda_j:a_j=a,\ c_j(h)=\rho\}$ converges in every
 $L^2\to C^k$ norm to the $\rho$-eigenspace of the $a$-quotient graph. This
 eigenspace is embedded in $\mathcal K$ as block-scalar endomorphisms
 constant on the contracted vertices. For $\rho=\frac12$ at $a=1$ it
 contains $\sum_{i\in\mathcal C}w_i\id_{F_i}$ for each tight component $\mathcal C$ with at
 least two vertices.
\end{enumerate}
\end{theorem}

\begin{proof}
\emph{Eigenvalues.} By Theorems~\ref{thm:diag-convergence}
and~\ref{thm:diag-spectrum}, every solution $g$ of the diagonal flow has
$\tau^{a_j}\nu_j(g(\tau))\to\rho_j(g)>0$. The exponent assigned to each
index is the same for all solutions, and $f=\tau^{-a}$ satisfies the shift
condition. Corollary~\ref{cor:coefficient-transfer} gives (1), with
$|\log(c_j(h)/\rho_j(g^{[T]}))|\le\delta_T$. At $a=1$ the value
$\rho_j(g)$ is the same for all $g$, which proves (2).

\emph{Edge limits.} By \eqref{eq:edge-asymptotics},
$\tau^{a_\alpha}k_\alpha(g^{[T]}(\tau))\to k^\infty_\alpha(g^{[T]})$.
Lemma~\ref{lem:instantaneous-edges} puts the lower and upper limits of
$\tau^{a_\alpha}k_\alpha(g_\tau)$ within $e^{\pm2b_T}$ of
$k^\infty_\alpha(g^{[T]})$ for every large $T$. Since $b_T\to0$, the limit
$k^\infty_\alpha(h)$ exists and equals $\lim_Tk^\infty_\alpha(g^{[T]})$. The
positive eigenvalues of each $a$-quotient graph depend continuously on the
conductances, so $c_j(h)$ is the corresponding eigenvalue for $k^\infty(h)$.

\emph{Unboundedness.} Fix $a>1$ with $d_a\ge1$ and constants $\theta_{\mathcal C}$ with
$\sum_{\mathcal C}M_{\mathcal C}\theta_{\mathcal C}=0$. For a large $s_0$, let $g_*$ be the diagonal map with
logarithms $w s_0+\zeta^*+\sum_{\mathcal C}\theta_{\mathcal C}\mathbf 1_{\mathcal C}$, and set
$h(0)=g_*^*h_0$. This is a metric on $E$.
\begin{itemize}
\item The diagonal flow from $g_*$ at time zero is a time shift of the
 solution considered at the end of the proof of
 Theorem~\ref{thm:diag-spectrum}. Hence its constants satisfy
 $Z_{\mathcal C}=\theta_{\mathcal C}+O(e^{4\max|\theta_{\mathcal C}|-\varepsilon_*s_0})$.
\item Its total conductance is
 $\mathfrak e_0\le Ce^{-s_0}+Ce^{4\max|\theta_{\mathcal C}|}e^{-(1+\varepsilon_*)s_0}$.
\item Compare $h$ on $[0,\infty)$ with the corrected diagonal metrics along
 the diagonal flow from $g_*$. The initial discrepancy is $O(\mathfrak e_0)$,
 and Lemma~\ref{lem:layer-energy} bounds the total residual by
 $O(\sqrt{\mathfrak e_0})$. The barriers give
 $e^{-b}\widehat h_{g(\tau)}\le h(\tau)\le e^b\widehat h_{g(\tau)}$ for
 all $\tau\ge0$, with $b=O(\mathfrak e_0+\sqrt{\mathfrak e_0})$ and constants independent of
 $g_*$.
\end{itemize}
Therefore $c_j(h)\in e^{\pm4b}\rho_j(g)$ for every $j$. Take
$s_0\gg4\max|\theta_{\mathcal C}|/\varepsilon_*$. Then
Theorem~\ref{thm:diag-spectrum}(3) shows that the sum of the coefficients
at exponent $a$ takes arbitrarily large values as $(\theta_{\mathcal C})$ varies.

\emph{Spectral subspaces.} Let $u$ be a normalized eigensection with
$a_j=a$ and $c_j(h)=\rho$, and
transport it to $(\db_0+N_{g_\tau},\widetilde h_\tau)$. Write $u=x+y$ with
$x\in\mathcal K$ and $y\perp\mathcal K$ in $h_0$.
\begin{itemize}
\item Changing $\widetilde h_\tau$ to $h_0$ alters pointwise norms by a
 factor $1+o(1)$. So $\lVert D_{g_\tau}u\rVert_{h_0}^2\le\lambda(1+o(1))$
 with $\lambda=O(\tau^{-a})$. By \eqref{eq:form-lower}, which uses the
 cancellation only in $h_0$,
 \[
  \lVert[N_{g_\tau},x]\rVert^2+\lVert y\rVert^2\le C\lambda.
 \]
\item For edges with $a_\alpha<a$, $k_\alpha(g_\tau)\ge c\tau^{-a_\alpha}$
 by $k^\infty_\alpha(h)>0$. So every limit $x_\infty$ of $x$ is constant on
 the contracted vertices. The identifications with the diagonal gauges consist of a map
 tending to the identity followed by a parallel diagonal unitary map, which
 fixes $\mathcal K$ pointwise. Full-projection convergence
 (Theorem~\ref{thm:transfer}) therefore gives $y\to0$ in every $C^k$
 norm.
\item Test the weak eigenvalue equation against a fixed block-scalar $z$
 constant on the contracted vertices. Then
 $\lVert[N_{g_\tau},z]\rVert^2=O(\tau^{-a})$. Replacing $\widetilde h_\tau$
 by $h_0$ costs $o(\tau^{-a})$. In $h_0$ the complementary term is
 $\langle[N_{g_\tau},y],[N_{g_\tau},z]\rangle$, since $D_0^*[N,z]=0$. It is
 at most
 \[
  \lVert N_{g_\tau}\rVert_{L^\infty}\lVert y\rVert\lVert[N_{g_\tau},z]\rVert
  =O(\tau^{-1/2})O(\tau^{-a/2})O(\tau^{-a/2})=o(\tau^{-a}),
 \]
 using $\mathfrak e_{g_\tau}=O(\tau^{-1})$, which follows from the decay bound of
 Lemma~\ref{lem:layer-energy} together with
 Lemma~\ref{lem:instantaneous-edges}.
\end{itemize}
Multiply by $\tau^a$ and pass to the limit, using the edge limits. Edges
with $a_\alpha<a$ do not contribute, since $z$ is constant across them, and
edges with $a_\alpha>a$ contribute $o(1)$. So $x_\infty$ is a
$\rho$-eigenvector of the $a$-quotient graph. The number of eigenvalues in
this group equals the multiplicity of $\rho$. Limits of orthonormal
eigensections remain orthonormal because $y\to0$. Compactness in
$\mathcal K$ and smooth convergence of $y$ give the stated $L^2\to C^k$
convergence. The last assertion of (4) is
Theorem~\ref{thm:diag-spectrum}(2).
\end{proof}

\begin{corollary}\label{cor:metric-asymptotics}
Under the hypotheses of Theorem~\ref{thm:complete-spectrum}, write
$h(\tau)=A_\tau g_\tau^*\widetilde h_\tau$ with the diagonal gauges $g_\tau$
of Lemma~\ref{lem:socle-matching} and $\widetilde h_\tau\to h_0$. The limits
$Z_{\mathcal C}(h)=\lim_TZ_{\mathcal C}(g^{[T]})$ exist, and, up to a common
scalar,
\[
 \log g_\tau=w\log\tau+y_\infty(h)+o(1),\qquad
 y_\infty(h)=\zeta^*+\sum_{\mathcal C}Z_{\mathcal C}(h)\mathbf 1_{\mathcal C}.
\]
\end{corollary}

\begin{proof}
The edge limits give convergence of
$\log d_i-\log d_j-(w_i-w_j)\log\tau$ along every edge. Since the graph is
connected, all such differences converge.
\end{proof}

Haiden, Katzarkov, Kontsevich, and Pandit \cite[Theorem~1.1]{HKKP2018} and
Lam \cite[Theorem~1.1]{Lam2026} describe the metric growth up to a bounded
term. Corollary~\ref{cor:metric-asymptotics} shows that in the
nondegenerate case this bounded term converges. The limits
$Z_{\mathcal C}(h)$ exist because the edge limits converge and the quotient
of $\Gamma$ by the tight components is connected. So the limit depends on the
initial metric only through the constants $Z_{\mathcal C}(h)$.

\section{Examples along the flow}\label{sec:flow-examples}

We work on a compact curve $X$ of genus at least two. There, stable bundles
of every rank and degree exist. For stable $F_i,F_j$ of equal slope that are
not isomorphic, $h^0(\Hom(F_j,F_i))=0$ and Riemann--Roch gives
$h^1(\Hom(F_j,F_i))=r_ir_j(g(X)-1)>0$. Every graph $\Gamma$ of the kind
considered below is therefore realized by nonzero harmonic entries.

\subsection{All edges tight}

When every edge is tight and $(\Gamma,(r_i))$ is nondegenerate, there is one tight
component and every small eigenvalue decays like $\tau^{-1}$. For two-step
extensions this is the \emph{positive balance condition}. It asks for
positive numbers $\sigma_{ij}$ on the edges with prescribed margins,
\begin{equation}\label{eq:positive-balance}
 \sum_j\sigma_{ij}=\frac{M_iM_J}{2|M|}\quad(i\in I),\qquad
 \sum_i\sigma_{ij}=\frac{M_jM_I}{2|M|}\quad(j\in J),
\end{equation}
where $M_I=\sum_{i\in I}M_i$ and $M_J=\sum_{j\in J}M_j$. Feasibility is a
classical matrix-scaling question \cite{SinkhornKnopp1967,RothblumSchneider1989}.

Condition \eqref{eq:positive-balance} holds exactly when every edge is tight
and $(\Gamma,(r_i))$ is nondegenerate. Let $\bar w$ take the value $-M_J/(2|M|)$ on
$I$ and $M_I/(2|M|)$ on $J$. It is admissible, with equality on every edge,
and \eqref{eq:positive-balance} is the multiplier equation \eqref{eq:kkt}
for $\bar w$ with multipliers $\sigma$. A pair satisfying these optimality
conditions of a convex program characterizes its minimizer, so $\bar w$ is the
weight grading, every edge is tight, and $(\Gamma,(r_i))$ is nondegenerate.
Conversely, if every edge is tight and $\Gamma$ is connected and bipartite, the
differences $w_j-w_i=\frac12$ and $\sum_iM_iw_i=0$ force $w=\bar w$, and
positive multipliers satisfy \eqref{eq:positive-balance}.

\begin{corollary}\label{cor:balanced}
Assume \eqref{eq:positive-balance}. There are unique conductances
$k^*_{ij}=\lVert\beta_{ij}\rVert^2e^{2(y^*_i-y^*_j)}$ with these margins,
normalized by $\sum_iM_iy^*_i=0$. Then
$\tau\lambda_j(h(\tau))\to\rho_j$ for every initial metric, where
$\rho_1\le\cdots\le\rho_{m-1}$ are the positive eigenvalues of the graph
with conductances $k^*$ and masses $M_i$. The spectral projections onto the
eigenvalues with a common limit converge to the
corresponding graph eigenspaces.
\end{corollary}

\begin{proof}
This is Theorem~\ref{thm:complete-spectrum} with $\mathcal T=\Gamma_1$ and $u^*=k^*$.
\end{proof}

\begin{example}\label{ex:four-cycle-flow}
Take four distinct degree-zero line bundles and all four entries
$\beta_{12},\beta_{14},\beta_{32},\beta_{34}$ nonzero, with $V=1$. The
quantity
\[
 \Theta=\frac{\lVert\beta_{12}\rVert^2\lVert\beta_{34}\rVert^2}
             {\lVert\beta_{14}\rVert^2\lVert\beta_{32}\rVert^2}
\]
is invariant under diagonal changes of the factor metrics. The balanced
matrix, with rows $1,3$ and columns $2,4$, is
\[
 \begin{pmatrix}x&y\\y&x\end{pmatrix},\qquad
 x=\frac{\sqrt\Theta}{4(1+\sqrt\Theta)},\qquad
 y=\frac1{4(1+\sqrt\Theta)}.
\]
Its positive graph eigenvalues are $2x$, $2y$ and $\frac12$. Thus every
initial metric gives
\[
 \lim_{\tau\to\infty}\tau(\lambda_1,\lambda_2,\lambda_3)
 =\Bigl(\frac{\min(1,\sqrt\Theta)}{2(1+\sqrt\Theta)},
        \frac{\max(1,\sqrt\Theta)}{2(1+\sqrt\Theta)},\frac12\Bigr).
\]
The constants retain the invariant $\Theta$ of the extension norms, but
not the extension classes themselves. The tight subgraph is a four-cycle,
so $u^*$ depends on the norms. The value $\frac12$ has eigenvector
$w=\frac14(-1,1,-1,1)$.
\end{example}

The three-factor extension $0\to L_1\oplus L_3\to E\to L_2\to0$ with
$\beta_{12},\beta_{32}\ne0$ and $V=1$ is of this kind. Its weight grading
is $(-\frac16,\frac13,-\frac16)$, and both multipliers equal $\frac16$. So
\[
 \lambda_1\sim\frac1{6\tau},\qquad\lambda_2\sim\frac1{2\tau},
\]
with limiting lines $\CC\diag(1,0,-1)$ and $\CC\diag(1,-2,1)$.

\subsection{A leading coefficient depending on the initial metric}

Let $V=1$, and choose pairwise nonisomorphic stable degree-zero bundles
$F_1,F_2,F_3,F_4$ of ranks $(2,1,1,2)$, with nonzero harmonic entries
$\beta_{12},\beta_{32},\beta_{34}$ and no others.

The weight grading is $w=(-\frac16,\frac13,-\frac13,\frac16)$. The edges
$12$ and $34$ are tight with multipliers $\frac13$. The edge $32$ has slack
$\frac13$, hence exponent $\frac43$; the edge $32$ is slack because
$M_1M_4=4>1=M_2M_3$. The multipliers on the tight edges are positive, so
$(\Gamma,(r_i))$ is nondegenerate. This is the region treated in
\cite[(3.24)--(3.26)]{HKKP2023}.

Theorem~\ref{thm:complete-spectrum} gives:
\begin{itemize}
\item two eigenvalues asymptotic to $1/(2\tau)$ (the tight components
 $\{1,2\}$ and $\{3,4\}$, each contributing $\frac12$), with limiting
 subspace $\operatorname{span}\{\diag(1,-2,0,0),\diag(0,0,2,-1)\}$;
\item one eigenvalue $\lambda_1\sim c(h(0))\tau^{-4/3}$, with limiting line
 $\CC\diag(1,1,-1,-1)$ and $c(h(0))=\frac23k^\infty_{32}(h)$.
\end{itemize}
The coefficient is unbounded on the metrics of one bundle. The explicit
family below makes this quantitative.

\begin{proposition}\label{prop:unequal-rank-family}
Normalize the three entries to have norm one, and for $\rho>0$ put
\[
 g_\rho=\diag(\rho^{1/2}\id_{F_1},\rho^{-1/2}\id_{F_2},
              \rho^{1/2}\id_{F_3},\rho^{-1/2}\id_{F_4}),\qquad
 h_\rho=g_\rho^*h_0.
\]
Then the flows starting at $h_\rho$ satisfy
\[
 c(h_\rho)=\frac2{27}\rho^{-2/3}(1+O(\rho)),\qquad\rho\downarrow0.
\]
\end{proposition}

\begin{proof}
Write $k_1=k_{12}$, $k_2=k_{32}$, $k_3=k_{34}$. The masses give
$p=(k_1/2,-k_1-k_2,k_2+k_3,-k_3/2)$ and
\[
 \dot k_1=-3k_1^2-2k_1k_2,\qquad\dot k_2=-2(k_1+k_3)k_2-4k_2^2,\qquad
 \dot k_3=-3k_3^2-2k_2k_3.
\]
The symmetric solution $k_1=k_3=u$, with initial values $u_0$ and
$r_0=k_2(T)/u_0$, is explicit. In the time $s$ with $ds/d\tau=u$ and $s(T)=0$, put
$D(s)=1+2r_0-2r_0e^{-s}$. Then
\[
 r(s)=\frac{r_0e^{-s}}{D(s)},\qquad
 u(s)=\frac{u_0e^{-3s}}{D(s)},\qquad
 \tau-T=\frac1{u_0}\Bigl(\frac{(1+2r_0)(e^{3s}-1)}3-r_0(e^{2s}-1)\Bigr),
\]
so
\[
 \tau^{4/3}k_2\longrightarrow B=\frac{r_0}{3^{4/3}u_0^{1/3}(1+2r_0)^{2/3}}.
\]
For $h_\rho$ all three initial conductances are $\rho^2$, so $u_0=\rho^2$ and
$r_0=1$, and the diagonal flow gives the coefficient
$\frac23B=2/(27\rho^{2/3})$. As in the proof of
Theorem~\ref{thm:complete-spectrum}, compare the flow from $h_\rho$ with the
corrected diagonal metrics along this solution, from time zero. Its initial
discrepancy from $h_\rho$ is $O(\rho^2)$, and its total residual is
$O(\rho)$. The barriers put the actual coefficient within
$e^{\pm C\rho}$ of the diagonal one.
\end{proof}

\subsection{A degenerate example: logarithmic separation}

Let $L_1,\ldots,L_4$ be pairwise nonisomorphic degree-zero line bundles,
with $V=1$, and take the zigzag $0\to L_1\oplus L_3\to E\to L_2\oplus L_4\to0$
with exactly $\beta_{12},\beta_{32},\beta_{34}$ nonzero. The masses are
equal, so $M_1M_4=M_2M_3$ and $(\Gamma,(r_i))$ is degenerate. The middle edge is
tight, but its multiplier is forced to vanish.
Theorem~\ref{thm:complete-spectrum} does not apply. The diagonal flow is
the scalar zigzag of \cite[Example~5.13]{HKKP2023}, and the transfer
theorem still applies.

\begin{theorem}\label{thm:zigzag-flow}
For every smooth initial metric,
\[
 \lambda_1(h(\tau))\sim\frac1{2\tau\log\tau},\qquad
 \lambda_2(h(\tau))\sim\lambda_3(h(\tau))\sim\frac1{2\tau},\qquad
 \liminf_{\tau\to\infty}\lambda_4(h(\tau))>0.
\]
The first line tends to $\CC\diag(1,1,-1,-1)$, and the second and third
together tend to
\[
 \operatorname{span}_{\CC}\{\diag(1,-1,0,0),\diag(0,0,1,-1)\}.
\]
The constants do not depend on the extension forms or the initial
metric.
\end{theorem}

\begin{proof}
With $k_1=k_{12}$, $k_2=k_{32}$, $k_3=k_{34}$, the diagonal flow is
\[
 \dot k_1=-2k_1(2k_1+k_2),\qquad\dot k_2=-2k_2(k_1+2k_2+k_3),\qquad
 \dot k_3=-2k_3(2k_3+k_2).
\]
Put $u=\sqrt{k_1k_3}$, $r=k_2/u$, $z=\sqrt{k_1/k_3}$, $\xi=(z-1)/(z+1)$,
and $ds/d\tau=2u$ with $s(T)=0$. Then
\[
 \frac{dr}{ds}=-r^2,\qquad\frac{d\xi}{ds}=-2\xi,\qquad
 \frac{d\log u}{ds}=-\Bigl(2+\frac{4\xi^2}{1-\xi^2}+r\Bigr).
\]
Hence $r=(s+r_0^{-1})^{-1}$ and $\xi=\xi_0e^{-2s}$ with $|\xi_0|<1$. There
is $C>0$ with $u=Ce^{-2s}(s+r_0^{-1})^{-1}(1+O(e^{-4s}))$, and
$\tau\sim e^{2s}(s+r_0^{-1})/(4C)$. Therefore
\[
 \tau k_1\to\tfrac14,\qquad\tau k_3\to\tfrac14,\qquad
 \tau\log\tau\,k_2\to\tfrac12,
\]
for every positive initial condition. The pairs $\{1,2\}$ and $\{3,4\}$
give two graph eigenvalues asymptotic to $1/(2\tau)$. On vectors constant on
the pairs, the remaining eigenvalue is asymptotic to $k_2$ (quotient masses
$2,2$). Eliminating the internal directions changes it by
$O(k_2^2/\min(k_1,k_3))=o(k_2)$. So
$\nu_1\sim(2\tau\log\tau)^{-1}$ and $\nu_2\sim\nu_3\sim(2\tau)^{-1}$.

Corollary~\ref{cor:coefficient-transfer} applies with
$f=(\tau\log\tau)^{-1}$ and $f=\tau^{-1}$, which proves the eigenvalue
statements. For the first line, a normalized eigensection has Rayleigh
quotient $O((\tau\log\tau)^{-1})$. Along the diagonal gauges $g_\tau$,
Lemma~\ref{lem:instantaneous-edges} gives $k_1,k_3\sim1/(4\tau)$. With the
form bound \eqref{eq:form-lower}, this forces $x_1=x_2$ and $x_3=x_4$ in
every limit. The
complementary two-dimensional space follows by orthogonality. Full-projection
convergence upgrades these identifications to $L^2\to C^k$ convergence.
\end{proof}

The four-factor first eigensection separates $L_1\oplus L_2$ from
$L_3\oplus L_4$, and the faster subspace separates the members of each
pair. Adjoining the identity to the full limiting small subspace recovers
all four summand projections.

\section{Two-step extensions in higher dimensions}\label{sec:higher-dim}

Let $(X,\omega)$ be a compact connected K\"ahler manifold of complex
dimension $n\ge1$, with the conventions of Section~\ref{sec:intro}. Let
$F_1,\ldots,F_m$ be pairwise nonisomorphic stable vector bundles of common
slope, with Hermitian--Einstein metrics $h_i$. Partition the indices as
$I\sqcup J$, and let
\[
 0\longrightarrow S=\bigoplus_{i\in I}F_i\longrightarrow E\longrightarrow
 U=\bigoplus_{j\in J}F_j\longrightarrow0
\]
be an extension, with harmonic representative $N=\sum\beta_{ij}$. Here each
entry $\beta_{ij}$ is a harmonic form in $\Omega^{0,1}(\Hom(F_j,F_i))$.
Assume that the graph $\Gamma$ of nonzero entries is connected.

\begin{theorem}\label{thm:higher-dim}
Let $E$ be as above, with masses $M_i=V\rk F_i$ and conductances
$k_{ij}=\lVert(gNg^{-1})_{ij}\rVert_{L^2}^2$. Along the normalized flow from
any smooth initial metric on $E$:
\begin{enumerate}
\item every small eigenvalue satisfies the two-sided comparison
 \eqref{eq:transfer} with the diagonal flow, with $\delta_T\to0$, and the
 full small spectral projection converges to $\mathcal K$ in every
 $L^2\to C^k$ norm;
\item if $(\Gamma,(r_i))$ is nondegenerate, conclusions (1)--(4) of
 Theorem~\ref{thm:complete-spectrum} hold;
\item for two factors, $\tau\lambda_1(h(\tau))\to\frac12$, as in
 Theorem~\ref{thm:flow-two-factor}.
\end{enumerate}
Corollaries~\ref{cor:coefficient-transfer} and~\ref{cor:balanced} and the
statements of Section~\ref{sec:flow-examples} hold on $X$ whenever the
required bundles and nonzero harmonic entries exist there.
\end{theorem}

The fixed-metric results of Sections~\ref{sec:two-factor}
and~\ref{sec:several}, including Theorem~\ref{thm:orders}, also hold on $X$
for two-step extensions. Section~\ref{sec:two-factor} is already stated
there. For more than two socle layers the families need not be integrable
when $n\ge2$. On
$\Sigma_1\times\Sigma_2$, for instance, entries $\beta_{12}$ and $\beta_{23}$ pulled
back from the two factors have a nonzero product class in
$H^{0,2}(\Hom(F_3,F_1))$.

\begin{proof}
We check each input used on curves.

\emph{Integrability.} Harmonic forms are $\db_0$-closed, and
$\beta_{ij}\wedge\beta_{kl}$ requires $j=k$, which cannot occur since
$j\in J$ and $k\in I$. So $\db_0N+N\wedge N=0$, and the same holds for
$gNg^{-1}$ and every entrywise rescaling.

\emph{Mean curvature.} The Chern connection of $(\db_0+N,h_0)$ is
$A_0+N-N^*$. Since $\db_0^*N=0$, the K\"ahler identity
$[\Lambda,\partial_0]=\sqrt{-1}\,\db_0^*$ gives $\Lambda\partial_0N=0$, and
likewise for $\db_0N^*$. Hence
$K-\kappa\id=\sqrt{-1}\Lambda(N-N^*)\wedge(N-N^*)$. In a local unitary
coframe with $N=\sum_\nu n_\nu\bar\theta^\nu$, this is
$\sum_\nu[n_\nu,n_\nu^*]$, as in \eqref{eq:two-C}. Its block traces give the
same $p_i$ as on a curve.

\emph{Fixed-metric comparison.} The proofs of Theorem~\ref{thm:relative}
and Proposition~\ref{prop:subspace} use the cancellation
$D_0^*[N,x]=[\db_0^*N,x]=0$, the spectral gap of $\Delta_0$ on $QH$,
elliptic estimates, and Sobolev embedding. All of these hold on $X$.

\emph{Order comparison.} Propositions~4.3 and~4.5 of \cite{HKKP2018},
used in \eqref{eq:flow-barriers}, are stated and proved on compact
K\"ahler manifolds of any dimension, by a pointwise maximum principle.

\emph{Correction.} The first variation of mean curvature at a
Hermitian--Einstein metric is $\Delta_\partial=\Delta_{\db}$ on
endomorphisms: the difference is $[K,\cdot]$, which vanishes there. The
mixed terms in Lemma~\ref{lem:general-correction} are
$O(\mathfrak e^{3/2})$ uniformly. So that lemma and Lemma~\ref{lem:layer-energy} hold
unchanged.

\emph{Convergence of the flow.} Every Jordan--H\"older filtration of $E$
has graded object isomorphic to $S\oplus U$. A saturated stable subsheaf
of slope $\mu$ in $E$ is reflexive, hence isomorphic to some $F_k$, and
$\Hom(F_k,E)=0$ for $k\in J$ since each extension class is nonzero. So the
Harder--Narasimhan--Seshadri graded object is locally free, and its singular
set is empty. By Sibley--Wentworth \cite[Theorem~1.1]{SibleyWentworth2015},
the flow converges modulo unitary gauge, smoothly on all of $X$, to an
admissible Yang--Mills connection $A_\infty$ on a bundle isomorphic to
$S\oplus U$. They credit the smooth convergence to Hong and Tian
\cite{HongTian2004}, and the
identification of the limit to Jacob and Sibley
\cite{Jacob2015,Sibley2015}.

$A_\infty$ is smooth and integrable. Its Hermitian--Einstein summands have
constants given by their slopes, and every holomorphic summand of
$\bigoplus F_i$ has slope $\mu$. So $A_\infty$ is Hermitian--Einstein, and it
is gauge equivalent to the product connection of $h_0$. For each $k$, the
infimum over unitary gauges of the $C^k$ distance to $A_\infty$ tends to
zero. Choosing near-optimal gauges with $k$ tending slowly to infinity gives,
for every large $\tau$, holomorphic isometries from $(E,h(\tau))$ to
$(E_0,h_0,\db^{(\tau)})$ with $\db^{(\tau)}\to\db_0$ in $C^\infty$. This is
the input used on curves.

\emph{Matching.} In Lemmas~\ref{lem:flow-matching}
and~\ref{lem:socle-matching}, the transported extension entry $\alpha_T$ is
$\db$-closed: in the block computation of $D_T^2=0$, the product
$\alpha_T\wedge\alpha_T$ lands in $\Hom(U,S)\circ\Hom(U,S)=0$. Hodge theory
on $X$ gives $\alpha_T=H\alpha_T+\db v_T$ with $v_T\to0$ in $C^\infty$.
Uniqueness up to diagonal maps uses only $L^2$-orthogonality of harmonic
and exact forms and $H^0(\Hom(U,S))=0$.

The remaining arguments are finite-dimensional or use only these inputs.
\end{proof}

\begin{example}\label{ex:products}
Let $X=\Sigma\times Y$, where $\Sigma$ is a curve of genus at least two and $Y$ is a
connected compact K\"ahler manifold, with a product K\"ahler form.
\begin{itemize}
\item For stable $F$ on $\Sigma$, the pullback of a Hermitian--Einstein metric is
 Hermitian--Einstein, since $\Lambda_\omega$ of a form pulled back from $\Sigma$
 is $\Lambda_{\omega_\Sigma}$ of that form. By the K\"unneth formula
 $H^0(X,\End\pi^*F)=\CC$, so $\pi^*F$ is stable.
\item Nonisomorphic stable bundles of equal slope pull back to such
 bundles, and
 \[
  H^1(X,\Hom(\pi^*G,\pi^*F))\supset H^1(\Sigma,\Hom(G,F))\ne0.
 \]
\item Harmonic forms pull back to harmonic forms.
\end{itemize}
Every example of Section~\ref{sec:flow-examples} therefore occurs on $X$,
with arbitrary initial metrics on $X$. The same holds for the sharpness
families of Section~\ref{sec:slope}. For the explicit constants,
$V=1$ means $\operatorname{Vol}(\Sigma)\operatorname{Vol}(Y)=1$.
\end{example}

For extensions with more than two steps in dimension $n\ge2$, the
Maurer--Cartan equation prevents all entries from being harmonic. A natural
replacement is the gauge $\db_0^*N=0$. In it, the entry $N_{13}$ acquires
the cup-product correction $-\db_0^*\mathbb G(N_{12}\wedge N_{23})$, with
$\mathbb G$ the Green operator, which exists
only when $[N_{12}]\cup[N_{23}]=0$. We do not treat this case.

\section{Repeated stable summands}\label{sec:repeated}

The relative estimate also applies when some stable summands are
isomorphic. The limiting kernel then contains matrices acting between
the repeated copies, and the comparison operator is a commutator
quadratic form on this larger space.

Retain the curve, scaling exponents, and harmonic upper-triangular family
\eqref{eq:several-family}, but allow repetitions among the stable
equal-slope bundles $F_1,\ldots,F_m$. Choose pairwise nonisomorphic
representatives $H_1,\ldots,H_s$ and identifications
\begin{equation}\label{eq:isotypical-splitting}
 G=\bigoplus_{i=1}^m F_i
   \simeq\bigoplus_{\alpha=1}^s H_\alpha\otimes\CC^{m_\alpha},
 \qquad \sum_\alpha m_\alpha=m.
\end{equation}
Here $H_\alpha\otimes\CC^{m_\alpha}$ is the direct sum of all copies
of $H_\alpha$, called its isotypical component. Fix a
Hermitian--Einstein metric on each $H_\alpha$ and give its copies the
transported metric. Thus the product metric on each isotypical
component is $h_{H_\alpha}\otimes I$. Write
$r_\alpha=\rk H_\alpha$.

Stability at equal slope gives
\begin{equation}\label{eq:matrix-algebra}
 \mathcal A:=H^0(X,\End G)
 =\bigoplus_{\alpha=1}^s
   \id_{H_\alpha}\otimes M_{m_\alpha}(\CC).
\end{equation}
Every element of $\mathcal A$ is parallel for the split Chern
connection: under the chosen identifications its nonzero entries
are constant multiples of the identity of a stable representative.
Consequently
\begin{equation}\label{eq:matrix-kernel}
 \begin{split}
 \mathcal K:=H^0(X,\End_0G)
 &=\left\{(X_\alpha)_\alpha:
             \sum_\alpha r_\alpha\tr X_\alpha=0\right\},\\
 d:=\dim_\CC\mathcal K&=\sum_\alpha m_\alpha^2-1,\qquad
 \lVert X\rVert^2=V\sum_\alpha r_\alpha\tr(X_\alpha^*X_\alpha).
 \end{split}
\end{equation}
Use $D_t=D_0+B_t$, $B_tx=[N_t,x]$, $\Delta_t=D_t^*D_t$,
$P,Q,\gamma,\eta_t$, and $p_*$ as before, now with this kernel.
Define the nonnegative self-adjoint operator $L_t$ on $\mathcal K$ by
\begin{equation}\label{eq:matrix-form}
 \langle x,L_tx\rangle=A_t(x):=\lVert[N_t,x]\rVert_{L^2}^2,
 \qquad L_t=PB_t^*B_tP|_{\mathcal K}.
\end{equation}
Let $\nu_1(t)\le\cdots\le\nu_d(t)$ denote its eigenvalues, including
zeros. Unlike \eqref{eq:graph-form}, this form may have cross terms
between extension entries contributing to the same matrix block.

\begin{theorem}[Comparison with the commutator form]\label{thm:matrix-relative}
If $\eta_t\le\sqrt\gamma/4$ and
$\gamma_t=(\sqrt\gamma-\eta_t)^2$, then exactly $d$ eigenvalues of
$\Delta_t$ lie below $\gamma/4$, counting zeros, and
\begin{equation}\label{eq:matrix-relative}
 \left(1-\frac{2\eta_t^2}{\gamma_t}\right)\nu_j(t)
 \le\lambda_j(t)\le\nu_j(t),\qquad 1\le j\le d,
 \qquad \lambda_{d+1}(t)\ge\frac\gamma4.
\end{equation}
In particular, $\lambda_j(t)/\nu_j(t)=1+O(t^{2p_*})$ for every
positive $\nu_j(t)$, uniformly over these $d$ modes.

For every $t>0$, the trace-free holomorphic endomorphism space of
$E_t$ has dimension
\begin{equation}\label{eq:matrix-centralizer}
 \dim_\CC\{x\in\mathcal K:[N_t,x]=0\}.
\end{equation}
Thus $E_t$ is simple if and only if the only elements of $\mathcal A$
commuting with $N_t$ are scalar identities.

For a unit small eigensection with eigenvalue $\lambda$, and for the
orthogonal projection $\Pi_t$ onto the full $d$-dimensional small
spectral subspace, one has, for every $k\ge0$,
\begin{equation}\label{eq:matrix-convergence}
 \lVert Qu\rVert_{C^k}\le C_k t^{p_*}\sqrt\lambda,
 \qquad
 \lVert\Pi_t-P\rVert_{L^2\to C^k}=O(t^{2p_*}).
\end{equation}
\end{theorem}

\begin{proof}
Multiplication on either side by a parallel endomorphism commutes
with the split Dolbeault Laplacian on bundle-valued forms. Since
$N_t$ is harmonic and $x\in\mathcal K$ is parallel, $B_tx=[N_t,x]$
is harmonic. In particular,
\[
 D_0^*B_tx=0\qquad(x\in\mathcal K).
\]
For $u=x+y$, $x\in\mathcal K$ and $y\in QH^1$, the same gap
estimate and cross-term bound as in \eqref{eq:complement-gap} give
\[
 \lVert D_tu\rVert^2
 \ge\left(1-\frac{2\eta_t^2}{\gamma_t}\right)A_t(x)
       +\frac{\gamma_t}{2}\lVert y\rVert^2.
\]
Also $\nu_d(t)\le\eta_t^2\le\gamma/16$. The min--max argument
in Theorem~\ref{thm:relative} therefore applies with $d$ in place
of $m-1$ and proves \eqref{eq:matrix-relative}. No commutativity of
$\mathcal A$ is used in this argument.

The displayed lower bound identifies the kernel for small $t$ with
$\{x\in\mathcal K:B_tx=0\}$. The parallel diagonal map $g_t$
preserves $\mathcal A$ by conjugation, preserves trace, and satisfies
$N_t=g_tN_1g_t^{-1}$. It therefore identifies the commutants for
different positive parameters. Since the bundles $E_t$ are also
holomorphically isomorphic, the kernel dimension formula holds for
all $t>0$.

To check smooth convergence, the complementary equation is still
\eqref{eq:schur}, with
$R_t=QB_t^*B_tP$ and $\lVert R_tx\rVert\le\eta_t\sqrt{A_t(x)}$.
Every $B_tx$ lies in the fixed finite-dimensional space of harmonic
$(0,1)$-forms with values in $\End_0G$. Equivalence of norms on
that space gives
\[
 \lVert B_tx\rVert_{H^\ell}\le C_\ell\sqrt{A_t(x)},\qquad
 \lVert R_tx\rVert_{H^\ell}
       \le C_\ell t^{p_*}\sqrt{A_t(x)}.
\]
This estimate remains valid in the presence of cancellations between
matrix entries. The uniform complementary inverse estimates and
Sobolev embedding used in Proposition~\ref{prop:subspace} now give
the first assertion of \eqref{eq:matrix-convergence}. The largest
small eigenvalue is $O(t^{2p_*})$; applying that assertion to an
orthonormal basis of the $d$ small eigensections and using the
projection formula for the resulting graph over $\mathcal K$
proves the second assertion.
\end{proof}

When all $m_\alpha=1$, \eqref{eq:matrix-form} is exactly the
weighted graph form. For repeated factors, connectedness of that
graph is insufficient for simplicity. For example, take
$G=H\oplus H$ with $H$ stable and
$N_t=t\beta E_{12}$, where $\beta$ is a nonzero harmonic scalar
$(0,1)$-form. The parallel nilpotent endomorphism $E_{12}$ commutes
with $N_t$, so this connected two-vertex family is not simple.
Here $E_{ij}$ denotes the identity map from the $j$th copy of $H$
to the $i$th, extended by zero on the other blocks.

The mean-curvature calculation \eqref{eq:several-curvature} uses
only harmonicity and the common Einstein constant, so it also applies
with repeated factors. In particular, the transported metrics
$g_t^*h_0$ have uniform mean-curvature error $O(t^{2p_*})$ on the
fixed bundle $E_1$.

\subsection{Recovery of factors and multiplicities}

After adjoining the scalar identity, the limiting small subspace is
the algebra $\mathcal A$ in \eqref{eq:matrix-algebra}. Its minimal
nonzero central idempotents are the projections onto the isotypical
components. The corresponding ideals have dimensions $m_\alpha^2$,
so their matrix sizes determine the multiplicities. Within such an
ideal, a minimal nonzero idempotent has rank one on the multiplicity
space and its holomorphic image is isomorphic to $H_\alpha$.
A decomposition of the central identity into orthogonal rank-one
projections gives $m_\alpha$ copies. Such a decomposition is not
canonical when $m_\alpha>1$.

The abstract algebra alone determines the multiplicities, but not the
isomorphism classes of the bundles $H_\alpha$; these come from its action
on $G$.

\subsection{A simple example with a repeated line bundle}\label{sec:repeated-example}

Let $X$ have genus at least two and volume one. Choose pairwise
nonisomorphic degree-zero line bundles $A,B,C$, with their
Hermitian--Einstein metrics, and put
$G=A\oplus B\oplus A\oplus C$, using the same metric on both
copies of $A$. Riemann--Roch gives
$h^1(X,L)=g(X)-1>0$ for every nontrivial degree-zero line bundle
$L$. Thus one can choose harmonic forms of unit $L^2$ norm
\[
 a\in\Omega^{0,1}(\Hom(B,A)),\quad
 b\in\Omega^{0,1}(\Hom(A,B)),\quad
 c\in\Omega^{0,1}(\Hom(C,A)).
\]
Set $N_{12}=ta$, $N_{23}=t^2b$, $N_{34}=tc$, and all other entries
to zero. These scales come from the scaling exponents $(4,3,1,0)$. An element of
$\mathcal K$ is
\[
 x=\begin{pmatrix}
 x_1&0&u&0\\0&x_2&0&0\\v&0&x_3&0\\0&0&0&x_4
 \end{pmatrix},\qquad \sum_{i=1}^4x_i=0,
 \qquad \lVert x\rVert^2=\sum_i|x_i|^2+|u|^2+|v|^2.
\]
The six nonzero commutator entries give
\begin{equation}\label{eq:repeated-example-form}
 \begin{split}
 A_t(x)={}&t^2\bigl(|x_2-x_1|^2+|x_4-x_3|^2+|u|^2+|v|^2\bigr)\\
         &+t^4\bigl(|x_3-x_2|^2+|v|^2\bigr).
 \end{split}
\end{equation}
This is positive definite on $\mathcal K$ for $t>0$, so
Theorem~\ref{thm:matrix-relative} proves that $E_t$ is simple.
Its upper-triangular filtration has stable degree-zero quotients
and a degree-zero line subbundle; hence it is strictly semistable.

The $u$ and $v$ modes have comparison eigenvalues $t^2$ and
$t^2+t^4$. The diagonal part is the four-vertex path with
conductances $t^2,t^4,t^2$. Its three positive eigenvalues are
\[
 2t^2,\qquad t^2+t^4\pm\sqrt{t^4+t^8}.
\]
Thus there are $2^2+1^2+1^2-1=5$ small positive Dolbeault
eigenvalues, in increasing order,
\begin{equation}\label{eq:repeated-example-rates}
 \begin{aligned}
 \lambda_1&=t^4(1+O(t^2)),\\
 \lambda_j&=t^2(1+O(t^2))&& (j=2,3),\\
 \lambda_j&=2t^2(1+O(t^2))&& (j=4,5).
 \end{aligned}
\end{equation}
The full limiting space contains the two off-diagonal maps between
the copies of $A$, in addition to the three trace-free diagonal
directions. After adjoining the identity, it is
$M_2(\CC)\oplus\CC\oplus\CC$, realized on $G$.

For completeness the limiting subspaces at each leading coefficient
can also be read from this example. The slow line is spanned by
$\diag(1,1,-1,-1)$. The coefficient-one subspace at order $t^2$
is spanned by $E_{13},E_{31}$, and the coefficient-two subspace
is spanned by $\diag(1,-1,0,0)$ and $\diag(0,0,1,-1)$.
Indeed, $t^{-2}L_t$ converges to the operator with exactly these
eigenspaces and respective eigenvalues $0,1,2$. In the complementary
equation \eqref{eq:schur}, the correction has operator norm
$O(t^4)$ by \eqref{eq:schur-bound}, which applies equally to
\eqref{eq:matrix-form}. After division by $t^2$ it tends to zero.
Together with \eqref{eq:matrix-convergence} and the known
multiplicities, this proves convergence of the three spectral projections
in every $L^2\to C^k$ norm. It does not single out a basis within
either two-dimensional limiting subspace.

\section{Spectral recovery for semistable bundles on curves}\label{sec:curve-recovery}

We now show that the harmonic families of Section~\ref{sec:several} can be
constructed from every simple strictly semistable bundle on a curve,
including those with repeated stable Jordan--H\"older factors. The
construction uses the normal form of Lemma~\ref{lem:normal-form}.

\subsection{A metric family realizing the graded object}

Let $E$ now be simple and strictly semistable. Fix a Jordan--H\"older
filtration \eqref{eq:curve-filtration} with stable quotients $F_i$,
allowing repetitions. On a smooth curve a saturated
filtration has locally free terms and quotients, so the lemma applies.
All quotients have slope $\mu(E)$. Choose Hermitian--Einstein metrics
$h_i$ on them, choosing the same metric on isomorphic copies under
fixed identifications, and put $h_0=\bigoplus_i h_i$ on
$G=\operatorname{gr}(E)=\bigoplus_iF_i$.
Write $G=\bigoplus_\alpha H_\alpha\otimes\CC^{m_\alpha}$ as in
\eqref{eq:isotypical-splitting}, and set
$d=\sum_\alpha m_\alpha^2-1$.
Apply Lemma~\ref{lem:normal-form} using these metrics to obtain
$\Phi$ and $\beta$.

Choose any strictly decreasing integers $\ell_1>\cdots>\ell_m=0$, scaling
exponents as in Section~\ref{sec:several}, and set
\begin{equation}\label{eq:recovery-gauges}
 g_t=\diag(t^{\ell_i}\id_{F_i}),\qquad
 \Psi_t=g_t\Phi^{-1}:E\longrightarrow G,\qquad
 h_t=\Psi_t^*h_0,\qquad t>0.
\end{equation}
Define $p_{ij}=\ell_i-\ell_j$ and
$p_* =\min_{\beta_{ij}\ne0}p_{ij}$. This minimum exists: if all
entries vanished, $E$ would be a direct sum of $m\ge2$ stable
bundles and hence would not be simple. Let $L_t$ be the operator
\eqref{eq:matrix-form} on $\mathcal K_G=H^0(X,\End_0G)$, with
$N_t=\sum_{i<j}t^{p_{ij}}\beta_{ij}$, and let $\nu_j(t)$ be its
$d$ eigenvalues in increasing order.

\begin{theorem}[Recovery along a constructed metric family]\label{thm:curve-recovery}
The metrics \eqref{eq:recovery-gauges} form an approximate
Hermitian--Einstein family on $E$. More precisely, with
$\kappa=2\pi\mu(E)/V$,
\begin{equation}\label{eq:recovery-curvature}
 \lVert K_{h_t}-\kappa\id_E\rVert_{L^\infty,\op,h_t}
 =O(t^{2p_*}).
\end{equation}
The commutator form $L_t$ is positive definite. Exactly $d$ eigenvalues of the
trace-free endomorphism Laplacian tend to zero, all are positive for
$t>0$, and
\begin{equation}\label{eq:recovery-rates}
 \frac{\lambda_j(h_t)}{\nu_j(t)}=1+O(t^{2p_*}),
 \qquad 1\le j\le d,\qquad
 \lambda_{d+1}(h_t)\ge\delta>0
\end{equation}
for sufficiently small $t$.

Let $\Pi_t^E$ be the orthogonal projection onto these $d$
eigenvalues, and let $P_G$ be the $L^2(h_0)$ projection onto
\[
 \mathcal K_G=H^0(X,\End_0G)
 =\left\{(X_\alpha)_\alpha:\sum_\alpha
       \rk(H_\alpha)\tr X_\alpha=0\right\}.
\]
Conjugation by $\Psi_t$ defines a unitary identification
$\mathcal U_tu=\Psi_tu\Psi_t^{-1}$ of the endomorphism $L^2$
spaces. For every $k\ge0$,
\begin{equation}\label{eq:recovery-projection}
 \lVert\mathcal U_t\Pi_t^E\mathcal U_t^{-1}-P_G
 \rVert_{L^2(h_0)\to C^k(h_0)}=O(t^{2p_*}).
\end{equation}
Adjoining $\CC\id_G$ to the full limiting space gives the algebra
$\bigoplus_\alpha M_{m_\alpha}(\CC)$ acting on $G$. Its minimal
nonzero central idempotents and matrix sizes recover the isotypical components and
their multiplicities; the holomorphic images of minimal nonzero
idempotents recover the stable factors up to isomorphism.

If the $F_i$ are pairwise nonisomorphic, then $d=m-1$, $L_t$ is the
connected weighted graph Laplacian \eqref{eq:graph-form}, and the
quotient graphs of Theorem~\ref{thm:orders} determine every leading
coefficient and corresponding limiting spectral subspace.
\end{theorem}

\begin{proof}
The holomorphic structure transported by $\Psi_t$ is
\begin{equation}\label{eq:recovery-operator}
 \Psi_t\db_E\Psi_t^{-1}
 =g_t(\db_0+\beta)g_t^{-1}
 =\db_0+\sum_{i<j}t^{p_{ij}}\beta_{ij}.
\end{equation}
This is precisely the family of Section~\ref{sec:repeated}, with
its fixed product metric. The map $\Psi_t$ is a holomorphic
isometry from $(E,h_t)$ to that bundle. Naturality of the Chern
connection and \eqref{eq:several-curvature} prove
\eqref{eq:recovery-curvature}; the pointwise operator norms agree
under this isometry.

Conjugation by a holomorphic isometry intertwines the Dolbeault
operators on endomorphisms and their adjoints. It preserves the
trace-free condition and the Hilbert--Schmidt metrics. Thus
$\mathcal U_t$ identifies the endomorphism Laplacians unitarily.
Since $E$ is simple, Theorem~\ref{thm:matrix-relative} makes $L_t$
positive definite and gives \eqref{eq:recovery-rates}, with
$\delta=\gamma/4$ in the notation of that theorem.
The same theorem proves \eqref{eq:recovery-projection}. When all
multiplicities equal one, Theorems~\ref{thm:relative} and
\ref{thm:orders} give connectedness and the graph contraction
description of individual decay orders and limiting spectral subspaces.

Finally, stability of the distinct representatives gives
\[
 \mathcal K_G\oplus\CC\id_G
 =H^0(X,\End G)=\bigoplus_\alpha
     \id_{H_\alpha}\otimes M_{m_\alpha}(\CC).
\]
The central idempotent for $\alpha$ has image
$H_\alpha\otimes\CC^{m_\alpha}$. Its matrix ideal has dimension
$m_\alpha^2$, and a minimal nonzero idempotent in that ideal has
image isomorphic to $H_\alpha$. Decomposing the central identity
into orthogonal rank-one projections recovers $m_\alpha$ copies.
Their direct sum over $\alpha$ is $\operatorname{gr}(E)$.
The Jordan--H\"older theorem makes
the isomorphism class of this graded bundle independent of the chosen
filtration \cite{Kobayashi1987}.
\end{proof}

The construction depends on a filtration, quotient metrics, and integer
scaling exponents; the scope of such statements is summarized in
Section~\ref{sec:earlier-work}.

\section{A slope inequality and its sharpness}\label{sec:slope}

This section proves Theorem~\ref{thm:main} and shows that its constants are
optimal in every rank.

\subsection{Projections associated with coherent subsheaves}

We recall the analytic input that permits the use of a possibly singular
subsheaf as a test section. The Sobolev spaces below are defined using any
smooth background connection on $\End E$.

\begin{lemma}[Subsheaf degree formula]\label{lem:projection}
Let $\mathcal S\subset E$ be a saturated coherent subsheaf of rank $s$.
Outside an analytic subset $Z\subset X$ of complex codimension at least
two, the sheaf $\mathcal S$ is a holomorphic subbundle of $E$. Its
$h$-orthogonal projection $P_{\mathcal S}$ extends as an element of
$W^{1,2}(X,\End E)$ and satisfies, almost everywhere,
\[
 P_{\mathcal S}^2=P_{\mathcal S}=P_{\mathcal S}^{*h},\qquad \tr P_{\mathcal S}=s,
 \qquad (\id_E-P_{\mathcal S})\db_{\End E}P_{\mathcal S}=0.
\]
Moreover,
\begin{equation}\label{eq:degree}
 2\pi\deg_\omega\mathcal S
 =\int_X\tr(K_hP_{\mathcal S})\dvol
  -\lVert\db_{\End E}P_{\mathcal S}\rVert_{L^2}^2.
\end{equation}
\end{lemma}

\begin{proof}
The sheaves $\mathcal S$ and $E/\mathcal S$ are torsion-free, so both
are locally free away from an analytic subset of codimension at least
two. On this complement, a local splitting of the resulting exact
sequence makes $\mathcal S$ a subbundle. The projection identities follow
there from orthogonality and holomorphicity.

For a smooth subbundle, the Gauss--Codazzi formula gives
\eqref{eq:degree}. With the Hilbert--Schmidt convention used here,
$\lVert\db_{\End E}P_{\mathcal S}\rVert_{L^2}^2$ is the squared $L^2$ norm of
the second fundamental form. The factor $2\pi$ follows from
$c_1=\frac{\sqrt{-1}}{2\pi}\tr F$ and the contraction identity
\[
 \alpha\wedge\frac{\omega^{n-1}}{(n-1)!}
 =(\Lambda_\omega\alpha)\frac{\omega^n}{n!}
\]
for a $(1,1)$-form $\alpha$.

For a saturated coherent subsheaf, we use the standard extension of this
formula in the theory of weakly holomorphic projections
\cite{UhlenbeckYau1986,Kobayashi1987}. In particular, the projection is
bounded and its second fundamental form is square-integrable. A
regularization argument establishing this integrability and preservation
of the degree is given in \cite[Sections~2--3]{Jacob2014}. Saturation
ensures that there is no divisorial contribution from the singular set.
The regularized degree formula consequently gives
\eqref{eq:degree} on $X\setminus Z$.

For completeness, square-integrability of the second fundamental form
also gives the required Sobolev statement. Since $P_{\mathcal S}$ is self-adjoint,
its $(1,0)$ covariant derivative is the adjoint of its $(0,1)$ derivative,
so its full first covariant derivative is square-integrable on
$X\setminus Z$. An analytic set of complex codimension at least two has
zero $W^{1,2}$ capacity. Applying cutoff functions with Dirichlet energy
tending to zero, and using boundedness of $P_{\mathcal S}$, extends it across $Z$
as a $W^{1,2}$ section. The integral formula is unchanged because $Z$
has measure zero.
\end{proof}

\begin{remark}\label{rem:weak}
The condition $(\id_E-P_{\mathcal S})\db_{\End E}P_{\mathcal S}=0$ says that the image of
$P_{\mathcal S}$ is holomorphic. It does not say that $P_{\mathcal S}$ is a holomorphic
endomorphism. The latter would require $\db_{\End E}P_{\mathcal S}=0$ and would
give a holomorphic direct-sum decomposition of $E$.
\end{remark}

The following elementary estimate uses the trace-free condition on the
curvature, in addition to the spectrum of the projection.

\begin{lemma}\label{lem:matrix}
Let $A$ be a self-adjoint trace-free endomorphism of an $r$-dimensional
Hermitian vector space, and let $\Pi$ be an orthogonal projection of
rank $s$, where $0<s<r$. Set $q=\Pi-\frac{s}{r}\id$. Then
\begin{equation}\label{eq:matrix}
 |q|^2=\frac{s(r-s)}{r},\qquad
 |\tr(Aq)|\le\min\{s,r-s\}\lVert A\rVert_{\op}.
\end{equation}
The coefficient $\min\{s,r-s\}$ in the second inequality is optimal
among such matrices and projections.
\end{lemma}

\begin{proof}
The eigenvalues of $q$ are $(r-s)/r$, with multiplicity $s$, and
$-s/r$, with multiplicity $r-s$, which proves the norm identity.
Since $\tr A=0$,
\[
 \tr(Aq)=\tr(A\Pi)=-\tr(A(\id-\Pi)).
\]
Estimating the two compressed traces by their respective dimensions
gives \eqref{eq:matrix}.

To see optimality, suppose first that $s\le r-s$ and take
$A=\id$ on $\operatorname{im}\Pi$ and
$A=-\frac{s}{r-s}\id$ on its orthogonal complement. Then
$\tr A=0$, $\lVert A\rVert_{\op}=1$, and $\tr(Aq)=s$.
For $s\ge r-s$, take $A=\frac{r-s}{s}\id$ on
$\operatorname{im}\Pi$ and $A=-\id$ on its complement.
\end{proof}

\subsection{Proof of the slope estimate}

\begin{proof}[Proof of Theorem~\ref{thm:main}]
Let $P_{\mathcal S}$ be the projection supplied by Lemma~\ref{lem:projection}, and
put
\[
 q=P_{\mathcal S}-\frac{s}{r}\id_E.
\]
This is a trace-free $W^{1,2}$ endomorphism and
$\db_{\End E}q=\db_{\End E}P_{\mathcal S}$. Subtracting
\eqref{eq:degree} from $s/r$ times the degree formula for $E$ gives
\begin{equation}\label{eq:identity}
 2\pi s\bigl(\mu_\omega(E)-\mu_\omega(\mathcal S)\bigr)
 =\lVert\db_{\End E}q\rVert_{L^2}^2
  -\int_X\tr(K_h^0q)\dvol.
\end{equation}
The scalar part of $K_h$ disappears because $\tr q=0$ pointwise;
no assumption on $\tr K_h$ is used.

The Rayleigh inequality in \eqref{eq:lambda} extends to trace-free
$W^{1,2}$ sections by density. Explicitly, a smooth approximation may
be made trace-free by subtracting its scalar part. Hence
Lemma~\ref{lem:matrix} gives
\begin{equation}\label{eq:rayleigh}
 \lVert\db_{\End E}q\rVert_{L^2}^2
 \ge\lambda_1(h)\lVert q\rVert_{L^2}^2
 =\lambda_1(h)V\frac{s(r-s)}{r}.
\end{equation}
The same lemma, applied pointwise to $A=K_h^0$, gives
\begin{equation}\label{eq:error}
 \left|\int_X\tr(K_h^0q)\dvol\right|
 \le V\min\{s,r-s\}\varepsilon_1(h).
\end{equation}
Combining \eqref{eq:identity}--\eqref{eq:error} and dividing by
$2\pi s$ proves \eqref{eq:main}, since
\[
 \frac{r\min\{s,r-s\}}{s(r-s)}
 =\frac{r}{\max\{s,r-s\}}.
\]

Finally,
\[
 \max_{1\le s\le r-1}\frac{r}{\max\{s,r-s\}}
 =\frac{r}{\lceil r/2\rceil}=c_r.
\]
Under \eqref{eq:criterion}, every saturated proper subsheaf has strictly
smaller slope. An arbitrary subsheaf has slope at most that of its
saturation, so $E$ is stable.
\end{proof}

\begin{remark}
The coefficient for all subsheaf ranks is $c_r=2$ in even rank and
$c_r=2r/(r+1)$ in odd rank. For example, rank three gives $c_3=3/2$.
For rank four the coefficient in \eqref{eq:main} is $4/3$ for
subsheaves of rank one or three, and $2$ for subsheaves of rank two.
Pointwise optimality in Lemma~\ref{lem:matrix} does not by itself give
optimality of the geometric estimate, which requires simultaneous equality
in the curvature estimate and the Rayleigh inequality.
Corollary~\ref{cor:sharpness} shows that both become asymptotically sharp on
suitable nonsplit extensions. Every coefficient in \eqref{eq:main}, and
$c_r$ itself, is optimal in every rank.
\end{remark}

\subsection{Consequences}

\begin{corollary}[A uniform slope margin]\label{cor:margin}
If $\gamma(h)=\lambda_1(h)-c_r\varepsilon_1(h)>0$, then every coherent
subsheaf $\mathcal S\subset E$ of rank $0<s<r$ satisfies
\[
 \mu_\omega(E)-\mu_\omega(\mathcal S)
 \ge\frac{V(r-s)}{2\pi r}\gamma(h)
 \ge\frac{V}{2\pi r}\gamma(h).
\]
In particular, if $E$ is simple and $K_h^0=0$, then
\begin{equation}\label{eq:HE-margin}
 \mu_\omega(E)-\mu_\omega(\mathcal S)
 \ge\frac{V(r-s)}{2\pi r}\lambda_1(h)>0.
\end{equation}
\end{corollary}

\begin{proof}
For a saturated subsheaf, replace the rank-dependent coefficient in
\eqref{eq:main} by $c_r$. For any other subsheaf, apply the result to its
saturation, which has the same rank and at least the same degree. If
$K_h^0=0$, use \eqref{eq:main} directly and the positivity of
$\lambda_1(h)$ for a simple bundle.
\end{proof}

\begin{remark}
An irreducible Hermitian--Einstein connection has only scalar
holomorphic endomorphisms by the Bochner formula, so
\eqref{eq:HE-margin} applies in the classical setting. Conversely, a
stable bundle admits a Hermitian--Einstein metric by
\cite{Donaldson1985,UhlenbeckYau1986}, and that metric satisfies
\eqref{eq:criterion}. Thus existence of a metric satisfying the
criterion is equivalent to stability. The existence direction here is
the classical correspondence; the estimate does not construct such a
metric from an arbitrary one.
\end{remark}

\begin{corollary}[Spectral collapse]\label{cor:collapse}
Suppose that $E$ is simple and strictly semistable. Fix a saturated
subsheaf $\mathcal S\subset E$ of rank $0<s<r$ with
$\mu_\omega(\mathcal S)=\mu_\omega(E)$. Then every smooth Hermitian
metric $h$ satisfies
\begin{equation}\label{eq:collapse}
 0<\lambda_1(h)
 \le\frac{r}{\max\{s,r-s\}}\varepsilon_1(h)
 \le c_r\varepsilon_1(h).
\end{equation}
Consequently, if $h_j$ is any sequence with
$\varepsilon_1(h_j)\longrightarrow0$, then
$\lambda_1(h_j)\longrightarrow0$.
\end{corollary}

\begin{proof}
Strict semistability provides an equal-slope proper subsheaf. Its
saturation still has equal slope, by semistability. Substituting this
subsheaf in \eqref{eq:main} proves the upper bound. Simplicity supplies
strict positivity for each fixed metric.
\end{proof}

Semistable bundles on compact K\"ahler manifolds admit approximate
Hermitian--Einstein metrics \cite{Jacob2014}. In the normalizations of
this paper, these are metrics satisfying
\[
 \lVert K_{h_j}-\kappa\id_E\rVert_{L^\infty,\op}\longrightarrow0,
 \qquad \kappa=\frac{2\pi\mu_\omega(E)}{V}.
\]
Their trace-free mean curvatures tend to zero, so
Corollary~\ref{cor:collapse} applies to every such sequence on a simple
strictly semistable bundle. The conclusion also holds under the weaker
assumption in the corollary, without convergence of the scalar part of
the mean curvature. For a nonsimple bundle, the lowest eigenvalue in
\eqref{eq:lambda} is already zero for every metric; no conclusion about
its first positive eigenvalue is asserted.

\begin{proposition}[Scalar changes of metric]\label{prop:conformal}
If $f$ is a smooth real function on $X$ and $\widetilde h=e^f h$, then
\[
 \lambda_1(\widetilde h)=\lambda_1(h),\qquad
 K_{\widetilde h}^0=K_h^0,\qquad
 \varepsilon_1(\widetilde h)=\varepsilon_1(h).
\]
\end{proposition}

\begin{proof}
Multiplication of $h$ by a scalar function leaves the induced metric on
$\End E$ unchanged, and the Dolbeault operator is fixed. It therefore
leaves both norms in \eqref{eq:lambda} unchanged. The Chern curvatures
of $\widetilde h$ and $h$ differ by a scalar $(1,1)$-form times
$\id_E$, so their trace-free mean curvatures agree. The operator norm
of an endomorphism is also unchanged by a scalar change of metric.
\end{proof}

\subsection{Sharpness in every rank}

\begin{corollary}[Geometric sharpness in every rank]\label{cor:sharpness}
Let $X$ be a curve of genus at least two, and let $r\ge2$ and
$1\le s\le r-1$. There are a simple strictly semistable bundle $E$ of rank
$r$ with a saturated subbundle of rank $s$ and slope $\mu(E)$, and metrics
$h_t$ on $E$ with $\varepsilon_1(h_t)\to0$, such that
\[
 \frac{\lambda_1(h_t)}{\varepsilon_1(h_t)}\longrightarrow\frac r{\max\{s,r-s\}},
 \qquad
 \frac{\lambda_1(h_t)}{\lVert K^0_{h_t}\rVert_{L^\infty,\op}}
 \longrightarrow\frac r{\max\{s,r-s\}}.
\]
Consequently the coefficient $r/\max\{s,r-s\}$ in \eqref{eq:main} cannot be
decreased for any $(r,s)$, and the constant $c_r$ in \eqref{eq:criterion} is
optimal in every rank, also for the criterion with the $L^\infty$ norm in
place of $\varepsilon_1$.
\end{corollary}

\begin{proof}
Choose stable $F$ of rank $a=s$ and $G$ of rank $b=r-s$, both of degree
zero, and nonisomorphic. A nonzero map between them would be an
isomorphism, so $h^0(\Hom(G,F))=0$. By Riemann--Roch,
$h^1(\Hom(G,F))=ab(g-1)>0$, so there is a nonzero harmonic $\beta$.

With $v=Rz$, consider on $E_t$ the metric $h_t'=h_0e^{-t^2v}$, where
$h_t'(x,y)=h_0(e^{-t^2v}x,y)$. As in the proof of
Lemma~\ref{lem:flow-approximation}, the first variation of mean curvature
at $h_0$ is $\Delta_0$. For $t\ne0$ the operator changes by $O(t)$. Hence
\[
 K_{h_t'}=\kappa\id+t^2(C-\Delta_0v)+O_{C^0}(t^3)
 =\kappa\id+ct^2q+O_{C^0}(t^3),
\]
since $\Delta_0v=z=C-cq$. The endomorphism $K^0_{h_t'}$ is
$h_t'$-self-adjoint, so its operator norm is its spectral radius. The
eigenvalues of $ct^2q$ are $ct^2b/r$ and $-ct^2a/r$. Therefore
\[
 \lVert K^0_{h_t'}\rVert_{\op}=\frac{ct^2\max\{a,b\}}r+O(t^3)
\]
uniformly on $X$, and the same holds for $\varepsilon_1$ and the $L^\infty$
norm. Since $e^{-Ct^2}h_0\le h_t'\le e^{Ct^2}h_0$,
\eqref{eq:metric-spectral-comparison} and Theorem~\ref{thm:two-factor} give
$\lambda_1(h_t')=ct^2(1+O(t^2))$. The ratios therefore tend to
$r/\max\{a,b\}$.

Transport by $g_t$ in \eqref{eq:two-gauge}. This gives the metrics
$g_t^*h_t'$ on the fixed bundle $E=E_1$, with the same $\lambda_1$,
$\varepsilon_1$, and $L^\infty$ norm. The subbundle $F\subset E$ has rank
$s$ and equal slope, so \eqref{eq:main} gives
$\lambda_1\le(r/\max\{s,r-s\})\varepsilon_1$ for every metric on $E$. A
smaller coefficient would contradict the limit. For $s=\lfloor r/2\rfloor$
the coefficient is $c_r$.
\end{proof}

\begin{remark}
The metrics $h_t'$ make both estimates in the proof of
Theorem~\ref{thm:main} asymptotically sharp. The Rayleigh bound is sharp
because $u_t\to q/\sqrt{V_q}$. The trace-free matrix estimate is sharp because
the correction replaces $C$ by its projection $cq$, which is constant on each
block with the larger absolute value on the smaller block.

Without the correction, on a curve,
$\lambda_1/\varepsilon_1\to(r/ab)\int\sum_k\sigma_k^2/\int\sigma_{\max}^2$,
where the $\sigma_k$ are the pointwise singular values of $\beta$. This
reaches $r/\max\{a,b\}$ exactly when the singular values agree almost
everywhere. It happens, for example, for line-bundle extensions, when
$\min\{a,b\}=1$, and for $\beta=\gamma\otimes\id_W$ on
$F=L_1\otimes W$, $G=L_2\otimes W$, with $L_1L_2^{-1}$ chosen so that
$W\otimes L_1L_2^{-1}\not\cong W$. It does not give the $L^\infty$
statement: in rank two, $|\beta|^2$ vanishes somewhere, since $\bar\beta$
corresponds to a holomorphic section of a line bundle of positive degree.

For line-bundle extensions, put $f=|\beta|^2-V^{-1}\lVert\beta\rVert_{L^2}^2$
and let $R_{\rm sc}$ be the inverse of the scalar $\db^*\db$ on mean-zero
functions. Since $z=2fq$ and $q$ is parallel,
$d=\frac4V\langle f,R_{\rm sc}f\rangle$. More generally, $d=0$ exactly when
$C=cq$. In that case $q$ is an exact eigensection with eigenvalue $ct^2$,
which is the first eigenvalue for small $|t|$.
\end{remark}


\section{Open problems}\label{sec:limits}

\begin{enumerate}
\item \emph{Arbitrary approximate Hermitian--Einstein sequences.}
 Corollary~\ref{cor:collapse} concerns eigenvalues only. For a general
 sequence the Hilbert norms and adjoints depend on $h_j$, and a statement
 about eigensections must specify a gauge, the limiting holomorphic
 structure, and the topology of convergence.
\item \emph{The degenerate case.} When a tight edge carries a vanishing
 multiplier, iterated logarithms can appear. Theorem~\ref{thm:transfer}
 reduces the question to the diagonal flow, but the sharp asymptotics are
 known here only for the zigzag of Theorem~\ref{thm:zigzag-flow}.
\item \emph{Individual coefficients.} At each faster rate we show that the
 sum of the coefficients depends on the initial metric. Each coefficient
 does so when the edges of that exponent form a forest on the tight
 components, parallel edges counting as cycles: prescribing the constants
 $Z_{\mathcal C}$ along the forest scales all these conductances by a common
 factor. The general case is open.
\item \emph{Repeated factors along the flow.} Then the diagonal flow is
 replaced by a moment-map flow for $\prod_\alpha GL(m_\alpha)$. Lam
 \cite{Lam2026} gives the metric asymptotics up to a bounded term; the
 spectral questions are open.
\item \emph{More than two steps in dimension at least two.} The
 Maurer--Cartan equation requires cup-product corrections, and the
 matching argument needs a Kuranishi-slice version.
\item \emph{Singular graded objects,} where the flow bubbles
 \cite{SibleyWentworth2015}.
\item \emph{Variation of the polarization.} Clarke--Tipler
 \cite{ClarkeTipler2023} study convergence of Hermitian Yang--Mills
 connections to the graded bundle as the polarization varies. Our
 fixed-polarization calculation gives no effective wall-crossing radius.
\end{enumerate}

\begin{thebibliography}{99}

\bibitem[AB83]{AtiyahBott1983}
Michael~F. Atiyah and Raoul Bott, \emph{The {Y}ang--{M}ills equations
over {R}iemann surfaces}, Philos. Trans. Roy. Soc. London Ser. A
\textbf{308} (1983), no.~1505, 523--615.

\bibitem[CT23]{ClarkeTipler2023}
Andrew Clarke and Carl Tipler, \emph{Semi-stability and local wall-crossing
for {H}ermitian {Y}ang--{M}ills connections}, arXiv:2304.05245 (2023).

\bibitem[Das92]{Daskalopoulos1992}
Georgios D. Daskalopoulos, \emph{The topology of the space of stable
bundles on a compact {R}iemann surface}, J. Differential Geom.
\textbf{36} (1992), no.~3, 699--746.

\bibitem[DW04]{DaskalopoulosWentworth2004}
Georgios~D. Daskalopoulos and Richard~A. Wentworth, \emph{Convergence
properties of the {Y}ang--{M}ills flow on {K}\"ahler surfaces}, J. Reine
Angew. Math. \textbf{575} (2004), 69--99; correction, ibid.
\textbf{606} (2007), 235--236.

\bibitem[Don83]{Donaldson1983}
Simon~K. Donaldson, \emph{A new proof of a theorem of {N}arasimhan and
{S}eshadri}, J. Differential Geom. \textbf{18} (1983), no.~2, 269--277.

\bibitem[Don85]{Donaldson1985}
Simon~K. Donaldson, \emph{Anti self-dual {Y}ang--{M}ills connections over
complex algebraic surfaces and stable vector bundles}, Proc. London Math.
Soc. (3) \textbf{50} (1985), 1--26.

\bibitem[HKKP18]{HKKP2018}
Fabian Haiden, Ludmil Katzarkov, Maxim Kontsevich, and Pranav Pandit,
\emph{Iterated logarithms and gradient flows}, arXiv:1802.04123 (2018).

\bibitem[HKKP23]{HKKP2023}
Fabian Haiden, Ludmil Katzarkov, Maxim Kontsevich, and Pranav Pandit,
\emph{Semistability, modular lattices, and iterated logarithms},
J. Differential Geom. \textbf{123} (2023), no.~1, 21--66,
arXiv:1706.01073.

\bibitem[HL10]{HuybrechtsLehn2010}
Daniel Huybrechts and Manfred Lehn, \emph{The geometry of moduli spaces
of sheaves}, second edition, Cambridge Mathematical Library, Cambridge
University Press, Cambridge, 2010.

\bibitem[HT04]{HongTian2004}
Min-Chun Hong and Gang Tian, \emph{Asymptotical behaviour of the
{Y}ang--{M}ills flow and singular {Y}ang--{M}ills connections}, Math. Ann.
\textbf{330} (2004), no.~3, 441--472.

\bibitem[Jac14]{Jacob2014}
Adam Jacob, \emph{Existence of approximate {H}ermitian--{E}instein
structures on semi-stable bundles}, Asian J. Math. \textbf{18} (2014),
no.~5, 859--884, arXiv:1012.1888.

\bibitem[Jac15]{Jacob2015}
Adam Jacob, \emph{The limit of the {Y}ang--{M}ills flow on semi-stable
bundles}, J. Reine Angew. Math. \textbf{709} (2015), 1--13,
arXiv:1104.4767.

\bibitem[Kat95]{Kato1995}
Tosio Kato, \emph{Perturbation theory for linear operators}, second edition,
Classics in Mathematics, Springer-Verlag, Berlin, 1995.

\bibitem[Kob87]{Kobayashi1987}
Shoshichi Kobayashi, \emph{Differential geometry of complex vector
bundles}, Publications of the Mathematical Society of Japan, vol.~15
(Kan\^o Memorial Lectures~5), Iwanami Shoten, Tokyo, and Princeton
University Press, Princeton, NJ, 1987.

\bibitem[Lam26]{Lam2026}
Shing Tak Lam, \emph{Semistability and asymptotics of geometric flows},
arXiv:2606.09235 (2026).

\bibitem[L{\"u}b83]{Lubke1983}
Martin L{\"u}bke, \emph{Stability of {E}instein--{H}ermitian vector
bundles}, Manuscripta Math. \textbf{42} (1983), 245--257.

\bibitem[Mum63]{Mumford1963}
David Mumford, \emph{Projective invariants of projective structures and
applications}, Proc. Internat. Congr. Mathematicians (Stockholm, 1962),
Inst. Mittag-Leffler, Djursholm, 1963, 526--530.

\bibitem[NLGSSS19]{Nacson2019}
Mor~Shpigel Nacson, Jason~D. Lee, Suriya Gunasekar, Pedro H.~P. Savarese, Nathan Srebro, and Daniel Soudry, \emph{Convergence of gradient
descent on separable data}, Proceedings of the 22nd International Conference
on Artificial Intelligence and Statistics (AISTATS 2019), Proc. Mach. Learn.
Res., vol.~89, 2019, arXiv:1803.01905.

\bibitem[NS65]{NarasimhanSeshadri1965}
M.~S. Narasimhan and C.~S. Seshadri, \emph{Stable and unitary vector
bundles on a compact {R}iemann surface}, Ann. of Math. (2) \textbf{82}
(1965), 540--567.

\bibitem[R{\aa}d92]{Rade1992}
Johan R{\aa}de, \emph{On the {Y}ang--{M}ills heat equation in two
and three dimensions}, J. Reine Angew. Math. \textbf{431} (1992),
123--163.

\bibitem[RS89]{RothblumSchneider1989}
Uriel~G. Rothblum and Hans Schneider, \emph{Scalings of matrices which
have prespecified row sums and column sums via optimization},
Linear Algebra Appl. \textbf{114/115} (1989), 737--764.

\bibitem[ST24]{SektnanTipler2024}
Lars Martin Sektnan and Carl Tipler, \emph{Hermitian {Y}ang--{M}ills
connections on pullback bundles}, Calc. Var. Partial Differential
Equations \textbf{63} (2024), article~13.

\bibitem[Ses67]{Seshadri1967}
C.~S. Seshadri, \emph{Space of unitary vector bundles on a compact
{R}iemann surface}, Ann. of Math. (2) \textbf{85} (1967), 303--336.

\bibitem[Sib15]{Sibley2015}
Benjamin Sibley, \emph{Asymptotics of the {Y}ang--{M}ills flow for
holomorphic vector bundles over {K}\"ahler manifolds: the canonical
structure of the limit}, J. Reine Angew. Math. \textbf{706} (2015),
123--191.

\bibitem[SW15]{SibleyWentworth2015}
Benjamin Sibley and Richard~A. Wentworth, \emph{Analytic cycles,
{B}ott--{C}hern forms, and singular sets for the {Y}ang--{M}ills flow on
{K}\"ahler manifolds}, Adv. Math. \textbf{279} (2015), 501--531.

\bibitem[Sim88]{Simpson1988}
Carlos~T. Simpson, \emph{Constructing variations of {H}odge structure
using {Y}ang--{M}ills theory and applications to uniformization},
J. Amer. Math. Soc. \textbf{1} (1988), no.~4, 867--918.

\bibitem[SK67]{SinkhornKnopp1967}
Richard Sinkhorn and Paul Knopp, \emph{Concerning nonnegative matrices
and doubly stochastic matrices}, Pacific J. Math. \textbf{21} (1967),
no.~2, 343--348.

\bibitem[SHNGS18]{Soudry2018}
Daniel Soudry, Elad Hoffer, Mor~Shpigel Nacson, Suriya Gunasekar, and
Nathan Srebro, \emph{The implicit bias of gradient descent on separable
data}, J. Mach. Learn. Res. \textbf{19} (2018), 1--57.

\bibitem[Tak72]{Takemoto1972}
Fumio Takemoto, \emph{Stable vector bundles on algebraic surfaces},
Nagoya Math. J. \textbf{47} (1972), 29--48.

\bibitem[UY86]{UhlenbeckYau1986}
Karen Uhlenbeck and Shing-Tung Yau, \emph{On the existence of
{H}ermitian--{Y}ang--{M}ills connections in stable vector bundles},
Comm. Pure Appl. Math. \textbf{39} (1986), S257--S293.

\end{thebibliography}

\bigskip
\noindent{\sc Bernd Johannes Wuebben, New York, NY,}
\texttt{wuebben@gmail.com}

\end{document}
